\chapter{\texorpdfstring{Elementary Algorithms}{1. Elementary Algorithms}}

\section{\texorpdfstring{Array Transformations}{1.1 Array Transformations}}
\setcounter{section}{1}
\setcounter{subsection}{0}
\subsection{Sorting Algorithms}
The following functions are equivalent to \inlinecode{std::sort()}, taking random-access iterators as a half-open range $[{\inlinecodemath{lo}}, {\inlinecodemath{hi}})$ to be sorted. The range is sorted into ascending order after the function call. Optionally, a comparison function object specifying a strict weak ordering may be specified to replace the default \inlinecode{operator\textless{}}.

\begin{compactitem}
  \item \inlinecode{quicksort(lo, hi, comp = std::less\textless{}\textgreater{})} sorts the range using quicksort.
  \item \inlinecode{mergesort(lo, hi, comp = std::less\textless{}\textgreater{})} sorts the range using merge sort, which is stable.
  \item \inlinecode{heapsort(lo, hi, comp = std::less\textless{}\textgreater{})} sorts the range using heapsort.
  \item \inlinecode{combsort(lo, hi, comp = std::less\textless{}\textgreater{})} sorts the range using comb sort.
  \item \inlinecode{insertion\_sort(lo, hi, comp = std::less\textless{}\textgreater{})} sorts the range using insertion sort, which is stable.
  \item \inlinecode{shellsort(lo, hi, comp = std::less\textless{}\textgreater{})} sorts the range using shell sort.
  \item \inlinecode{counting\_sort(lo, hi, keys, key)} sorts the range by the index that \inlinecode{key(element)} returns in $[0, {\inlinecodemath{keys}})$, which is stable; unlike the shared interface above, it takes no comparator.
  \item \inlinecode{radix\_sort(lo, hi)} sorts an integer range using least-significant-byte radix sort; unlike the shared interface above, it takes no comparator.
  \item \inlinecode{bitonic\_sort(lo, hi, comp = std::less\textless{}\textgreater{})} sorts a range of any length using a compare-exchange sequence generated from that length.
\end{compactitem}

These functions are not meant to compete with standard library implementations in terms of speed. Instead, they demonstrate how common sorting algorithms can be concisely implemented in C++.

\begin{lstlisting}[language=C++]
#include <algorithm>
#include <climits>
#include <functional>
#include <iterator>
#include <type_traits>
#include <vector>
\end{lstlisting}
Quicksort repeatedly selects a pivot and partitions the range so that elements comparing less than the pivot precede it, elements comparing equal stay in the middle, and elements comparing greater follow it. Divide and conquer is then applied to the two outer ranges until the original range is sorted. The swaps used during partitioning make quicksort unstable. Despite having a worst case of $O(n^{2})$, quicksort is often faster in practice than merge sort and heapsort, which both have a worst case time complexity of $O(n \log n)$.

The pivot chosen in this implementation is always a middle element of the range to be sorted. Shuffling the input first or choosing each pivot randomly makes consistently unbalanced partitions unlikely. Median-of-three is another practical choice, while median-of-medians can guarantee a balanced pivot at greater constant-factor cost.

\textbf{Time Complexity}: $O(n)$ best (all equal), $O(n \log n)$ average, and $O(n^{2})$ worst.

\textbf{Space Complexity}: $O(\log n)$ auxiliary stack space.

\begin{lstlisting}[language=C++]
template<typename It, typename Compare = std::less<>>
void quicksort(It lo, It hi, Compare comp = Compare{}) {
  using T = typename std::iterator_traits<It>::value_type;
  while (hi - lo >= 2) {
    T pivot = *(lo + (hi - lo) / 2);
    It lt = lo, mid = lo, gt = hi;
    while (mid != gt) {
      if (comp(*mid, pivot)) {
        std::iter_swap(lt++, mid++);
      } else if (comp(pivot, *mid)) {
        std::iter_swap(mid, --gt);
      } else {
        ++mid;
      }
    }
    if (lt - lo < hi - gt) {
      quicksort(lo, lt, comp);
      lo = gt;
    } else {
      quicksort(gt, hi, comp);
      hi = lt;
    }
  }
}
\end{lstlisting}
Merge sort first divides a list into $n$ sublists of one element each, then recursively merges the sublists into sorted order until only a single sorted sublist remains. Merge sort is a stable sort, meaning that it preserves the relative order of elements which compare equal by \inlinecode{operator\textless{}} or the custom comparator given.

An analogous function in the C++ standard library is \inlinecode{std::stable\_sort()}, except that the implementation here requires sufficient memory to be available. When $O(n)$ auxiliary memory is not available, \inlinecode{std::stable\_sort()} falls back to a time complexity of $O(n \log^{2} n)$ whereas the implementation here will simply fail.

\textbf{Time Complexity}: $O(n \log n)$ in all cases.

\textbf{Space Complexity}: $O(\log n)$ auxiliary stack space and $O(n)$ auxiliary heap space.

\begin{lstlisting}[language=C++]
template<typename It, typename Compare = std::less<>>
void mergesort(It lo, It hi, Compare comp = Compare{}) {
  if (hi - lo < 2) {
    return;
  }
  It mid = lo + (hi - lo - 1) / 2, a = lo, c = mid + 1;
  mergesort(lo, mid + 1, comp);
  mergesort(mid + 1, hi, comp);
  using T = typename std::iterator_traits<It>::value_type;
  std::vector<T> merged;
  merged.reserve(hi - lo);
  while (a <= mid && c < hi) {
    merged.push_back(comp(*c, *a) ? *c++ : *a++);
  }
  merged.insert(merged.end(), a, mid + 1);
  merged.insert(merged.end(), c, hi);
  std::copy(merged.begin(), merged.end(), lo);
}
\end{lstlisting}
Heapsort first rearranges an array to satisfy the max-heap property. Then, it repeatedly pops the max element of the heap (the left, unsorted subrange), moving it to the beginning of the right, sorted subrange until the entire range is sorted. Heapsort has a better worst case time complexity than quicksort and also a better space complexity than merge sort, but its swaps make it unstable.

The C++ standard library equivalent is calling \inlinecode{std::make\_heap(lo, hi)}, followed by \inlinecode{std::sort\_heap(lo, hi)}.

\inlinecode{sift\_down()} restores the heap property below one node. The first loop applies it to each non-leaf from right to left, which builds the heap in $O(n)$ time. The second loop repeatedly moves the maximum to the back of the unsorted range and restores the heap among the remaining elements.

\textbf{Time Complexity}: $O(n \log n)$ in all cases.

\textbf{Space Complexity}: $O(1)$ auxiliary.

\begin{lstlisting}[language=C++]
template<typename It, typename Compare = std::less<>>
void heapsort(It lo, It hi, Compare comp = Compare{}) {
  if (hi - lo < 2) {
    return;
  }
  using T = typename std::iterator_traits<It>::value_type;
  auto sift_down = [&](It root, It end) {
    T value = *root;
    auto parent = root - lo, n = end - lo;
    while (2 * parent + 1 < n) {
      auto child = 2 * parent + 1;
      if (child + 1 < n && comp(*(lo + child), *(lo + child + 1))) {
        child++;
      }
      if (!comp(value, *(lo + child))) {
        break;
      }
      *(lo + parent) = *(lo + child);
      parent = child;
    }
    *(lo + parent) = value;
  };
  for (It root = lo + (hi - lo) / 2; root != lo;) {
    sift_down(--root, hi);
  }
  for (It end = hi; --end != lo;) {
    std::iter_swap(lo, end);
    sift_down(lo, end);
  }
}
\end{lstlisting}
Comb sort is an improved bubble sort. While bubble sort compares only adjacent elements, comb sort compares elements separated by a fixed gap, decreasing that gap after every pass. Large gaps move small values near the end toward the front quickly, avoiding the slow movement of such values in bubble sort. These nonadjacent swaps make comb sort unstable. The shrink factor $1.3$ is a common empirical choice.

\textbf{Time Complexity}: $O(n \log n)$ best and $O(n^{2})$ worst.

\textbf{Space Complexity}: $O(1)$ auxiliary.

\begin{lstlisting}[language=C++]
template<typename It, typename Compare = std::less<>>
void combsort(It lo, It hi, Compare comp = Compare{}) {
  int gap = static_cast<int>(hi - lo);
  bool swapped = true;
  while (gap > 1 || swapped) {
    if (gap > 1) {
      gap = gap * 10 / 13;
    }
    swapped = false;
    for (It it = lo; it + gap < hi; ++it) {
      if (comp(*(it + gap), *it)) {
        std::iter_swap(it, it + gap);
        swapped = true;
      }
    }
  }
}
\end{lstlisting}
Insertion sort builds the sorted range one element at a time. It scans left to right, and for each element shifts the larger elements of the already-sorted prefix one position to the right to open a slot where the element belongs. Although its average and worst cases are quadratic, it is simple, stable, in-place, and online (it can sort a stream as elements arrive).

Its strength is being adaptive: on an already-sorted or nearly-sorted range, each element travels only a short distance, giving $O(n)$ best-case time and $O(n + d)$ in general for $d$ total inversions. This is why it outperforms the asymptotically faster sorts on small or nearly-sorted ranges, and why hybrid sorts such as introsort and Timsort fall back to it once a subrange is small enough.

The comparison \inlinecode{comp(key, *(j - 1))} is strict, so an element never moves past an earlier element it compares equal to, which keeps the sort stable.

\textbf{Time Complexity}: $O(n + d)$, i.e. $O(n)$ best and $O(n^{2})$ average/worst.

\textbf{Space Complexity}: $O(1)$ auxiliary.

\begin{lstlisting}[language=C++]
template<typename It, typename Compare = std::less<>>
void insertion_sort(It lo, It hi, Compare comp = Compare{}) {
  using T = typename std::iterator_traits<It>::value_type;
  for (It i = lo; i != hi; ++i) {
    T key = *i;
    It j = i;
    while (j != lo && comp(key, *(j - 1))) {
      *j = *(j - 1);
      --j;
    }
    *j = key;
  }
}
\end{lstlisting}
Shell sort is insertion sort performed on subsequences of elements a fixed gap apart, repeating with smaller gaps until the final pass uses a gap of one and is an ordinary insertion sort. The earlier passes are what make that final pass cheap: each one leaves the range closer to sorted, and insertion sort runs in time proportional to the number of remaining inversions. Comparing elements a gap apart also moves a value many positions in one step, which is what insertion sort alone cannot do. Like comb sort above, the nonadjacent swaps make it unstable.

The gap sequence determines the bound, and choosing it well is the whole difficulty. This implementation uses Ciura's empirical sequence, extended upward by a factor of $2.25$, which is the fastest known for small and medium ranges though it has no proven bound. The theoretical choices are Sedgewick's sequence at $O(n^{4/3})$ and Pratt's at $O(n \log^{2} n)$, the latter being slower in practice despite the better bound.

\textbf{Time Complexity}: $O(n \log n)$ best, and $O(n^{4/3})$ or better in practice for the gaps used here.

\textbf{Space Complexity}: $O(1)$ auxiliary.

\begin{lstlisting}[language=C++]
template<typename It, typename Compare = std::less<>>
void shellsort(It lo, It hi, Compare comp = Compare{}) {
  static const int GAPS[] = {1, 4, 10, 23, 57, 132, 301, 701, 1577, 3548, 7983, 17962};
  int n = static_cast<int>(hi - lo), count = sizeof(GAPS) / sizeof(GAPS[0]);
  for (int g = count - 1; g >= 0; g--) {
    int gap = GAPS[g];
    if (gap >= n) {
      continue;
    }
    for (It it = lo + gap; it < hi; ++it) {
      typename std::iterator_traits<It>::value_type key = *it;
      It j = it;
      while (j - lo >= gap && comp(key, *(j - gap))) {
        *j = *(j - gap);
        j -= gap;
      }
      *j = key;
    }
  }
}
\end{lstlisting}
Counting sort sorts values from a small integer range by counting how many times each value occurs, converting those counts into starting offsets with a prefix sum, and then placing each element at the offset for its key. It never compares two elements, so the $O(n \log n)$ lower bound for comparison sorts does not apply, and it runs in linear time whenever the range of keys is proportional to the number of elements.

Placing elements in a second pass over the input from right to left keeps the sort stable, which is what makes it usable as the inner pass of the radix sort below. The \inlinecode{key} function maps an element to an index in $[0, {\inlinecodemath{keys}})$, so records can be sorted by one field while carrying the rest.

\textbf{Time Complexity}: $O(n + k)$ for $n$ elements with $k$ possible keys.

\textbf{Space Complexity}: $O(n + k)$ auxiliary.

\begin{lstlisting}[language=C++]
template<typename It, typename KeyFn>
void counting_sort(It lo, It hi, int keys, KeyFn key) {
  if (hi - lo < 2) {
    return;
  }
  std::vector<int> count(keys);
  for (It it = lo; it != hi; ++it) {
    count[key(*it)]++;
  }
  for (int i = 1; i < keys; ++i) {
    count[i] += count[i - 1];
  }
  using T = typename std::iterator_traits<It>::value_type;
  std::vector<T> res(hi - lo);
  for (It it = hi; it != lo;) {
    --it;
    res[--count[key(*it)]] = *it;
  }
  std::copy(res.begin(), res.end(), lo);
}
\end{lstlisting}
Radix sort is used to sort integer elements with a constant number of bits in linear time. This implementation works on ranges pointing to any signed or unsigned integer primitive. Signed values are handled by sorting on a key that flips the sign bit, which maps the two's-complement order onto the unsigned order so the most negative value sorts first.

Digits are processed from least significant to most significant, and each counting-sort pass is stable. After one pass, the values are sorted by the processed digit; stability ensures that a later pass preserves the ordering already established by all less-significant digits. Inductively, after the final pass the values are sorted by the entire key.

This implementation uses one byte per digit, so digits can be extracted with shifts and masks and the counting table has $2^8 = 256$ entries.

\textbf{Time Complexity}: $O(n \cdot w)$ for $n$ integers of $w$ bits each.

\textbf{Space Complexity}: $O(n + 2^{b})$ auxiliary for a radix of $b$ bits, i.e. $O(n)$ for constant $b$.

\begin{lstlisting}[language=C++]
template<typename It>
void radix_sort(It lo, It hi) {
  if (hi - lo < 2) {
    return;
  }
  const int radix_bits = 8;
  const int radix_base = 1 << radix_bits;  // e.g. 2^8 = 256
  const int radix_mask = radix_base - 1;   // e.g. 2^8 - 1 = 0xFF
  using T = typename std::iterator_traits<It>::value_type;
  static_assert(std::is_integral_v<T> && !std::is_same_v<T, bool>);
  using U = typename std::make_unsigned_t<T>;
  const int num_bits = sizeof(T) * CHAR_BIT;
  // Sort on an unsigned key. For signed types, flipping the sign bit sends the most negative value
  // to 0, mapping the signed order onto the unsigned order; logical shifts then extract each digit.
  auto key = [](T x) -> U {
    U u = static_cast<U>(x);
    return std::is_signed<T>::value ? (u ^ (U{1} << (num_bits - 1))) : u;
  };
  std::vector<T> buf(hi - lo);
  for (int pos = 0; pos < num_bits; pos += radix_bits) {
    std::vector<int> count(radix_base);
    for (It it = lo; it != hi; ++it) {
      count[(key(*it) >> pos) & radix_mask]++;
    }
    std::vector<T *> bucket(radix_base);
    T *curr = buf.data();
    for (int i = 0; i < radix_base; curr += count[i++]) {
      bucket[i] = curr;
    }
    for (It it = lo; it != hi; ++it) {
      *bucket[(key(*it) >> pos) & radix_mask]++ = *it;
    }
    std::copy(buf.begin(), buf.end(), lo);
  }
}
\end{lstlisting}
Bitonic sort is a sorting network: a fixed sequence of compare-exchanges on fixed positions, each swapping its pair when out of order. The sequence depends only on how many elements there are, never on what they are, so no comparison branches on data. It sorts the first half ascending and the second half descending, leaving a bitonic sequence that rises and then falls, which a fixed pattern of halving compare-exchanges merges. Splitting each merge at the largest power of two below its length is what admits lengths that are not powers of two, which the textbook formulation requires.

Its $O(n \log^{2} n)$ comparators exceed the $O(n \log n)$ comparisons of a comparison sort, so run one at a time it loses to every sort above, as the benchmark below shows. What it buys instead is that each group of comparators is independent and may run at once, which is why networks are the standard choice across SIMD lanes and on a GPU. The same independence pays off at a small size fixed at compile time: unrolled into branchless minimum and maximum, a network sorts eight integers about an order of magnitude faster than a general sort. Driving that same comparator sequence from a lookup table forfeits it, since each comparator then costs an indexed load and an unpredictable branch.

\textbf{Time Complexity}: $O(n \log^{2} n)$ in all cases.

\textbf{Space Complexity}: $O(\log n)$ auxiliary stack space.

\begin{lstlisting}[language=C++]
template<typename It, typename Compare = std::less<>>
void bitonic_sort(It lo, It hi, Compare comp = Compare{}) {
  auto merge = [&](auto &&merge, It lo, int n, bool up) -> void {
    if (n < 2) {
      return;
    }
    int m = 1;
    while (2 * m < n) {
      m *= 2;
    }
    for (int i = 0; i < n - m; i++) {
      if (comp(*(lo + i + m), *(lo + i)) == up) {
        std::iter_swap(lo + i, lo + i + m);
      }
    }
    merge(merge, lo, m, up);
    merge(merge, lo + m, n - m, up);
  };
  auto rec = [&](auto &&rec, It lo, int n, bool up) -> void {
    if (n < 2) {
      return;
    }
    rec(rec, lo, n / 2, !up);
    rec(rec, lo + n / 2, n - n / 2, up);
    merge(merge, lo, n, up);
  };
  rec(rec, lo, static_cast<int>(hi - lo), true);
}
\end{lstlisting}
\Needspace*{4\baselineskip}
\textbf{Example Usage}\par
\begin{lstlisting}[language=C++]
#include <cassert>
#include <ctime>
#include <iomanip>
#include <iostream>
#include <random>
#include <utility>
#include <vector>
using namespace std;

template<typename It>
void print_range(It lo, It hi) {
  while (lo != hi) {
    cout << *lo++ << " ";
  }
  cout << endl;
}

int main() {
  {  // Can be used to sort arrays like std::sort().
    int a[]{32, 71, 12, 45, 26, 80, 53, 33};
    quicksort(begin(a), end(a));
    assert(is_sorted(begin(a), end(a)));
  }
  {  // STL containers work too.
    vector<int> a{32, 71, 12, 45, 26, 80, 53, 33};
    quicksort(a.begin(), a.end());
    assert(is_sorted(a.begin(), a.end()));
  }
  {  // Reverse iterators work as expected.
    vector<int> a{32, 71, 12, 45, 26, 80, 53, 33};
    heapsort(a.rbegin(), a.rend());
    assert(is_sorted(a.rbegin(), a.rend()));
  }
  {  // Insertion sort is adaptive: an already-sorted range is left untouched in O(n).
    vector<int> a{12, 26, 32, 33, 45, 53, 71, 80};
    insertion_sort(a.begin(), a.end());
    assert(is_sorted(a.begin(), a.end()));
  }
  {  // We can sort doubles just as well.
    vector<double> a{1.1, -5.0, 6.23, 4.123, 155.2};
    combsort(a.begin(), a.end());
    assert(is_sorted(a.begin(), a.end()));
  }
  {  // Shell sort moves values far in one step, so it is not adaptive but is nearly in-place.
    vector<int> a{32, 71, 12, 45, 26, 80, 53, 33};
    shellsort(a.begin(), a.end());
    assert(is_sorted(a.begin(), a.end()));
  }
  {  // Counting sort keys records by one field, and equal keys keep their input order.
    vector<pair<int, char>> a{{2, 'a'}, {0, 'b'}, {2, 'c'}, {1, 'd'}, {0, 'e'}};
    counting_sort(a.begin(), a.end(), 3, [](const pair<int, char> &x) { return x.first; });
    vector<pair<int, char>> expected{{0, 'b'}, {0, 'e'}, {1, 'd'}, {2, 'a'}, {2, 'c'}};
    assert(a == expected);
  }
  {  // radix_sort() handles signed integers (including negatives), unlike a plain counting sort.
    vector<int> a{32, -71, 12, -45, 26, -80, 53, 33};
    radix_sort(a.begin(), a.end());
    assert(is_sorted(a.begin(), a.end()));
  }
  {  // Bitonic sort is oblivious too, and takes any length rather than only a power of two.
    vector<int> a{32, 71, 12, 45, 26, 80, 53, 33, 19};
    bitonic_sort(a.begin(), a.end());
    assert(is_sorted(a.begin(), a.end()));
  }
  {  // Empty and singleton ranges are valid no-ops.
    vector<int> empty, single{42};
    quicksort(empty.begin(), empty.end());
    mergesort(empty.begin(), empty.end());
    heapsort(empty.begin(), empty.end());
    insertion_sort(empty.begin(), empty.end());
    combsort(empty.begin(), empty.end());
    radix_sort(empty.begin(), empty.end());
    mergesort(single.begin(), single.end());
    assert(empty.empty() && single[0] == 42);
  }

  // Stable sort.
  const vector<double> a{3.14, 1.41, 2.72, 4.67, 1.73, 1.32, 1.62, 2.58};
  {
    vector<double> v(a);
    cout << "mergesort() with default comparisons: ";
    mergesort(v.begin(), v.end());
    print_range(v.begin(), v.end());
  }
  {
    vector<double> v(a);
    cout << "mergesort() with integer comparisons: ";
    mergesort(v.begin(), v.end(), [](double i, double j) {
      return static_cast<int>(i) < static_cast<int>(j);
    });
    print_range(v.begin(), v.end());
  }
  cout << "------" << endl;

  mt19937 rng(1234567);  // Fixed seed for reproducibility.
  vector<int> data(5000000);
  const int maxval = 10000000;
  for (int &x : data) {
    x = static_cast<int>(rng() % maxval);  // limit magnitude for counting sort
  }
  cout << "Sorting five million integers..." << endl;
  cout.precision(3);
  auto benchmark = [&](const string &name, auto sort) {
    vector<int> v = data;
    clock_t start = clock();
    sort(v.begin(), v.end());
    double t = static_cast<double>(clock() - start) / CLOCKS_PER_SEC;
    cout << setw(18) << left << name + "(): " << fixed << t << "s" << endl;
    assert(is_sorted(v.begin(), v.end()));
  };
  benchmark("std::sort", [](auto lo, auto hi) { sort(lo, hi); });
  benchmark("quicksort", [](auto lo, auto hi) { quicksort(lo, hi); });
  benchmark("mergesort", [](auto lo, auto hi) { mergesort(lo, hi); });
  benchmark("heapsort", [](auto lo, auto hi) { heapsort(lo, hi); });
  benchmark("combsort", [](auto lo, auto hi) { combsort(lo, hi); });
  benchmark("shellsort", [](auto lo, auto hi) { shellsort(lo, hi); });
  benchmark("counting_sort", [](auto lo, auto hi) {
    counting_sort(lo, hi, maxval, [](int x) { return x; });
  });
  benchmark("radix_sort", [](auto lo, auto hi) { radix_sort(lo, hi); });
  benchmark("bitonic_sort", [](auto lo, auto hi) { bitonic_sort(lo, hi); });
  return 0;
}
\end{lstlisting}
\begin{exampleoutput}
mergesort() with default comparisons: 1.32 1.41 1.62 1.73 2.58 2.72 3.14 4.67
mergesort() with integer comparisons: 1.41 1.73 1.32 1.62 2.72 2.58 3.14 4.67
------
Sorting five million integers...
std::sort():      0.285s
quicksort():      0.325s
mergesort():      0.517s
heapsort():       0.509s
combsort():       0.469s
shellsort():      0.574s
counting_sort():  0.044s
radix_sort():     0.024s
bitonic_sort():   1.040s
\end{exampleoutput}
\subsection{Partitioning and Selection}
Partition a range around a pivot, or select the element with a given sorted rank without fully sorting the range. A comparator \inlinecode{comp} defines the ordering (defaults to \inlinecode{std::less\textless{}\textgreater{}}). Three-way partitioning rearranges the elements into consecutive groups that precede, are equivalent to, and follow the pivot. \inlinecode{std::partition()} can be used to place elements satisfying a unary predicate before those that do not, but forms only two groups and does not preserve relative order. Forming three groups can be done by first partitioning by \inlinecode{comp(x, pivot)}, then partitioning the remaining suffix by \inlinecode{!comp(pivot, x)}; using \inlinecode{std::stable\_partition()} for both calls preserves relative order within each group.

\inlinecode{partition\_three\_way()} implements the Dutch National Flag algorithm, forming the three groups in one pass with $O(1)$ auxiliary space. At every iteration, $[{\inlinecodemath{lo}}, {\inlinecodemath{lt}})$ contains elements preceding the pivot, $[{\inlinecodemath{lt}}, {\inlinecodemath{cur}})$ contains equivalent elements, $[{\inlinecodemath{cur}}, {\inlinecodemath{gt}})$ remains unclassified, and $[{\inlinecodemath{gt}}, {\inlinecodemath{hi}})$ contains elements following the pivot. The current element \inlinecode{*cur} is compared with the pivot. If it precedes the pivot, it is swapped with \inlinecode{*lt} and both \inlinecode{lt} and \inlinecode{cur} advance. If it is equivalent, only \inlinecode{cur} advances. Otherwise, it is swapped with \inlinecode{*--gt} without advancing \inlinecode{cur}, since the incoming element remains unclassified.

Quickselect repeatedly applies this partition and continues only into the group containing the desired rank. This is the same task as \inlinecode{std::nth\_element()}: after the call, \inlinecode{*nth} is the value that would appear there in comparator order, no earlier value follows it, and no later value precedes it. Choosing pivots uniformly at random gives expected linear time, while the three-way split avoids unnecessary work on duplicate-heavy inputs.

\begin{compactitem}
  \item \inlinecode{partition\_three\_way(lo, hi, pivot, comp = std::less\textless{}\textgreater{}())} rearranges $[{\inlinecodemath{lo}}, {\inlinecodemath{hi}})$ in-place and returns a pair of iterators (\inlinecode{mid1}, \inlinecode{mid2}). The three resulting ranges contain elements that precede, are equivalent to, and follow \inlinecode{pivot} according to \inlinecode{comp}, respectively. The resulting partition is not stable.
  \item \inlinecode{quickselect(lo, nth, hi, comp = std::less\textless{}\textgreater{}())} rearranges $[{\inlinecodemath{lo}}, {\inlinecodemath{hi}})$ in-place around the 0-based rank represented by iterator \inlinecode{nth} according to \inlinecode{comp}. This requires random-access iterators.
\end{compactitem}

Random pivots leave an adversary free to force $O(n^{2})$, which the median-of-medians rule removes. Split the range into groups of five, sort each, and pivot on the median of the group medians: at least three elements in half of the groups precede it, so every partition discards at least three tenths of the range. Finding that median is itself a selection, solved by recursing on the $n/5$ medians gathered to the front, and $T(n) = T(n/5) + T(7n/10) + O(n)$ is linear because the two fractions sum to less than one, which five is the smallest group size to achieve. Each group is sorted directly, since at five elements \inlinecode{std::sort} is an inlined insertion sort. This guarantee is not free, but on random input the cost over a random pivot is modest.

\begin{compactitem}
  \item \inlinecode{quickselect\_linear(lo, nth, hi, comp = std::less\textless{}\textgreater{}())} does the same in linear time even in the worst case, choosing each pivot by the median-of-medians rule instead of at random.
\end{compactitem}

\Needspace*{3\baselineskip}
\textbf{Time Complexity}:
\begin{compactitem}
  \item $O(n)$ per call to \inlinecode{partition\_three\_way()}, where $n$ is the range length.
  \item $O(n)$ expected per call to \inlinecode{quickselect()}, and $O(n^{2})$ in the worst case.
  \item $O(n)$ per call to \inlinecode{quickselect\_linear()} in all cases.
\end{compactitem}

\Needspace*{3\baselineskip}
\textbf{Space Complexity}:
\begin{compactitem}
  \item $O(1)$ auxiliary for \inlinecode{partition\_three\_way()} and \inlinecode{quickselect()}.
  \item $O(\log n)$ auxiliary stack space for \inlinecode{quickselect\_linear()}.
\end{compactitem}

\begin{lstlisting}[language=C++]
#include <algorithm>
#include <chrono>
#include <functional>
#include <iterator>
#include <random>
#include <utility>

template<typename It, typename T, typename Compare = std::less<>>
std::pair<It, It> partition_three_way(It lo, It hi, const T &pivot, Compare comp = Compare{}) {
  It lt = lo, cur = lo, gt = hi;
  while (cur != gt) {
    if (comp(*cur, pivot)) {
      std::iter_swap(lt++, cur++);
    } else if (comp(pivot, *cur)) {
      std::iter_swap(cur, --gt);
    } else {
      ++cur;
    }
  }
  return {lt, gt};
}

template<typename It, typename Compare = std::less<>>
void quickselect(It lo, It nth, It hi, Compare comp = Compare{}) {
  using T = typename std::iterator_traits<It>::value_type;
  static std::mt19937 rng(std::chrono::steady_clock::now().time_since_epoch().count());
  while (hi - lo > 1) {
    std::uniform_int_distribution<int> dist(0, hi - lo - 1);
    T pivot = *(lo + dist(rng));
    auto [lt, gt] = partition_three_way(lo, hi, pivot, comp);
    if (nth < lt) {
      hi = lt;
    } else if (nth >= gt) {
      lo = gt;
    } else {
      return;
    }
  }
}

template<typename It, typename Compare = std::less<>>
void quickselect_linear(It lo, It nth, It hi, Compare comp = Compare{}) {
  while (hi - lo > 5) {
    int groups = 0;
    for (It g = lo; g < hi; g += 5, groups++) {
      It ge = hi - g < 5 ? hi : g + 5;
      std::sort(g, ge, comp);  // At most 5 elements, so this is an O(1) step.
      std::iter_swap(lo + groups, g + (ge - g) / 2);
    }
    It mid = lo + groups / 2;
    quickselect_linear(lo, mid, lo + groups, comp);
    using T = typename std::iterator_traits<It>::value_type;
    T pivot = *mid;
    std::pair<It, It> parts = partition_three_way(lo, hi, pivot, comp);
    if (nth < parts.first) {
      hi = parts.first;
    } else if (nth < parts.second) {
      return;
    } else {
      lo = parts.second;
    }
  }
  std::sort(lo, hi, comp);
}
\end{lstlisting}
\Needspace*{4\baselineskip}
\textbf{Example Usage}\par
\begin{lstlisting}[language=C++]
#include <cassert>
#include <vector>
using namespace std;

int main() {
  vector<int> values{2, 0, 1, 2, 1, 0, 1};
  // On a range containing only 0, 1, and 2, partitioning around 1 sorts the range.
  partition_three_way(values.begin(), values.end(), 1);
  assert(is_sorted(values.begin(), values.end()));

  vector<int> b{4, 2, 5, 3, 3, 1};
  auto [mid1, mid2] = partition_three_way(b.begin(), b.end(), 3);
  for (auto it = b.begin(); it != mid1; ++it) {
    assert(*it < 3);
  }
  for (auto it = mid1; it != mid2; ++it) {
    assert(*it == 3);
  }
  for (auto it = mid2; it != b.end(); ++it) {
    assert(*it > 3);
  }

  vector<int> a{5, 6, 4, 3, 2, 6, 7, 9, 3};
  int n = static_cast<int>(a.size());
  quickselect(a.begin(), a.begin() + n / 2, a.end());
  assert(a[n / 2] == 5);
  // Values left of the median are <=, and values right are >= (the exact order within each side is
  // randomized, since the pivot is chosen at random).
  for (int i = 0; i < n / 2; i++) {
    assert(a[i] <= a[n / 2]);
  }
  for (int i = n / 2 + 1; i < n; i++) {
    assert(a[i] >= a[n / 2]);
  }

  // The median-of-medians variant gives the same answer with a worst case that is still linear.
  vector<int> c{5, 6, 4, 3, 2, 6, 7, 9, 3};
  quickselect_linear(c.begin(), c.begin() + n / 2, c.end());
  assert(c[n / 2] == 5);

  vector<int> desc{1, 4, 2, 3};
  quickselect(desc.begin(), desc.begin() + 1, desc.end(), greater<int>());
  assert(desc[1] == 3);
  return 0;
}
\end{lstlisting}
\subsection{Array Rotation}
Rotates the half-open iterator range $[{\inlinecodemath{lo}}, {\inlinecodemath{hi}})$ left around the split point \inlinecode{mid}, where ${\inlinecodemath{lo}} \leq {\inlinecodemath{mid}} \leq {\inlinecodemath{hi}}$. Writing the adjacent subranges $[{\inlinecodemath{lo}}, {\inlinecodemath{mid}})$ and $[{\inlinecodemath{mid}}, {\inlinecodemath{hi}})$ as \inlinecode{A} and \inlinecode{B}, respectively, the operation transforms \inlinecode{A B} into \inlinecode{B A} while preserving the relative order within each subrange. The standard library function \inlinecode{std::rotate(lo, mid, hi)} performs the same rearrangement.

All three versions below operate in place, using different algorithms and iterator capabilities.

\begin{compactitem}
  \item \inlinecode{rotate1(lo, mid, hi)} requires ForwardIterators and repeatedly swaps the next elements of the two unfinished subranges. When the right iterator reaches \inlinecode{hi}, it wraps to the current \inlinecode{mid}; when the left iterator reaches \inlinecode{mid}, that boundary advances to the right iterator.
  \item \inlinecode{rotate2(lo, mid, hi)} requires BidirectionalIterators, applying a trick with three reversals. Writing the input subranges as \inlinecode{A B}, two initial reversals yield \inlinecode{reverse(A) reverse(B)}; then the whole range is reversed to yield \inlinecode{B A}, restoring the order within each subrange.
  \item \inlinecode{rotate3(lo, mid, hi)} requires random-access iterators, applying a juggling algorithm which first divides the range into \inlinecode{gcd(hi - lo, mid - lo)} sets and then rotates the corresponding elements in each set. The method follows the permutation that sends each position to its rotated destination. These form that many disjoint cycles, so walking each cycle once moves every element to its final position.
\end{compactitem}

\textbf{Time Complexity}: $O(n)$ per call to all versions, where $n$ is the distance between \inlinecode{lo} and \inlinecode{hi}.

\textbf{Space Complexity}: $O(1)$ auxiliary for all versions.

\begin{lstlisting}[language=C++]
#include <algorithm>
#include <numeric>

template<typename It>
void rotate1(It lo, It mid, It hi) {
  if (lo == mid || mid == hi) {
    return;
  }
  It next = mid;
  while (lo != next) {
    std::iter_swap(lo++, next++);
    if (next == hi) {
      next = mid;
    } else if (lo == mid) {
      mid = next;
    }
  }
}

template<typename It>
void rotate2(It lo, It mid, It hi) {
  if (lo == mid || mid == hi) {
    return;
  }
  std::reverse(lo, mid);
  std::reverse(mid, hi);
  std::reverse(lo, hi);
}

template<typename It>
void rotate3(It lo, It mid, It hi) {
  if (lo == mid || mid == hi) {
    return;
  }
  int n = static_cast<int>(hi - lo);
  int jump = static_cast<int>(mid - lo);
  int g = std::gcd(jump, n), cycle = n / g;
  for (int i = 0; i < g; i++) {
    int curr = i, next;
    for (int j = 0; j < cycle - 1; j++) {
      next = curr + jump;
      if (next >= n) {
        next -= n;
      }
      std::iter_swap(lo + curr, lo + next);
      curr = next;
    }
  }
}
\end{lstlisting}
\Needspace*{4\baselineskip}
\textbf{Example Usage}\par
\begin{lstlisting}[language=C++]
#include <algorithm>
#include <cassert>
#include <vector>
using namespace std;

int main() {
  vector<int> a0(10000), a1, a2, a3;
  iota(a0.begin(), a0.end(), 0);
  a1 = a2 = a3 = a0;
  int mid = 5678;
  std::rotate(a0.begin(), a0.begin() + mid, a0.end());
  rotate1(a1.begin(), a1.begin() + mid, a1.end());
  rotate2(a2.begin(), a2.begin() + mid, a2.end());
  rotate3(a3.begin(), a3.begin() + mid, a3.end());
  assert(a0 == a1 && a0 == a2 && a0 == a3);

  vector<int> no_op{1, 2, 3};
  rotate1(no_op.begin(), no_op.begin(), no_op.end());
  assert((no_op == vector<int>{1, 2, 3}));
  rotate2(no_op.begin(), no_op.end(), no_op.end());
  assert((no_op == vector<int>{1, 2, 3}));
  rotate3(no_op.begin(), no_op.begin(), no_op.end());
  assert((no_op == vector<int>{1, 2, 3}));

  vector<int> a{2, 4, 2, 0, 5, 10, 7, 3, 7, 1};

  // Insertion sort.
  for (auto i = a.begin(); i != a.end(); ++i) {
    rotate1(upper_bound(a.begin(), i, *i), i, i + 1);
  }
  assert((a == vector<int>{0, 1, 2, 2, 3, 4, 5, 7, 7, 10}));

  // Simple rotation to the left.
  rotate2(a.begin(), a.begin() + 1, a.end());
  assert((a == vector<int>{1, 2, 2, 3, 4, 5, 7, 7, 10, 0}));

  // Simple rotation to the right.
  rotate3(a.rbegin(), a.rbegin() + 1, a.rend());
  assert((a == vector<int>{0, 1, 2, 2, 3, 4, 5, 7, 7, 10}));
  return 0;
}
\end{lstlisting}
\subsection{Coordinate Compression}
Given a range of $n$ numerical elements, reassign each element to an integer in the domain $[0, k)$, where $k$ is the number of distinct elements in the original range, while preserving the initial relative ordering of elements. That is, if $a$ is the array of original values and $b$ is the array of compressed values, then every pair of indices $i, j$ in $[0, n)$ shall satisfy $a[i] < a[j]$ if and only if $b[i] < b[j]$.

The two functions below take a ForwardIterator range, rewrite it in place, and discard the mapping. The comparator \inlinecode{comp} defines the value ordering: \inlinecode{comp(a, b)} is true when \inlinecode{a} precedes \inlinecode{b}.

\begin{compactitem}
  \item \inlinecode{compress1(lo, hi, comp = std::less\textless{}\textgreater{}())} performs the compression by sorting the array, removing duplicates, and binary searching for the position of each original value.
  \item \inlinecode{compress2(lo, hi, comp = std::less\textless{}\textgreater{}())} achieves the same result by inserting all values into an \inlinecode{std::map}, which automatically removes duplicate values and supports efficient lookups of the compressed values.
\end{compactitem}

\textbf{Time Complexity}: $O(n \log n)$ per call to \inlinecode{compress1(lo, hi)} and \inlinecode{compress2(lo, hi)}, where $n$ is the distance between \inlinecode{lo} and \inlinecode{hi}.

\textbf{Space Complexity}: $O(n)$ auxiliary for \inlinecode{compress1()} and \inlinecode{compress2()}.

\begin{lstlisting}[language=C++]
#include <algorithm>
#include <cassert>
#include <functional>
#include <iterator>
#include <map>
#include <vector>

template<typename It, typename Compare = std::less<>>
void compress1(It lo, It hi, Compare comp = Compare{}) {
  using T = typename std::iterator_traits<It>::value_type;
  std::vector<T> v(lo, hi);
  std::sort(v.begin(), v.end(), comp);
  v.erase(
      std::unique(
          v.begin(), v.end(), [&](const T &a, const T &b) { return !comp(a, b) && !comp(b, a); }
      ),
      v.end()
  );
  for (It it = lo; it != hi; ++it) {
    *it = static_cast<int>(std::lower_bound(v.begin(), v.end(), *it, comp) - v.begin());
  }
}

template<typename It, typename Compare = std::less<>>
void compress2(It lo, It hi, Compare comp = Compare{}) {
  using T = typename std::iterator_traits<It>::value_type;
  std::map<T, int, Compare> m(comp);
  for (It it = lo; it != hi; ++it) {
    m[*it] = 0;
  }
  int id = 0;
  for (auto &[key, val] : m) {
    val = id++;
  }
  for (It it = lo; it != hi; ++it) {
    *it = m[*it];
  }
}
\end{lstlisting}
\inlinecode{Compressor} retains the sorted table of distinct values so that arbitrary values can be mapped to and from compressed ranks long after construction, such as for offline queries arriving separately from the array being compressed. It uses \inlinecode{std::less\textless{}T\textgreater{}} by default; to customize the ordering, instantiate \inlinecode{Compressor\textless{}T, Compare\textgreater{}} and pass the comparator to either constructor.

\begin{compactitem}
  \item \inlinecode{Compressor\textless{}T\textgreater{}()} constructs an empty compressor. Register values with \inlinecode{add(x)}, then call \inlinecode{build()} once before the first query.
  \item \inlinecode{Compressor\textless{}T\textgreater{}(lo, hi)} constructs a compressor from the half-open iterator range $[{\inlinecodemath{lo}}, {\inlinecodemath{hi}})$.
  \item \inlinecode{size()} returns the number of distinct registered values $k$.
  \item \inlinecode{value(r)} returns the original value with rank \inlinecode{r}, inverting \inlinecode{rank()}.
  \item \inlinecode{contains(x)} returns whether \inlinecode{x} is a registered value.
  \item \inlinecode{rank(x)} returns the compressed value (rank) of \inlinecode{x} in $[0, k)$. \inlinecode{x} must have been registered.
\end{compactitem}

\Needspace*{3\baselineskip}
\textbf{Time Complexity}:
\begin{compactitem}
  \item $O(n \log n)$ per call to \inlinecode{Compressor(lo, hi)}, where $n$ is the range length.
  \item $O(1)$ amortized per call to \inlinecode{add()}, and $O(m \log m)$ per call to \inlinecode{build()}, where $m$ is the total number of values registered.
  \item $O(\log k)$ per call to \inlinecode{rank()} and \inlinecode{contains()}, where $k$ is the number of distinct values.
  \item $O(1)$ per call to \inlinecode{size()} and \inlinecode{value()}.
\end{compactitem}

\textbf{Space Complexity}: $O(k)$ storage, where $k$ is the number of distinct registered values.

\begin{lstlisting}[language=C++]
template<typename T, typename Compare = std::less<T>>
class Compressor {
  Compare comp;
  std::vector<T> v;

 public:
  explicit Compressor(Compare comp = Compare{}) : comp(std::move(comp)) {}

  template<typename It>
  Compressor(It lo, It hi, Compare comp = Compare{}) : comp(std::move(comp)), v(lo, hi) {
    build();
  }

  void build() {
    std::sort(v.begin(), v.end(), comp);
    v.erase(
        std::unique(
            v.begin(), v.end(), [&](const T &a, const T &b) { return !comp(a, b) && !comp(b, a); }
        ),
        v.end()
    );
  }

  void add(const T &x) { v.push_back(x); }
  int size() const { return static_cast<int>(v.size()); }
  const T &value(int r) const { return v[r]; }
  bool contains(const T &x) const { return std::binary_search(v.begin(), v.end(), x, comp); }

  int rank(const T &x) const {
    int r = static_cast<int>(std::lower_bound(v.begin(), v.end(), x, comp) - v.begin());
    assert(r < size() && !comp(x, v[r]));  // x must be registered.
    return r;
  }
};
\end{lstlisting}
\Needspace*{4\baselineskip}
\textbf{Example Usage}\par
\begin{lstlisting}[language=C++]
#include <cassert>
using namespace std;

int main() {
  {
    vector<int> a{1, 30, 30, 7, 9, 8, 99, 99};
    compress1(a.begin(), a.end());
    assert((a == vector<int>{0, 4, 4, 1, 3, 2, 5, 5}));
  }
  {
    vector<int> a{1, 30, 30, 7, 9, 8, 99, 99};
    compress2(a.begin(), a.end());
    assert((a == vector<int>{0, 4, 4, 1, 3, 2, 5, 5}));
  }
  {  // Non-integral types work too, as long as ints can be assigned to them.
    vector<double> a{0.5, -1.0, 3, -1.0, 20, 0.5};
    compress1(a.begin(), a.end());
    assert((a == vector<double>{1, 0, 2, 0, 3, 1}));
  }
  {
    vector<int> a{10, 5, 30, 5, 20};
    Compressor<int> cc(a.begin(), a.end());
    assert(cc.size() == 4);
    assert(cc.rank(5) == 0);
    assert(cc.rank(30) == 3);
    assert(cc.value(2) == 20);
    assert(cc.contains(10) && !cc.contains(7));
  }
  {  // Register values incrementally, then build once before querying.
    Compressor<double> cc;
    cc.add(0.5);
    cc.add(-1.0);
    cc.add(0.5);
    cc.build();
    assert(cc.size() == 2);
    assert(cc.rank(0.5) == 1);
    assert(cc.value(0) == -1.0);
  }
  {
    vector<int> a{10, 30, 20, 10};
    Compressor<int, greater<int>> cc(a.begin(), a.end());
    assert(cc.rank(30) == 0 && cc.rank(10) == 2);
    assert(cc.value(1) == 20);
    compress1(a.begin(), a.end(), greater<int>());
    assert((a == vector<int>{2, 0, 1, 2}));
  }
  return 0;
}
\end{lstlisting}
\subsection{Counting Inversions}
The number of inversions for a sequence $a$ is the number of ordered pairs $(i, j)$ such that $i < j$ and $a[i] > a[j]$. This is roughly how ``close'' an array is to being sorted, but is \textit{not} the minimum number of swaps required to sort the array. If the array is sorted, then the inversion count is $0$. If the array is sorted in decreasing order, then the inversion count is maximal. The following two functions are each techniques to efficiently count inversions. In the merge sort approach, whenever the merge step emits an element from the right half, that element jumps ahead of every unmerged left-half element, and exactly that many inversions are added to the count.

\begin{compactitem}
  \item \inlinecode{inversions(lo, hi, comp = std::less\textless{}\textgreater{})} uses merge sort to return the number of inversions given two random-access iterators as a range $[{\inlinecodemath{lo}}, {\inlinecodemath{hi}})$. The input range will be sorted after the function call. Optionally, a comparison function object specifying a strict weak ordering may be specified to replace the default \inlinecode{operator\textless{}}.
  \item \inlinecode{inversions(a)} uses coordinate compression and a Fenwick tree to return the number of inversions in an integer vector without modifying it.
\end{compactitem}

\Needspace*{3\baselineskip}
\textbf{Time Complexity}:
\begin{compactitem}
  \item $O(n \log n)$ per call to \inlinecode{inversions(lo, hi)}, where $n$ is the distance between \inlinecode{lo} and \inlinecode{hi}.
  \item $O(n \log n)$ per call to \inlinecode{inversions(a)}.
\end{compactitem}

\Needspace*{3\baselineskip}
\textbf{Space Complexity}:
\begin{compactitem}
  \item $O(\log n)$ auxiliary stack space and $O(n)$ auxiliary heap space for \inlinecode{inversions(lo, hi)}.
  \item $O(n)$ auxiliary heap space for \inlinecode{inversions(a)}.
\end{compactitem}

\begin{lstlisting}[language=C++]
#include <algorithm>
#include <cstdint>
#include <functional>
#include <iterator>
#include <vector>

template<typename It, typename Compare = std::less<>>
int64_t inversions(It lo, It hi, Compare comp = Compare{}) {
  int n = static_cast<int>(hi - lo);
  if (n < 2) {
    return 0;
  }
  It mid = lo + (n - 1) / 2, a = lo, c = mid + 1;
  int64_t res = inversions(lo, mid + 1, comp) + inversions(mid + 1, hi, comp);
  using T = typename std::iterator_traits<It>::value_type;
  std::vector<T> merged;
  merged.reserve(n);
  while (a <= mid && c < hi) {
    if (comp(*c, *a)) {
      merged.push_back(*(c++));
      res += (mid - a) + 1;
    } else {
      merged.push_back(*(a++));
    }
  }
  if (a > mid) {
    for (It it = c; it != hi; ++it) {
      merged.push_back(*it);
    }
  } else {
    for (It it = a; it <= mid; ++it) {
      merged.push_back(*it);
    }
  }
  for (It it = lo; it != hi; ++it) {
    *it = merged[it - lo];
  }
  return res;
}

int64_t inversions(const std::vector<int> &a) {
  int n = static_cast<int>(a.size());
  std::vector<int> values(a);
  std::sort(values.begin(), values.end());
  values.resize(std::unique(values.begin(), values.end()) - values.begin());
  std::vector<int> bit(values.size() + 1);
  int64_t res = 0;
  for (int i = n - 1; i >= 0; i--) {
    int id =
        static_cast<int>(std::lower_bound(values.begin(), values.end(), a[i]) - values.begin()) + 1;
    for (int j = id - 1; j > 0; j -= j & -j) {
      res += bit[j];
    }
    for (int j = id; j < static_cast<int>(bit.size()); j += j & -j) {
      bit[j]++;
    }
  }
  return res;
}
\end{lstlisting}
\Needspace*{4\baselineskip}
\textbf{Example Usage}\par
\begin{lstlisting}[language=C++]
#include <cassert>
using namespace std;

int main() {
  {
    vector<int> a{6, 9, 1, 14, 8, 12, 3, 2};
    assert(inversions(a.begin(), a.end()) == 16);
    assert(is_sorted(a.begin(), a.end()));
  }
  {
    vector<int> a{6, 9, 1, 14, 8, 12, 3, 2};
    assert(inversions(a) == 16);
    assert((a == vector<int>{6, 9, 1, 14, 8, 12, 3, 2}));
  }
  {
    vector<int> a{2, 2, 1};
    assert(inversions(a) == 2);  // Equal elements are not inversions.
  }
  return 0;
}
\end{lstlisting}

\section{\texorpdfstring{1D Scan and Window Techniques}{1.2 1D Scan and Window Techniques}}
\setcounter{section}{2}
\setcounter{subsection}{0}
\subsection{Prefix Sums}
Precomputes prefix sums so range-sum queries can be answered in constant time. This is one of the most common array preprocessing tools, and also appears as a building block for subarray sums, difference arrays, and two-dimensional grids. Each table entry stores the sum of all elements before it, so any range sum is the difference of two entries; in two dimensions, rectangle sums combine four entries by inclusion-exclusion.

Prefix sums can also be combined with a frequency table to count subarrays having a target sum. When the current prefix sum is $s$, every earlier prefix sum equal to $s - t$ identifies a subarray with sum $t$. Recording frequencies rather than only whether a prefix exists counts duplicate sums correctly.

\begin{compactitem}
  \item \inlinecode{prefix\_sums(a)} returns the prefix sum array \inlinecode{pref} of length $n + 1$, with \inlinecode{pref[0] = 0} and \inlinecode{pref[i + 1]} equal to the sum of \inlinecode{a[0]} through \inlinecode{a[i]}.
  \item \inlinecode{range\_sum(pref, lo, hi)} returns the sum of range $[{\inlinecodemath{lo}}, {\inlinecodemath{hi}}]$.
  \item \inlinecode{prefix\_sums\_2d(a)} returns a two-dimensional prefix sum table for matrix \inlinecode{a}.
  \item \inlinecode{rect\_sum(pref, r1, c1, r2, c2)} returns the sum of the rectangle with rows $[{\inlinecodemath{r1}}, {\inlinecodemath{r2}}]$ and columns $[{\inlinecodemath{c1}}, {\inlinecodemath{c2}}]$.
  \item \inlinecode{count\_subarrays\_with\_sum(a, target)} returns the number of contiguous subarrays whose sum is \inlinecode{target}.
\end{compactitem}

Overflow warning: All prefix sums, target differences, inclusion-exclusion results, and returned counts must fit in \inlinecode{int64\_t}.

\Needspace*{3\baselineskip}
\textbf{Time Complexity}:
\begin{compactitem}
  \item $O(n)$ per call to \inlinecode{prefix\_sums(a)}, where $n$ is the array size.
  \item $O(R \cdot C)$ per call to \inlinecode{prefix\_sums\_2d(a)}, where $R$ and $C$ are the number of rows and columns of \inlinecode{a}, respectively.
  \item $O(1)$ per range or rectangle query.
  \item $O(n)$ expected per call to \inlinecode{count\_subarrays\_with\_sum(a, target)}.
\end{compactitem}

\Needspace*{3\baselineskip}
\textbf{Space Complexity}:
\begin{compactitem}
  \item $O(n)$ for the array returned by \inlinecode{prefix\_sums(a)}.
  \item $O(R \cdot C)$ for the table returned by \inlinecode{prefix\_sums\_2d(a)}.
  \item $O(1)$ auxiliary for \inlinecode{range\_sum()} and \inlinecode{rect\_sum()}.
  \item $O(n)$ auxiliary for \inlinecode{count\_subarrays\_with\_sum(a, target)}.
\end{compactitem}

\begin{lstlisting}[language=C++]
#include <cassert>
#include <cstdint>
#include <unordered_map>
#include <vector>

std::vector<int64_t> prefix_sums(const std::vector<int> &a) {
  std::vector<int64_t> pref(a.size() + 1);
  for (int i = 0; i < static_cast<int>(a.size()); i++) {
    pref[i + 1] = pref[i] + a[i];
  }
  return pref;
}

int64_t range_sum(const std::vector<int64_t> &pref, int lo, int hi) {
  assert(0 <= lo && lo <= hi && hi < static_cast<int>(pref.size()) - 1);
  return pref[hi + 1] - pref[lo];
}

std::vector<std::vector<int64_t>> prefix_sums_2d(const std::vector<std::vector<int>> &a) {
  int rows = static_cast<int>(a.size());
  int cols = a.empty() ? 0 : static_cast<int>(a[0].size());
  std::vector<std::vector<int64_t>> pref(rows + 1, std::vector<int64_t>(cols + 1));
  for (int i = 0; i < rows; i++) {
    for (int j = 0; j < cols; j++) {
      pref[i + 1][j + 1] = a[i][j] + pref[i][j + 1] + pref[i + 1][j] - pref[i][j];
    }
  }
  return pref;
}

int64_t rect_sum(const std::vector<std::vector<int64_t>> &pref, int r1, int c1, int r2, int c2) {
  assert(
      !pref.empty() && !pref[0].empty() && 0 <= r1 && r1 <= r2 &&
      r2 < static_cast<int>(pref.size()) - 1 && 0 <= c1 && c1 <= c2 &&
      c2 < static_cast<int>(pref[0].size()) - 1
  );
  return pref[r2 + 1][c2 + 1] - pref[r1][c2 + 1] - pref[r2 + 1][c1] + pref[r1][c1];
}

int64_t count_subarrays_with_sum(const std::vector<int> &a, int64_t target) {
  std::unordered_map<int64_t, int64_t> count{{0, 1}};
  int64_t sum = 0, result = 0;
  for (int x : a) {
    sum += x;
    int64_t needed = sum - target;
    if (auto it = count.find(needed); it != count.end()) {
      result += it->second;
    }
    count[sum]++;
  }
  return result;
}
\end{lstlisting}
\Needspace*{4\baselineskip}
\textbf{Example Usage}\par
\begin{lstlisting}[language=C++]
using namespace std;

int main() {
  vector<int> a{3, -1, 4, 1, 5};
  auto pref = prefix_sums(a);
  assert(range_sum(pref, 0, 4) == 12);  // Whole array.
  assert(range_sum(pref, 1, 3) == 4);   // -1 + 4 + 1.
  assert(range_sum(pref, 2, 2) == 4);   // Single element.
  assert(count_subarrays_with_sum(a, 4) == 2);
  assert(count_subarrays_with_sum(vector<int>{0, 0, 0}, 0) == 6);

  vector<vector<int>> grid{
      {1, 2, 3},
      {4, 5, 6},
  };
  auto pre2 = prefix_sums_2d(grid);
  assert(rect_sum(pre2, 0, 1, 1, 2) == 16);  // Rows 0-1 and columns 1-2.
  assert(rect_sum(pre2, 0, 0, 1, 2) == 21);  // Whole grid.
  assert(rect_sum(pre2, 1, 1, 1, 1) == 5);   // Single cell.
  return 0;
}
\end{lstlisting}
\subsection{Difference Array}
Applies many range-add updates to an array in linear total time using a difference array. Instead of immediately adding \inlinecode{delta} to every element of range $[{\inlinecodemath{lo}}, {\inlinecodemath{hi}}]$, we can add \inlinecode{delta} at \inlinecode{diff[lo]} and subtract it at \inlinecode{diff[hi + 1]}; a final prefix sum reconstructs the updated values. This idea extends to two dimensions: a rectangle update touches just the four corner cells of a difference grid by inclusion-exclusion, and a final two-dimensional prefix sum reconstructs the grid.

\begin{compactitem}
  \item \inlinecode{DifferenceArray(n)} constructs an initially zero array of size \inlinecode{n}.
  \item \inlinecode{add(lo, hi, delta)} adds \inlinecode{delta} to every position in range $[{\inlinecodemath{lo}}, {\inlinecodemath{hi}}]$.
  \item \inlinecode{build()} returns the final array after all range updates.
  \item \inlinecode{DifferenceArray2D(rows, cols)} constructs an initially zero grid with \inlinecode{rows} rows and \inlinecode{cols} columns.
  \item \inlinecode{add(r1, c1, r2, c2, delta)} adds \inlinecode{delta} to every cell of the rectangle with rows $[{\inlinecodemath{r1}}, {\inlinecodemath{r2}}]$ and columns $[{\inlinecodemath{c1}}, {\inlinecodemath{c2}}]$.
  \item \inlinecode{build()} returns the final grid after all rectangle updates.
\end{compactitem}

Overflow warning: All stored differences and reconstructed values must fit in \inlinecode{int64\_t}.

\Needspace*{3\baselineskip}
\textbf{Time Complexity}:
\begin{compactitem}
  \item $O(1)$ per call to \inlinecode{add()} of either version.
  \item $O(n)$ per call to \inlinecode{build()} of \inlinecode{DifferenceArray}, where $n$ is the array size.
  \item $O(R \cdot C)$ per call to \inlinecode{build()} of \inlinecode{DifferenceArray2D}, where $R$ and $C$ are the number of rows and columns, respectively.
\end{compactitem}

\Needspace*{3\baselineskip}
\textbf{Space Complexity}:
\begin{compactitem}
  \item $O(n)$ for storage and $O(n)$ for the array returned by \inlinecode{DifferenceArray::build()}.
  \item $O(R \cdot C)$ for storage and $O(R \cdot C)$ for the grid returned by \inlinecode{DifferenceArray2D::build()}.
\end{compactitem}

\begin{lstlisting}[language=C++]
#include <cassert>
#include <cstdint>
#include <vector>

class DifferenceArray {
  std::vector<int64_t> diff;

 public:
  explicit DifferenceArray(int n) {
    assert(n >= 0);
    diff.assign(n + 1, 0);
  }

  void add(int lo, int hi, int64_t delta) {
    assert(0 <= lo && lo <= hi && hi + 1 < static_cast<int>(diff.size()));
    diff[lo] += delta;  // Overflow warning.
    diff[hi + 1] -= delta;
  }

  std::vector<int64_t> build() const {
    std::vector<int64_t> res(static_cast<int>(diff.size()) - 1);
    int64_t cur = 0;
    for (int i = 0; i < static_cast<int>(res.size()); i++) {
      cur += diff[i];  // Overflow warning.
      res[i] = cur;
    }
    return res;
  }
};

class DifferenceArray2D {
  std::vector<std::vector<int64_t>> diff;

 public:
  DifferenceArray2D(int rows, int cols) {
    assert(rows >= 0 && cols >= 0);
    diff.assign(rows + 1, std::vector<int64_t>(cols + 1));
  }

  void add(int r1, int c1, int r2, int c2, int64_t delta) {
    int rows = static_cast<int>(diff.size()) - 1;
    int cols = static_cast<int>(diff[0].size()) - 1;
    assert(0 <= r1 && r1 <= r2 && r2 < rows);
    assert(0 <= c1 && c1 <= c2 && c2 < cols);
    diff[r1][c1] += delta;  // Overflow warning.
    diff[r1][c2 + 1] -= delta;
    diff[r2 + 1][c1] -= delta;
    diff[r2 + 1][c2 + 1] += delta;
  }

  std::vector<std::vector<int64_t>> build() const {
    int rows = static_cast<int>(diff.size()) - 1;
    int cols = static_cast<int>(diff[0].size()) - 1;
    std::vector<std::vector<int64_t>> res(rows, std::vector<int64_t>(cols));
    for (int i = 0; i < rows; i++) {
      for (int j = 0; j < cols; j++) {
        int64_t up = i > 0 ? res[i - 1][j] : 0;
        int64_t left = j > 0 ? res[i][j - 1] : 0;
        int64_t diag = i > 0 && j > 0 ? res[i - 1][j - 1] : 0;
        res[i][j] = diff[i][j] + up + left - diag;  // Overflow warning.
      }
    }
    return res;
  }
};
\end{lstlisting}
\Needspace*{4\baselineskip}
\textbf{Example Usage}\par
\begin{lstlisting}[language=C++]
#include <cassert>
using namespace std;

int main() {
  DifferenceArray d(5);
  d.add(0, 2, 2);   // Add 2 to indices 0-2.
  d.add(1, 4, 4);   // Add 4 to indices 1-4.
  d.add(2, 3, -1);  // Subtract 1 from indices 2-3.
  d.add(4, 4, 3);   // Single index update.
  assert((d.build() == vector<int64_t>{2, 6, 5, 3, 7}));

  DifferenceArray2D d2(3, 4);
  d2.add(0, 0, 1, 2, 1);  // Add 1 to rows 0-1 and columns 0-2.
  d2.add(1, 1, 2, 3, 2);  // Add 2 to rows 1-2 and columns 1-3.
  d2.add(2, 0, 2, 0, 5);  // Single cell update.
  vector<vector<int64_t>> expected2{
      {1, 1, 1, 0},
      {1, 3, 3, 2},
      {5, 2, 2, 2},
  };
  assert(d2.build() == expected2);
  return 0;
}
\end{lstlisting}
\subsection{Two Pointers and Sliding Window}
Uses two moving indices to avoid trying all pairs of positions. This technique has two common forms: opposite-end pointers over sorted arrays, and a sliding window over contiguous subarrays. For sorted 2-sum, if the current sum is too small, pairing the left value with any smaller right value also fails, so the left pointer can safely advance; if the sum is too large, the symmetric argument lets the right pointer retreat. Each step therefore discards an entire row or column of candidate pairs.

A sliding window applies the same idea when extending and shrinking a range changes its state monotonically. With nonnegative values, extending a window cannot decrease its sum and removing its leftmost value cannot increase it, so no shorter feasible window is skipped. For a bound such as at most $k$ distinct values, the left pointer advances only until the window becomes valid again. Since neither endpoint moves backward, each element enters and leaves the window at most once.

For fixed-length windows, a monotonic deque of indices can report every minimum or maximum in linear total time. It stores only candidates that could still become the answer, ordered from the current extreme at the front to the least useful value at the back. A new element removes every older value that it dominates, while expired indices are removed from the front. Each index is again pushed and popped at most once.

\begin{lstlisting}[language=C++]
#include <algorithm>
#include <cassert>
#include <cstdint>
#include <deque>
#include <functional>
#include <tuple>
#include <unordered_map>
#include <utility>
#include <vector>
\end{lstlisting}
The examples below cover sorted 2-sum/3-sum variants, finding the shortest subarray with sum at least a target when all values are nonnegative, and finding the longest subarray with at most $k$ distinct values. They also cover the minimum or maximum of every fixed-length window and an online dynamic programming recurrence.

\begin{compactitem}
  \item \inlinecode{two\_sum\_sorted(a, target)} returns two indices whose values sum to \inlinecode{target}, or $(-1, -1)$ if no such pair exists. The input array must be sorted.
  \item \inlinecode{three\_sum(a, target)} returns a tuple of three original indices whose values sum to \inlinecode{target}, or $(-1, -1, -1)$ if none exist.
  \item \inlinecode{min\_subarray\_at\_least(a, target)} returns a tuple (\inlinecode{length}, \inlinecode{lo}, \inlinecode{hi}), the minimum length and inclusive endpoints of a contiguous subarray with sum at least \inlinecode{target}, or length $-1$ if none exists. Values in \inlinecode{a} must be nonnegative. If \inlinecode{target} is nonpositive, it returns the empty subarray.
  \item \inlinecode{max\_subarray\_at\_most\_k\_distinct(a, k)} returns a tuple (\inlinecode{length}, \inlinecode{lo}, \inlinecode{hi}), the maximum length and inclusive endpoints of a contiguous subarray containing at most \inlinecode{k} distinct values. If \inlinecode{k} is nonpositive, it returns the empty subarray.
  \item \inlinecode{sliding\_window\_extrema(a, k, comp = std::less\textless{}\textgreater{})} returns the extreme value in each window of length \inlinecode{k}. With the default \inlinecode{less\textless{}\textgreater{}} comparator it returns minimums; passing \inlinecode{greater\textless{}\textgreater{}} returns maximums.
\end{compactitem}

\Needspace*{3\baselineskip}
\textbf{Time Complexity}:
\begin{compactitem}
  \item $O(n)$ per call to \inlinecode{two\_sum\_sorted()} and \inlinecode{min\_subarray\_at\_least()}, where $n$ is the array size.
  \item $O(n^{2})$ per call to \inlinecode{three\_sum()}.
  \item $O(n)$ expected per call to \inlinecode{max\_subarray\_at\_most\_k\_distinct()}.
  \item $O(n)$ per call to \inlinecode{sliding\_window\_extrema()}.
\end{compactitem}

\Needspace*{3\baselineskip}
\textbf{Space Complexity}:
\begin{compactitem}
  \item $O(1)$ auxiliary for \inlinecode{two\_sum\_sorted(a, target)}.
  \item $O(n)$ auxiliary for \inlinecode{three\_sum(a, target)}.
  \item $O(1)$ auxiliary for \inlinecode{min\_subarray\_at\_least(a, target)}.
  \item $O(k)$ auxiliary for \inlinecode{max\_subarray\_at\_most\_k\_distinct(a, k)}.
  \item $O(k)$ auxiliary and $O(n)$ for the result of \inlinecode{sliding\_window\_extrema()}.
\end{compactitem}

\begin{lstlisting}[language=C++]
std::pair<int, int> two_sum_sorted(const std::vector<int> &a, int64_t target) {
  for (int lo = 0, hi = static_cast<int>(a.size()) - 1; lo < hi;) {
    int64_t sum = static_cast<int64_t>(a[lo]) + a[hi];
    if (sum == target) {
      return {lo, hi};
    } else if (sum < target) {
      lo++;
    } else {
      hi--;
    }
  }
  return {-1, -1};
}

std::tuple<int, int, int> three_sum(const std::vector<int> &a, int64_t target) {
  int n = static_cast<int>(a.size());
  std::vector<std::pair<int, int>> sorted;
  sorted.reserve(a.size());
  for (int i = 0; i < n; i++) {
    sorted.emplace_back(a[i], i);
  }
  std::sort(sorted.begin(), sorted.end());
  for (int i = 0; i < n; i++) {
    int lo = i + 1, hi = n - 1;
    while (lo < hi) {
      int64_t sum = static_cast<int64_t>(sorted[i].first) + sorted[lo].first + sorted[hi].first;
      if (sum == target) {
        return {sorted[i].second, sorted[lo].second, sorted[hi].second};
      } else if (sum < target) {
        lo++;
      } else {
        hi--;
      }
    }
  }
  return {-1, -1, -1};
}

std::tuple<int, int, int> min_subarray_at_least(const std::vector<int> &a, int64_t target) {
  if (target <= 0) {
    return {0, 0, -1};
  }
  int best = static_cast<int>(a.size()) + 1, best_lo = 0, best_hi = -1;
  int64_t sum = 0;
  for (int lo = 0, hi = 0; hi < static_cast<int>(a.size()); hi++) {
    sum += a[hi];  // Overflow warning.
    while (sum >= target) {
      if (hi - lo + 1 < best) {
        best = hi - lo + 1;
        best_lo = lo;
        best_hi = hi;
      }
      sum -= a[lo++];
    }
  }
  if (best == static_cast<int>(a.size()) + 1) {
    return {-1, 0, -1};
  }
  return {best, best_lo, best_hi};
}

std::tuple<int, int, int> max_subarray_at_most_k_distinct(const std::vector<int> &a, int k) {
  if (k <= 0) {
    return {0, 0, -1};
  }
  std::unordered_map<int, int> count;
  int best = 0, best_lo = 0, best_hi = -1;
  for (int lo = 0, hi = 0; hi < static_cast<int>(a.size()); hi++) {
    count[a[hi]]++;
    while (static_cast<int>(count.size()) > k) {
      if (--count[a[lo]] == 0) {
        count.erase(a[lo]);
      }
      lo++;
    }
    if (best < hi - lo + 1) {
      best = hi - lo + 1;
      best_lo = lo;
      best_hi = hi;
    }
  }
  return {best, best_lo, best_hi};
}

template<typename T, typename Compare = std::less<>>
std::vector<T> sliding_window_extrema(const std::vector<T> &a, int k, Compare comp = Compare{}) {
  int n = static_cast<int>(a.size());
  assert(1 <= k && k <= n);
  std::deque<int> window;
  std::vector<T> res;
  res.reserve(n - k + 1);
  for (int i = 0; i < n; i++) {
    while (!window.empty() && !comp(a[window.back()], a[i])) {
      window.pop_back();
    }
    window.push_back(i);
    if (window.front() <= i - k) {
      window.pop_front();
    }
    if (i >= k - 1) {
      res.push_back(a[window.front()]);
    }
  }
  return res;
}
\end{lstlisting}
A monotone queue can be used to maintain the minimum or maximum value in an online sliding window. This is useful for dynamic programming recurrences where each transition may only come from one of the last $w$ states, such as $dp(i) = a_i + \min_{j \in [i - w, i - 1]} dp(j)$.

The queue stores candidate (\inlinecode{index}, \inlinecode{value}) pairs in monotone order. Expired indices are removed from the front, and dominated values are removed from the back before inserting a new candidate. An index is an ordered \inlinecode{int} timestamp used only for expiration: indices need not be contiguous or nonnegative, but they must be pushed in strictly increasing order, and expiration thresholds must not decrease.

\begin{compactitem}
  \item \inlinecode{MonotoneQueue\textless{}T\textgreater{}()} constructs an empty queue for minimum queries. Instantiate \inlinecode{MonotoneQueue\textless{}T, std::greater\textless{}T\textgreater{}\textgreater{}} for maximum queries.
  \item \inlinecode{empty()} returns whether the queue has no active candidates.
  \item \inlinecode{push(index, value)} inserts candidate \inlinecode{value} with timestamp \inlinecode{index}, removing dominated candidates from the back.
  \item \inlinecode{expire(first\_valid)} removes candidates whose indices are less than \inlinecode{first\_valid}.
  \item \inlinecode{top()} returns the best active (\inlinecode{index}, \inlinecode{value}) pair. The queue must be nonempty.
\end{compactitem}

\Needspace*{3\baselineskip}
\textbf{Time Complexity}:
\begin{compactitem}
  \item $O(n)$ for any sequence of $n$ calls to \inlinecode{push(index, value)} and \inlinecode{expire(first\_valid)}.
  \item $O(1)$ amortized per call to \inlinecode{push(index, value)} and \inlinecode{expire(first\_valid)}.
  \item $O(1)$ worst-case per call to \inlinecode{top()} and \inlinecode{empty()}.
\end{compactitem}

\textbf{Space Complexity}: $O(w)$ for storage of a sliding window of width $w$.

\begin{lstlisting}[language=C++]
template<typename T, typename Compare = std::less<T>>
class MonotoneQueue {
  std::deque<std::pair<int, T>> q;
  Compare better;

 public:
  bool empty() const { return q.empty(); }

  void push(int index, const T &value) {
    while (!q.empty() && !better(q.back().second, value)) {
      q.pop_back();
    }
    q.emplace_back(index, value);
  }

  void expire(int first_valid) {
    while (!q.empty() && q.front().first < first_valid) {
      q.pop_front();
    }
  }

  std::pair<int, T> top() const {
    assert(!q.empty());
    return q.front();
  }
};
\end{lstlisting}
\Needspace*{4\baselineskip}
\textbf{Example Usage}\par
\begin{lstlisting}[language=C++]
using namespace std;

int main() {
  vector<int> c{-3, -1, 2, 4, 8, 11};
  assert(two_sum_sorted(c, 10) == make_pair(1, 5));  // -1 + 11.
  assert(two_sum_sorted(c, 7) == make_pair(1, 4));   // -1 + 8.
  assert(two_sum_sorted(c, 100) == make_pair(-1, -1));

  vector<int> d{4, -1, 8, 2, -3, 11};
  auto [i, j, k] = three_sum(d, 9);
  assert(d[i] + d[j] + d[k] == 9);
  assert((three_sum(d, 30) == make_tuple(-1, -1, -1)));

  vector<int> a{2, 3, 1, 2, 4, 3};
  auto [length, lo, hi] = min_subarray_at_least(a, 7);
  assert(length == 2 && lo == 4 && hi == 5);  // [4, 3].
  auto [empty_length, empty_lo, empty_hi] = min_subarray_at_least(a, 0);
  assert(empty_length == 0 && empty_lo == 0 && empty_hi == -1);

  vector<int> b{1, 2, 1, 3, 4, 3, 5};
  auto [longest, longest_lo, longest_hi] = max_subarray_at_most_k_distinct(b, 2);
  assert(longest == 3 && longest_lo == 0 && longest_hi == 2);  // [1, 2, 1].

  vector<int> e{1, 3, -1, -3, 5, 3, 6, 7};
  assert((sliding_window_extrema(e, 3) == vector<int>{-1, -3, -3, -3, 3, 3}));
  assert((sliding_window_extrema(e, 3, greater<>()) == vector<int>{3, 3, 5, 5, 6, 7}));

  // Toy DP: dp[i] = f[i] + min(dp[j]) for j in [i - w, i - 1].
  // Scanning the window costs O(w) per state; the monotone queue makes it O(1) amortized.
  vector<int> f{4, 2, 7, 1, 3, 6};
  int w = 3;
  vector<int> dp(f.size());
  MonotoneQueue<int> mq;
  dp[0] = f[0];
  mq.push(0, dp[0]);
  for (int i = 1; i < static_cast<int>(f.size()); i++) {
    mq.expire(i - w);
    dp[i] = f[i] + mq.top().second;
    mq.push(i, dp[i]);
  }
  assert((dp == vector<int>{4, 6, 11, 5, 8, 11}));
  return 0;
}
\end{lstlisting}
\subsection{Monotonic Stack}
A monotonic stack maintains its elements in sorted order as values are pushed, popping away any entries that would violate the ordering. Scanning an array once while keeping such a stack of indices answers, for every position, where the nearest smaller or larger neighbor lies. These ``nearest smaller/greater'' relationships underlie many array problems, the most famous being the largest rectangle in a histogram. Each index is pushed and popped at most once, so a full scan runs in linear time. A largest all-zero submatrix can then be found by sweeping rows, turning each row into a histogram of consecutive zeros ending at that row.

For any histogram bar, the nearest strictly lower bars to its left and right are the first positions that prevent a rectangle at the current bar's height from extending farther. The widest such rectangle therefore spans the positions strictly between those boundaries. For a binary matrix, the height at each column counts consecutive zeros ending in the current row; every all-zero rectangle has some bottom row, so solving one histogram per row considers every possible rectangle.

All functions below take a vector \inlinecode{a} of $n$ comparable values and return a vector of $n$ indices. A ``less'' query uses a strictly smaller neighbor and a ``greater'' query uses a strictly larger one; changing the comparison from strict to non-strict (e.g. \inlinecode{\textgreater{}=} to \inlinecode{\textgreater{}}) toggles how ties are handled.

\begin{compactitem}
  \item \inlinecode{prev\_less(a)} returns, for each $i$, the largest index $j < i$ with $a[j] < a[i]$, or $-1$ if there's no such index.
  \item \inlinecode{next\_less(a)} returns, for each $i$, the smallest index $j > i$ with $a[j] < a[i]$, or $n$ if there's no such index.
  \item \inlinecode{prev\_greater(a)} returns, for each $i$, the largest index $j < i$ with $a[j] > a[i]$, or $-1$ if there's no such index.
  \item \inlinecode{next\_greater(a)} returns, for each $i$, the smallest index $j > i$ with $a[j] > a[i]$, or $n$ if there's no such index.
  \item \inlinecode{largest\_histogram\_rectangle(heights)} returns the maximum area and its inclusive left and right endpoints in a histogram where bars have width $1$ and nonnegative heights given by \inlinecode{heights}.
  \item \inlinecode{largest\_zero\_submatrix(a)} returns the maximum area and its inclusive top, left, bottom, and right boundaries among all-zero rectangular submatrices of the 0/1 matrix \inlinecode{a}.
\end{compactitem}

Overflow warning: Histogram areas must fit in the height value type.

\Needspace*{3\baselineskip}
\textbf{Time Complexity}:
\begin{compactitem}
  \item $O(n)$ per call to all one-dimensional functions, where $n$ is the size of the input.
  \item $O(R \cdot C)$ per call to \inlinecode{largest\_zero\_submatrix()}, where $R$ and $C$ are the matrix dimensions.
\end{compactitem}

\Needspace*{3\baselineskip}
\textbf{Space Complexity}:
\begin{compactitem}
  \item $O(n)$ auxiliary and $O(n)$ for the returned indices from each neighbor query.
  \item $O(n)$ auxiliary for \inlinecode{largest\_histogram\_rectangle()}.
  \item $O(C)$ auxiliary for \inlinecode{largest\_zero\_submatrix()}.
\end{compactitem}

\begin{lstlisting}[language=C++]
#include <cassert>
#include <stack>
#include <tuple>
#include <vector>

template<typename T>
std::vector<int> prev_less(const std::vector<T> &a) {
  int n = static_cast<int>(a.size());
  std::vector<int> res(n);
  std::stack<int> s;
  for (int i = 0; i < n; i++) {
    while (!s.empty() && a[s.top()] >= a[i]) {
      s.pop();
    }
    res[i] = s.empty() ? -1 : s.top();
    s.push(i);
  }
  return res;
}

template<typename T>
std::vector<int> next_less(const std::vector<T> &a) {
  int n = static_cast<int>(a.size());
  std::vector<int> res(n);
  std::stack<int> s;
  for (int i = n - 1; i >= 0; i--) {
    while (!s.empty() && a[s.top()] >= a[i]) {
      s.pop();
    }
    res[i] = s.empty() ? n : s.top();
    s.push(i);
  }
  return res;
}

template<typename T>
std::vector<int> prev_greater(const std::vector<T> &a) {
  int n = static_cast<int>(a.size());
  std::vector<int> res(n);
  std::stack<int> s;
  for (int i = 0; i < n; i++) {
    while (!s.empty() && a[s.top()] <= a[i]) {
      s.pop();
    }
    res[i] = s.empty() ? -1 : s.top();
    s.push(i);
  }
  return res;
}

template<typename T>
std::vector<int> next_greater(const std::vector<T> &a) {
  int n = static_cast<int>(a.size());
  std::vector<int> res(n);
  std::stack<int> s;
  for (int i = n - 1; i >= 0; i--) {
    while (!s.empty() && a[s.top()] <= a[i]) {
      s.pop();
    }
    res[i] = s.empty() ? n : s.top();
    s.push(i);
  }
  return res;
}

template<typename T>
std::tuple<T, int, int> largest_histogram_rectangle(const std::vector<T> &heights) {
  int n = static_cast<int>(heights.size());
  std::vector<int> left = prev_less(heights), right = next_less(heights);
  T best{};
  int lo = 0, hi = -1;
  for (int i = 0; i < n; i++) {
    T area = heights[i] * (right[i] - left[i] - 1);  // Overflow warning.
    if (area > best) {
      best = area;
      lo = left[i] + 1;
      hi = right[i] - 1;
    }
  }
  return {best, lo, hi};
}

std::tuple<int, int, int, int, int> largest_zero_submatrix(
    const std::vector<std::vector<char>> &a
) {
  if (a.empty()) {
    return {0, 0, 0, -1, -1};
  }
  int rows = static_cast<int>(a.size()), cols = static_cast<int>(a[0].size());
  int best = 0, rlo = 0, clo = 0, rhi = -1, chi = -1;
  std::vector<int> height(cols);
  for (int i = 0; i < rows; i++) {
    for (int j = 0; j < cols; j++) {
      height[j] = a[i][j] ? 0 : height[j] + 1;
    }
    auto [area, lo, hi] = largest_histogram_rectangle(height);
    if (area > best) {
      best = area;
      rlo = i - area / (hi - lo + 1) + 1;
      clo = lo;
      rhi = i;
      chi = hi;
    }
  }
  return {best, rlo, clo, rhi, chi};
}
\end{lstlisting}
\Needspace*{4\baselineskip}
\textbf{Example Usage}\par
\begin{lstlisting}[language=C++]
using namespace std;

int main() {
  vector<int> a{2, 1, 4, 3};
  assert((prev_less(a) == vector<int>{-1, -1, 1, 1}));
  assert((next_less(a) == vector<int>{1, 4, 3, 4}));
  assert((prev_greater(a) == vector<int>{-1, 0, -1, 2}));
  assert((next_greater(a) == vector<int>{2, 2, 4, 4}));

  vector<int> hist{2, 1, 5, 6, 2, 3};
  // The best histogram rectangle uses bars 2 and 3: min height 5 times width 2.
  assert((largest_histogram_rectangle(hist) == make_tuple(10, 2, 3)));
  assert((largest_histogram_rectangle(vector<int>{}) == make_tuple(0, 0, -1)));

  vector<vector<char>> grid{
      {1, 0, 1, 1, 0, 0},
      {1, 0, 0, 1, 0, 0},
      {0, 0, 0, 0, 0, 1},
      {1, 0, 0, 1, 0, 0},
      {1, 0, 1, 0, 0, 1}
  };
  // The largest zero rectangle has 3 rows and 2 columns.
  assert((largest_zero_submatrix(grid) == make_tuple(6, 1, 1, 3, 2)));
  assert((largest_zero_submatrix({}) == make_tuple(0, 0, 0, -1, -1)));
  return 0;
}
\end{lstlisting}
\subsection{Aggregate Queue}
Maintain a first-in, first-out queue that additionally answers, in $O(1)$ amortized time, the fold of an associative operation over all elements currently in the queue. This is the ``sliding window aggregation'' (SWAG) technique: by pushing new elements and popping old ones as a window advances, it reports an aggregate of every window in linear total time.

Unlike the monotonic deque used for sliding-window extrema (see 1.2.3), the operation need only be associative, not order-based, so the same structure handles min, max, sum, gcd, matrix products, and other monoids. It is sometimes called a ``monotonic queue,'' but that term usually denotes the monotonic deque of 1.2.3 (which handles only min/max); the two are different structures. The queue is built from two stacks: a back stack receives pushes while a front stack serves pops, and each stack stores running folds so the two halves combine in $O(1)$. Whenever the front stack empties, the back stack is flipped onto it; this costs $O(1)$ amortized per element, since each element is moved between stacks at most once.

The fold is defined by an associative \inlinecode{combine(a, b)} function. The default code below returns the ``min'' of the elements; for ``sum'', \inlinecode{combine(a, b)} should return \inlinecode{a + b}. The fold is taken in queue order (front to back), so non-commutative operations such as matrix products are also supported.

\begin{compactitem}
  \item \inlinecode{AggregateQueue\textless{}T\textgreater{}()} constructs an empty queue.
  \item \inlinecode{empty()} returns whether the queue is empty.
  \item \inlinecode{size()} returns the number of elements in the queue.
  \item \inlinecode{push(x)} appends \inlinecode{x} to the back of the queue.
  \item \inlinecode{pop()} removes the element at the front of the queue, which must be non-empty.
  \item \inlinecode{front()} returns the element at the front of the queue, which must be non-empty.
  \item \inlinecode{fold()} returns \inlinecode{combine()} applied to all elements in queue order; the queue must be non-empty.
\end{compactitem}

\Needspace*{3\baselineskip}
\textbf{Time Complexity}:
\begin{compactitem}
  \item $O(1)$ per call to the constructor, \inlinecode{empty()}, \inlinecode{size()}, \inlinecode{front()}, and \inlinecode{fold()}.
  \item $O(1)$ amortized per call to \inlinecode{push()} and \inlinecode{pop()}.
\end{compactitem}

\Needspace*{3\baselineskip}
\textbf{Space Complexity}:
\begin{compactitem}
  \item $O(n)$ for storage of the queue elements, where $n$ is the number of elements.
  \item $O(1)$ auxiliary per operation, amortized (a \inlinecode{pop()} may flip the back stack onto the front stack).
\end{compactitem}

\begin{lstlisting}[language=C++]
#include <algorithm>
#include <cassert>
#include <utility>
#include <vector>

template<typename T>
class AggregateQueue {
  static T combine(const T &a, const T &b) { return std::min(a, b); }

  // Each stack entry is (value, fold), where fold combines the value with all entries between it
  // and its end of the queue (the front for front_stack, the back for back_stack).
  std::vector<std::pair<T, T>> front_stack, back_stack;

 public:
  bool empty() const { return front_stack.empty() && back_stack.empty(); }
  int size() const { return static_cast<int>(front_stack.size() + back_stack.size()); }

  void push(const T &x) {
    T acc = back_stack.empty() ? x : combine(back_stack.back().second, x);
    back_stack.emplace_back(x, acc);
  }

  void pop() {
    assert(!empty());
    if (front_stack.empty()) {
      // Flip the back stack onto the front stack, reversing order so the oldest element ends up on
      // top, and recompute the front folds in queue order.
      while (!back_stack.empty()) {
        T x = back_stack.back().first;
        back_stack.pop_back();
        T acc = front_stack.empty() ? x : combine(x, front_stack.back().second);
        front_stack.emplace_back(x, acc);
      }
    }
    front_stack.pop_back();
  }

  T front() const {
    assert(!empty());
    return front_stack.empty() ? back_stack.front().first : front_stack.back().first;
  }

  T fold() const {
    assert(!empty());
    if (front_stack.empty()) {
      return back_stack.back().second;
    }
    if (back_stack.empty()) {
      return front_stack.back().second;
    }
    return combine(front_stack.back().second, back_stack.back().second);
  }
};
\end{lstlisting}
\Needspace*{4\baselineskip}
\textbf{Example Usage}\par
\begin{lstlisting}[language=C++]
using namespace std;

int main() {
  // Used directly as a FIFO queue with an O(1) min query.
  AggregateQueue<int> q;
  assert(q.empty());
  q.push(5);
  q.push(3);
  q.push(8);
  assert(!q.empty());
  assert(q.front() == 5 && q.size() == 3);
  assert(q.fold() == 3);  // min(5, 3, 8)
  q.pop();                // Removes 5.
  assert(q.fold() == 3);  // min(3, 8)
  q.pop();                // Removes 3.
  assert(q.fold() == 8);
  q.pop();  // Removes 8.
  assert(q.empty());

  // Canonical use: the minimum of every width-3 window, advancing the window by push/pop.
  vector<int> a{4, 2, 5, 1, 3, 6};
  AggregateQueue<int> w;
  vector<int> mins;
  for (int i = 0; i < static_cast<int>(a.size()); i++) {
    w.push(a[i]);
    if (w.size() > 3) {
      w.pop();
    }
    if (w.size() == 3) {
      mins.push_back(w.fold());
    }
  }
  assert((mins == vector<int>{2, 1, 1, 1}));  // min of [4,2,5], [2,5,1], [5,1,3], [1,3,6]
  return 0;
}
\end{lstlisting}

\section{\texorpdfstring{Dynamic Programming}{1.3 Dynamic Programming}}
\setcounter{section}{3}
\setcounter{subsection}{0}
\subsection{Maximum Subarray Sum (Kadane)}
Given an array of numbers, determine the maximum possible sum of any contiguous subarray. Kadane's algorithm scans the array while maintaining a nonnegative running sum. At each position, it adds the current value, records the largest sum seen, and resets the running sum to zero if it becomes negative. Discarding a negative sum is safe because it could only reduce every subarray extending it. This can be adapted to compute the maximal submatrix sum as well.

Two related scans maintain slightly different states. For a circular array, the best wrapped subarray is the whole-array sum minus the minimum-sum subarray. For a maximum-product subarray, a negative value swaps the roles of the minimum and maximum products ending at the previous position, so both products must be tracked.

\begin{compactitem}
  \item \inlinecode{max\_subarray\_sum(a)} returns the sum and inclusive endpoints (\inlinecode{sum}, \inlinecode{begin}, \inlinecode{end}) for the maximal-sum subarray of vector \inlinecode{a}. This implementation requires operators \inlinecode{+} and \inlinecode{\textless{}} on the value type. By convention, the empty subarray is allowed, so an input containing only negative values returns sum $0$ and endpoints $[0, -1]$.
  \item \inlinecode{max\_circular\_subarray\_sum(a)} returns a tuple (\inlinecode{sum}, \inlinecode{begin}, \inlinecode{end}) containing the maximal circular-subarray sum and its inclusive endpoints. If \inlinecode{begin} $\leq$ \inlinecode{end}, the result is an ordinary range; if \inlinecode{begin} \textgreater{} \inlinecode{end}, it wraps from \inlinecode{begin} through the end of \inlinecode{a} and continues through \inlinecode{end}. The empty subarray is allowed under the same convention as \inlinecode{max\_subarray\_sum()}.
  \item \inlinecode{max\_product\_subarray(a)} returns a tuple (\inlinecode{product}, \inlinecode{begin}, \inlinecode{end}) containing the product and inclusive endpoints of a nonempty maximal-product subarray, or product $0$ and endpoints $[0, -1]$ for an empty input.
  \item \inlinecode{max\_submatrix\_sum(a)} returns the sum and inclusive boundaries (\inlinecode{sum}, \inlinecode{r1}, \inlinecode{c1}, \inlinecode{r2}, \inlinecode{c2}) for the largest rectangular submatrix of a matrix \inlinecode{a}. This implementation requires operators \inlinecode{+} and \inlinecode{\textless{}} on the matrix value type. By convention, the empty submatrix is allowed, so a matrix containing only negative values returns sum $0$ with empty row and column intervals.
\end{compactitem}

Overflow warning: All resulting sums and products must fit in the value type.

\Needspace*{3\baselineskip}
\textbf{Time Complexity}:
\begin{compactitem}
  \item $O(n)$ per call to \inlinecode{max\_subarray\_sum()}, \inlinecode{max\_circular\_subarray\_sum()}, and \inlinecode{max\_product\_subarray()}, where $n$ is the size of \inlinecode{a}.
  \item $O(R \cdot C^{2})$ per call to \inlinecode{max\_submatrix\_sum()}, where $R$ is the number of rows and $C$ is the number of columns in the matrix.
\end{compactitem}

\Needspace*{3\baselineskip}
\textbf{Space Complexity}:
\begin{compactitem}
  \item $O(1)$ auxiliary for each subarray function.
  \item $O(R)$ auxiliary for \inlinecode{max\_submatrix\_sum()}.
\end{compactitem}

\begin{lstlisting}[language=C++]
#include <tuple>
#include <utility>
#include <vector>

template<typename T>
std::tuple<T, int, int> max_subarray_sum(const std::vector<T> &a) {
  int n = static_cast<int>(a.size());
  if (n == 0) {
    return {T{}, 0, -1};
  }
  int curr_begin = 0, begin = 0, end = -1;
  T sum = 0, max_sum = 0;
  for (int i = 0; i < n; i++) {
    sum += a[i];  // Overflow warning.
    if (sum < 0) {
      sum = 0;
      curr_begin = i + 1;
    } else if (max_sum < sum) {
      max_sum = sum;
      begin = curr_begin;
      end = i;
    }
  }
  return {max_sum, begin, end};
}

template<typename T>
std::tuple<T, int, int> max_circular_subarray_sum(const std::vector<T> &a) {
  auto [max_sum, begin, end] = max_subarray_sum(a);
  int n = static_cast<int>(a.size());
  if (n == 0) {
    return {max_sum, begin, end};
  }
  int curr_begin = 0, min_begin = 0, min_end = -1;
  T total = 0, min_sum = 0, curr_sum = 0;
  for (int i = 0; i < n; i++) {
    total += a[i];     // Overflow warning.
    curr_sum += a[i];  // Overflow warning.
    if (0 < curr_sum) {
      curr_sum = 0;
      curr_begin = i + 1;
    } else if (curr_sum < min_sum) {
      min_sum = curr_sum;
      min_begin = curr_begin;
      min_end = i;
    }
  }
  T wrapped_sum = total - min_sum;
  if (max_sum < wrapped_sum) {
    max_sum = wrapped_sum;
    begin = (min_end + 1) % n;
    end = (min_begin + n - 1) % n;
  }
  return {max_sum, begin, end};
}

template<typename T>
std::tuple<T, int, int> max_product_subarray(const std::vector<T> &a) {
  int n = static_cast<int>(a.size());
  if (n == 0) {
    return {T{}, 0, -1};
  }
  T max_end = a[0], min_end = a[0], best = a[0];
  int max_begin = 0, min_begin = 0, begin = 0, end = 0;
  for (int i = 1; i < n; i++) {
    if (a[i] < 0) {
      std::swap(max_end, min_end);
      std::swap(max_begin, min_begin);
    }
    T next_max = max_end * a[i];  // Overflow warning.
    T next_min = min_end * a[i];  // Overflow warning.
    if (next_max < a[i]) {
      max_end = a[i];
      max_begin = i;
    } else {
      max_end = next_max;
    }
    if (a[i] < next_min) {
      min_end = a[i];
      min_begin = i;
    } else {
      min_end = next_min;
    }
    if (best < max_end) {
      best = max_end;
      begin = max_begin;
      end = i;
    }
  }
  return {best, begin, end};
}

template<typename T>
std::tuple<T, int, int, int, int> max_submatrix_sum(const std::vector<std::vector<T>> &a) {
  if (a.empty() || a[0].empty()) {
    return {T{}, 0, 0, -1, -1};
  }
  int rows = static_cast<int>(a.size()), cols = static_cast<int>(a[0].size());
  int r1 = 0, c1 = 0, r2 = -1, c2 = -1;
  T max_sum = 0;
  for (int clo = 0; clo < cols; clo++) {
    std::vector<T> sums(rows);
    for (int chi = clo; chi < cols; chi++) {
      for (int i = 0; i < rows; i++) {
        sums[i] += a[i][chi];  // Overflow warning.
      }
      auto [sum, rlo, rhi] = max_subarray_sum(sums);
      if (max_sum < sum) {
        max_sum = sum;
        r1 = rlo;
        c1 = clo;
        r2 = rhi;
        c2 = chi;
      }
    }
  }
  return {max_sum, r1, c1, r2, c2};
}
\end{lstlisting}
\Needspace*{4\baselineskip}
\textbf{Example Usage}\par
\begin{lstlisting}[language=C++]
#include <cassert>
#include <iostream>
using namespace std;

int main() {
  {
    // All negative values, so the empty subarray is maximal.
    auto [sum, begin, end] = max_subarray_sum(vector<int>{-2, -1, -3});
    assert(sum == 0 && begin == 0 && end == -1);
  }
  {
    vector<int> a{-2, -1, -3, 4, -1, 2, 1, -5, 4};
    auto [sum, begin, end] = max_subarray_sum(a);
    assert(sum == 6);
    cout << "Maximal sum subarray:" << endl;
    for (int i = begin; i <= end; i++) {
      cout << a[i] << " ";
    }
    cout << endl;
  }
  {
    vector<int> a{5, -3, 5};
    auto [sum, begin, end] = max_circular_subarray_sum(a);
    assert(sum == 10 && begin == 2 && end == 0);
  }
  {
    vector<int> a{2, 3, -2, 4};
    auto [product, begin, end] = max_product_subarray(a);
    assert(product == 6 && begin == 0 && end == 1);
  }
  {
    vector<vector<int>> a{
        {0, -2, -7, 0, 5},
        {9, 2, -6, 2, -4},
        {-4, 1, -4, 1, 0},
        {-1, 8, 0, -2, 3},
    };
    auto [sum, r1, c1, r2, c2] = max_submatrix_sum(a);
    assert(sum == 15);
    cout << "\nMaximal sum submatrix:" << endl;
    for (int i = r1; i <= r2; i++) {
      for (int j = c1; j <= c2; j++) {
        cout << a[i][j] << " ";
      }
      cout << endl;
    }
  }
  return 0;
}
\end{lstlisting}
\begin{exampleoutput}
Maximal sum subarray:
4 -1 2 1

Maximal sum submatrix:
9 2
-4 1
-1 8
\end{exampleoutput}
\subsection{Longest Increasing Subsequence}
Given an array of elements, determine a longest subsequence where all elements are in strictly increasing order. The subsequence is not necessarily contiguous or unique, so only one such answer will be found. The answer is computed by scanning left to right while maintaining, for each length, the smallest possible tail value of an increasing subsequence of that length. Each element extends or improves one of these tails, located in $O(\log n)$ by binary search; predecessor links then reconstruct the subsequence. The comparator \inlinecode{comp} defines the ordering and defaults to \inlinecode{std::less\textless{}\textgreater{}}.

\begin{compactitem}
  \item \inlinecode{longest\_increasing\_subsequence(a, comp = std::less\textless{}\textgreater{}())} returns the indices of one longest strictly increasing subsequence of \inlinecode{a} according to \inlinecode{comp}, in order.
\end{compactitem}

\textbf{Time Complexity}: $O(n \log n)$ per call, where $n$ is the size of \inlinecode{a}.

\textbf{Space Complexity}: $O(n)$ auxiliary and $O(n)$ for the returned indices.

\begin{lstlisting}[language=C++]
#include <algorithm>
#include <functional>
#include <vector>

template<typename T, typename Compare = std::less<>>
std::vector<int> longest_increasing_subsequence(const std::vector<T> &a, Compare comp = Compare{}) {
  int n = static_cast<int>(a.size());
  if (n == 0) {
    return {};
  }
  int len = 0;
  // prev[i] = index of the predecessor of element i in its LIS (or -1 for the first LIS element).
  // tail[i] = index of the last element of the best known LIS of length i + 1.
  std::vector<int> prev(n), tail(n);
  for (int i = 0; i < n; i++) {
    // Find the tail to extend or improve. Comparing by value through the stored indices keeps tail
    // index-based for reconstruction. Switch to an upper-bound search for a non-decreasing result.
    int pos = static_cast<int>(
        std::lower_bound(
            tail.begin(), tail.begin() + len, i, [&](int t, int x) { return comp(a[t], a[x]); }
        ) -
        tail.begin()
    );
    if (pos == len) {
      len++;
    }
    prev[i] = pos > 0 ? tail[pos - 1] : -1;
    tail[pos] = i;
  }
  // Optional: reconstruct one longest increasing subsequence.
  std::vector<int> res(len);
  for (int i = tail[len - 1]; i != -1; i = prev[i]) {
    res[--len] = i;
  }
  return res;
}
\end{lstlisting}
\Needspace*{4\baselineskip}
\textbf{Example Usage}\par
\begin{lstlisting}[language=C++]
#include <cassert>
using namespace std;

int main() {
  vector<int> a{-2, -5, 1, 9, 10, 8, 11, 10, 13, 11};
  assert((longest_increasing_subsequence(a) == vector<int>{1, 2, 3, 4, 6, 8}));

  vector<int> b{1, 4, 3, 2};
  assert((longest_increasing_subsequence(b, greater<int>()) == vector<int>{1, 2, 3}));
  return 0;
}
\end{lstlisting}
\subsection{0-1 Knapsack}
Solves the 0-1 knapsack problem: given items with nonnegative integer weights and values, choose each item at most once to maximize total value without exceeding a capacity limit.

The one-dimensional dynamic programming array \inlinecode{dp[w]} stores the best value that can be achieved with total weight at most \inlinecode{w} after processing some prefix of the items. Iterating weights in decreasing order is what enforces the 0-1 constraint; it prevents the same item from being reused during the current item update.

\begin{compactitem}
  \item \inlinecode{knapsack\_01(weight, value, capacity)} returns a pair (\inlinecode{best\_value}, \inlinecode{items}) containing the maximum value and the selected item indices in increasing order, where item \inlinecode{i} has weight \inlinecode{weight[i]} and value \inlinecode{value[i]}. All weights must be nonnegative integers. Items with weight $0$ are allowed and can each be taken once. The take/skip decision for every item and capacity is stored to reconstruct the optimal subset.
\end{compactitem}

\textbf{Time Complexity}: $O(n \cdot W)$ per call, where $n$ is the number of items and $W$ is the capacity.

\textbf{Space Complexity}: $O(n \cdot W)$ auxiliary.

\begin{lstlisting}[language=C++]
#include <algorithm>
#include <cassert>
#include <cstdint>
#include <utility>
#include <vector>

std::pair<int64_t, std::vector<int>> knapsack_01(
    const std::vector<int> &weight, const std::vector<int64_t> &value, int capacity
) {
  int n = static_cast<int>(weight.size());
  assert(static_cast<int>(value.size()) == n && capacity >= 0);
  for (int i = 0; i < n; i++) {
    assert(weight[i] >= 0 && value[i] >= 0);
  }
  std::vector<int64_t> dp(capacity + 1);
  std::vector<std::vector<char>> take(n, std::vector<char>(capacity + 1));
  for (int i = 0; i < n; i++) {
    for (int w = capacity; w >= weight[i]; w--) {
      int64_t candidate = dp[w - weight[i]] + value[i];  // Overflow warning.
      if (dp[w] < candidate) {
        dp[w] = candidate;
        take[i][w] = true;
      }
    }
  }
  // Optional: reconstruct one optimal subset.
  std::vector<int> items;
  for (int i = n - 1, w = capacity; i >= 0; i--) {
    if (take[i][w]) {
      items.push_back(i);
      w -= weight[i];
    }
  }
  std::reverse(items.begin(), items.end());
  return {dp[capacity], items};
}
\end{lstlisting}
\Needspace*{4\baselineskip}
\textbf{Example Usage}\par
\begin{lstlisting}[language=C++]
#include <cassert>
using namespace std;

int main() {
  vector<int> weight{3, 4, 5, 9};
  vector<int64_t> value{4, 5, 7, 10};
  auto [best_value, items] = knapsack_01(weight, value, 8);
  assert(best_value == 11 && (items == vector<int>{0, 2}));
  return 0;
}
\end{lstlisting}
\subsection{Unbounded Knapsack and Coin Change}
Solves unbounded knapsack and two related coin change problems, where each item or coin may be used any number of times. In the maximum-value version, each item has a weight and value, and the goal is to maximize total value subject to a capacity limit. In coin change variants, the goal is either to count unordered ways to form a target sum or to minimize the number of coins used.

For unbounded maximum-value knapsack, capacities are processed in increasing order. Thus, when computing \inlinecode{dp[w]}, every smaller-capacity state \inlinecode{dp[w - weight[i]]} is already final and may itself contain item \inlinecode{i}, allowing unlimited reuse. For unordered coin change, coins are the outer loop so each multiset of coins is counted once, independent of ordering. For minimum coins, each target sum chooses the best previous sum \inlinecode{sum - coin}, and the predecessor coin reconstructs one optimal multiset.

\begin{compactitem}
  \item \inlinecode{unbounded\_knapsack(weight, value, capacity)} returns a pair (\inlinecode{best\_value}, \inlinecode{count}) containing the maximum value achievable with total weight at most \inlinecode{capacity} and \inlinecode{count[i]}, the number of copies selected of each item \inlinecode{i}.
  \item \inlinecode{count\_coin\_change(coins, target)} returns the number of unordered ways to make sum \inlinecode{target} using the given coin denominations.
  \item \inlinecode{min\_coin\_change(coins, target)} returns a pair (\inlinecode{count}, \inlinecode{used}), where \inlinecode{count} is the minimum number of coins needed to make sum \inlinecode{target}, and \inlinecode{used[i]} is the number of copies of \inlinecode{coins[i]} in one optimal multiset. If \inlinecode{target} cannot be formed, \inlinecode{count} is $-1$ and \inlinecode{used} is empty.
\end{compactitem}

\inlinecode{capacity} and \inlinecode{target} must be nonnegative integers. Weights and coin denominations must be positive. Coin denominations must also be distinct when counting ways; duplicate denominations are treated as different coin types and overcount identical multisets.

\Needspace*{3\baselineskip}
\textbf{Time Complexity}:
\begin{compactitem}
  \item $O(n \cdot W)$ per call to \inlinecode{unbounded\_knapsack()}, where $n$ is the number of items and $W$ is \inlinecode{capacity}.
  \item $O(c \cdot t)$ per call to \inlinecode{count\_coin\_change()} and \inlinecode{min\_coin\_change()}, where $c$ is the number of coin denominations and $t$ is the target.
\end{compactitem}

\Needspace*{3\baselineskip}
\textbf{Space Complexity}:
\begin{compactitem}
  \item $O(W)$ auxiliary and $O(n)$ for the returned \inlinecode{count} from \inlinecode{unbounded\_knapsack()}.
  \item $O(t)$ auxiliary for \inlinecode{count\_coin\_change()}.
  \item $O(t)$ auxiliary and $O(c)$ for the returned \inlinecode{used} from \inlinecode{min\_coin\_change()}.
\end{compactitem}

\begin{lstlisting}[language=C++]
#include <cassert>
#include <cstdint>
#include <utility>
#include <vector>

std::pair<int64_t, std::vector<int>> unbounded_knapsack(
    const std::vector<int> &weight, const std::vector<int64_t> &value, int capacity
) {
  int n = static_cast<int>(weight.size());
  assert(static_cast<int>(value.size()) == n && capacity >= 0);
  for (int w : weight) {
    assert(w > 0);
  }
  std::vector<int64_t> dp(capacity + 1);
  std::vector<int> choice(capacity + 1, -1);
  for (int w = 1; w <= capacity; w++) {
    for (int i = 0; i < n; i++) {
      if (weight[i] <= w) {
        int64_t candidate = dp[w - weight[i]] + value[i];  // Overflow warning.
        if (dp[w] < candidate) {
          dp[w] = candidate;
          choice[w] = i;
        }
      }
    }
  }
  // Optional: reconstruct one optimal multiset of items.
  std::vector<int> count(n);
  for (int w = capacity; choice[w] != -1; w -= weight[choice[w]]) {
    count[choice[w]]++;
  }
  return {dp[capacity], count};
}

int64_t count_coin_change(const std::vector<int> &coins, int target) {
  assert(target >= 0);
  std::vector<int64_t> dp(target + 1);
  dp[0] = 1;
  for (int coin : coins) {
    assert(coin > 0);
    for (int sum = coin; sum <= target; sum++) {
      dp[sum] += dp[sum - coin];  // Overflow warning.
    }
  }
  return dp[target];
}

std::pair<int, std::vector<int>> min_coin_change(const std::vector<int> &coins, int target) {
  assert(target >= 0);
  for (int coin : coins) {
    assert(coin > 0);
  }
  int inf = target + 1;
  std::vector<int> dp(target + 1, inf), choice(target + 1, -1);
  dp[0] = 0;
  for (int sum = 1; sum <= target; sum++) {
    for (int i = 0; i < static_cast<int>(coins.size()); i++) {
      if (coins[i] <= sum && dp[sum] > dp[sum - coins[i]] + 1) {
        dp[sum] = dp[sum - coins[i]] + 1;
        choice[sum] = i;
      }
    }
  }
  if (dp[target] == inf) {
    return {-1, {}};
  }
  // Optional: reconstruct one minimum-size multiset of coins.
  std::vector<int> used(coins.size());
  for (int sum = target; sum > 0; sum -= coins[choice[sum]]) {
    used[choice[sum]]++;
  }
  return {dp[target], used};
}
\end{lstlisting}
\Needspace*{4\baselineskip}
\textbf{Example Usage}\par
\begin{lstlisting}[language=C++]
#include <cassert>
using namespace std;

int main() {
  vector<int> weight{3, 4, 5};
  vector<int64_t> value{4, 5, 7};
  auto [best_value, count] = unbounded_knapsack(weight, value, 10);
  assert(best_value == 14 && (count == vector<int>{0, 0, 2}));

  vector<int> coins{1, 2, 5};
  assert(count_coin_change(coins, 5) == 4);  // 5, 2+2+1, 2+1+1+1, 1*5.

  auto [coin_count, used] = min_coin_change(coins, 11);
  assert(coin_count == 3 && (used == vector<int>{1, 0, 2}));

  auto [impossible, empty_used] = min_coin_change(vector<int>{4, 7}, 6);
  assert(impossible == -1 && empty_used.empty());
  return 0;
}
\end{lstlisting}
\subsection{Matrix Chain Parenthesization}
Given a fixed-order chain of compatible matrices, finds a parenthesization that minimizes the number of scalar multiplications needed to compute their product. The algorithm examines only the matrix dimensions; it does not multiply matrix entries. Matrix multiplication is associative, but multiplying an $a$-by-$b$ matrix with a $b$-by-$c$ matrix costs $abc$ scalar multiplications, so different parenthesizations can have very different costs. This is a canonical example of interval dynamic programming.

Write $d_i$ for \inlinecode{dimensions[i]}, so matrix $i$ has $d_i$ rows and $d_{i+1}$ columns. For a chain of matrices, $dp(l, r)$ is the minimum cost of multiplying matrices $l$ through $r$. Trying every final split $k$ gives the recurrence $dp(l, r) = \min_{l \le k < r}(dp(l, k) + dp(k + 1, r) + d_l d_{k+1} d_{r+1})$. The final product multiplies a $d_l$-by-$d_{k+1}$ matrix with a $d_{k+1}$-by-$d_{r+1}$ matrix. Processing intervals in increasing order of length ensures that both subintervals have already been solved. The same pattern applies whenever a contiguous interval is formed by combining two smaller intervals.

\begin{compactitem}
  \item \inlinecode{matrix\_chain\_order(dimensions)} returns a pair (\inlinecode{operations}, \inlinecode{parenthesization}) containing the minimum number of scalar multiplications and one optimal parenthesization. Matrix \inlinecode{i} has dimensions \inlinecode{dimensions[i]} by \inlinecode{dimensions[i + 1]}; all dimensions must be positive.
\end{compactitem}

Overflow warning: The minimum cost and all intermediate products must fit in \inlinecode{int64\_t}.

\textbf{Time Complexity}: $O(n^{3})$ per call, where $n$ is the number of matrices. The dimension values affect the returned cost but not the running time.

\textbf{Space Complexity}: $O(n^{2})$ auxiliary for the dynamic programming and split tables.

\begin{lstlisting}[language=C++]
#include <cassert>
#include <cstdint>
#include <string>
#include <utility>
#include <vector>

std::pair<int64_t, std::string> matrix_chain_order(const std::vector<int> &dimensions) {
  assert(dimensions.size() >= 2);
  for (int d : dimensions) {
    assert(d > 0);
  }
  int n = static_cast<int>(dimensions.size()) - 1;
  std::vector<std::vector<int64_t>> dp(n, std::vector<int64_t>(n));
  std::vector<std::vector<int>> split(n, std::vector<int>(n));
  for (int length = 2; length <= n; length++) {
    for (int lo = 0; lo + length <= n; lo++) {
      int hi = lo + length - 1;
      for (int k = lo; k < hi; k++) {
        int64_t candidate = dp[lo][k] + dp[k + 1][hi] +
                            static_cast<int64_t>(dimensions[lo]) * dimensions[k + 1] *
                                dimensions[hi + 1];  // Overflow warning.
        if (k == lo || candidate < dp[lo][hi]) {
          dp[lo][hi] = candidate;
          split[lo][hi] = k;
        }
      }
    }
  }
  // Optional: reconstruct one optimal parenthesization.
  auto rec = [&](auto &&rec, int lo, int hi) -> std::string {
    if (lo == hi) {
      return "A" + std::to_string(lo);
    }
    int k = split[lo][hi];
    return "(" + rec(rec, lo, k) + " * " + rec(rec, k + 1, hi) + ")";
  };
  return {dp[0][n - 1], rec(rec, 0, n - 1)};
}
\end{lstlisting}
\Needspace*{4\baselineskip}
\textbf{Example Usage}\par
\begin{lstlisting}[language=C++]
using namespace std;

int main() {
  auto [ops, parenthesization] = matrix_chain_order({40, 20, 30, 10, 30});
  assert(ops == 26000);
  assert(parenthesization == "((A0 * (A1 * A2)) * A3)");

  auto [single_ops, single_parenthesization] = matrix_chain_order({3, 5});
  assert(single_ops == 0 && single_parenthesization == "A0");
  return 0;
}
\end{lstlisting}
\subsection{Digit DP}
Digit dynamic programming counts integers in a large bounded range whose decimal representations satisfy a digit-based property. The concrete routine below counts integers with a specified digit sum and provides the standard template to adapt for other finite-state properties.

To count valid integers in $[0, x]$, process the digits of $x$ from left to right. The state records the current position, accumulated digit sum, and whether the chosen prefix still equals the prefix of $x$. Once that ``tight'' condition is false, the remaining choices depend only on the position and sum, so those states can be memoized. Subtracting the prefix counts for \inlinecode{hi} and \inlinecode{lo - 1} gives the answer for the requested range.

Leading zeros do not affect a digit sum, so every integer can be treated as a zero-padded string of the same length as the bound. For properties where leading zeros matter, add a \inlinecode{started} flag to the state. More generally, the sum can be replaced by any small state with a transition for appending a digit.

\begin{compactitem}
  \item \inlinecode{count\_digit\_sum(lo, hi, target)} returns the number of integers in the inclusive range $[{\inlinecodemath{lo}}, {\inlinecodemath{hi}}]$ whose decimal digits sum to \inlinecode{target}. The bounds and \inlinecode{target} must be nonnegative.
\end{compactitem}

\textbf{Time Complexity}: $O(d \cdot s)$ per call, where $d$ is the number of digits in \inlinecode{hi} and $s$ is \inlinecode{target}.

\textbf{Space Complexity}: $O(d \cdot s)$ auxiliary heap space and $O(d)$ call stack space.

\begin{lstlisting}[language=C++]
#include <cassert>
#include <cstdint>
#include <string>
#include <vector>

int64_t count_digit_sum(int64_t lo, int64_t hi, int target) {
  assert(0 <= lo && lo <= hi && target >= 0);
  auto count_up_to = [&](int64_t bound) {
    if (bound < 0) {
      return 0LL;
    }
    std::string digits = std::to_string(bound);
    int n = static_cast<int>(digits.size());
    if (target > 9 * n) {
      return 0LL;
    }
    std::vector<std::vector<int64_t>> memo(n, std::vector<int64_t>(target + 1, -1));
    auto rec = [&](auto &&rec, int pos, int sum, bool tight) -> int64_t {
      if (sum > target || sum + 9 * (n - pos) < target) {
        return 0;
      }
      if (pos == n) {
        return sum == target;
      }
      if (!tight && memo[pos][sum] != -1) {
        return memo[pos][sum];
      }
      int bound_digit = digits[pos] - '0';
      int limit = tight ? bound_digit : 9;
      int64_t res = 0;
      for (int digit = 0; digit <= limit; digit++) {
        res += rec(rec, pos + 1, sum + digit, tight && digit == bound_digit);
      }
      if (!tight) {
        memo[pos][sum] = res;
      }
      return res;
    };
    return rec(rec, 0, 0, true);
  };
  return count_up_to(hi) - count_up_to(lo - 1);
}
\end{lstlisting}
\Needspace*{4\baselineskip}
\textbf{Example Usage}\par
\begin{lstlisting}[language=C++]
#include <cassert>
using namespace std;

int main() {
  assert(count_digit_sum(0, 20, 2) == 3);   // 2, 11, and 20.
  assert(count_digit_sum(10, 30, 3) == 3);  // 12, 21, and 30.
  assert(count_digit_sum(0, 99, 9) == 10);
  assert(count_digit_sum(0, 0, 0) == 1);
  assert(count_digit_sum(0, 999, 28) == 0);
  return 0;
}
\end{lstlisting}
\subsection{Egg Dropping}
Given $e$ identical eggs and a building of $f$ floors, find the highest floor from which an egg survives a drop, using as few drops as possible in the worst case. An egg that survives can be reused; an egg that breaks is gone. With one egg, the only safe strategy is to climb one floor at a time, needing $f$ drops. With unlimited eggs, binary search needs $\lceil \log_2(f + 1) \rceil$ drops. The problem is what happens in between.

The natural recurrence minimizes over the drop floor and maximizes over the two outcomes, costing $O(e \cdot f^{2})$. Removing that minimization is a technique worth knowing beyond this puzzle: when the quantity being minimized is small and bounded, promote it to a state dimension and solve for the old one. So ask how many floors $d$ drops with $e$ eggs can cover, not how many drops $f$ floors need. A drop splits the searchable floors into those below it, the drop floor, and those above, giving $dp(d, e) = dp(d-1, e-1) + dp(d-1, e) + 1$ with nothing left to minimize. Each entry is then $O(1)$, and the answer is the smallest $d$ whose reach is at least $f$.

That table is $dp(d, e) = \sum_{i=1}^{e} \binom{d}{i}$ in the binomial coefficients of section 6.2.1, which is what keeps $d$ small: enough eggs sum to $2^d - 1$, while one egg leaves only $\binom{d}{1} = d$. Two eggs reach $91$ floors in 13 drops and $105$ in 14, the classic answer. The same inversion fits any question of how many tests distinguish $n$ cases, such as finding one defective item with a limited number of destructive tests.

\begin{compactitem}
  \item \inlinecode{min\_drops(eggs, floors)} returns the number of drops needed in the worst case, where \inlinecode{eggs} is positive and \inlinecode{floors} is nonnegative.
  \item \inlinecode{drop\_floor(eggs, floors)} returns the floor to drop from first under an optimal strategy, which is $dp(d-1, e-1) + 1$ floors above the bottom of the current range, or $0$ when no drop is needed.
\end{compactitem}

\textbf{Time Complexity}: $O(e \cdot d)$ per call, where $e$ is \inlinecode{eggs} and $d$ is the returned number of drops, which is at most $f$ and is $O(\log f)$ once $e \geq \log_2(f)$.

\textbf{Space Complexity}: $O(e)$ auxiliary.

\begin{lstlisting}[language=C++]
#include <cassert>
#include <cstdint>
#include <vector>

int min_drops(int eggs, int floors) {
  assert(eggs > 0 && floors >= 0);
  // reach[e] is the number of floors distinguishable by the current number of drops with e eggs.
  std::vector<int64_t> reach(eggs + 1);
  int drops = 0;
  while (reach[eggs] < floors) {
    drops++;
    for (int e = eggs; e > 0; e--) {  // Descending, so reach[e - 1] is still the previous row.
      reach[e] += reach[e - 1] + 1;
    }
  }
  return drops;
}

int drop_floor(int eggs, int floors) {
  assert(eggs > 0 && floors >= 0);
  int drops = min_drops(eggs, floors);
  if (drops == 0) {
    return 0;
  }
  std::vector<int64_t> reach(eggs + 1);
  for (int i = 0; i < drops - 1; i++) {  // Rebuild the table one drop short of the answer.
    for (int e = eggs; e > 0; e--) {
      reach[e] += reach[e - 1] + 1;
    }
  }
  return static_cast<int>(reach[eggs - 1]) + 1;
}
\end{lstlisting}
\Needspace*{4\baselineskip}
\textbf{Example Usage}\par
\begin{lstlisting}[language=C++]
#include <cassert>

int main() {
  assert(min_drops(1, 0) == 0);
  assert(min_drops(1, 10) == 10);  // One egg forces a linear scan.
  assert(min_drops(2, 10) == 4);
  assert(min_drops(2, 100) == 14);  // The classic two-egg, hundred-floor puzzle.
  assert(min_drops(3, 100) == 9);
  assert(min_drops(50, 1000) == 10);  // Enough eggs to binary search.

  // With two eggs and a hundred floors, the first drop is from floor 14.
  assert(drop_floor(2, 100) == 14);
  assert(drop_floor(1, 10) == 1);
  assert(drop_floor(2, 0) == 0);
  return 0;
}
\end{lstlisting}

\section{\texorpdfstring{Grid Algorithms}{1.4 Grid Algorithms}}
\setcounter{section}{4}
\setcounter{subsection}{0}
\subsection{Grid Transformations}
Transforms a rectangular grid by transposing, rotating, or reflecting its cells. Transposition exchanges rows and columns. Rotation is clockwise by default, while a negative angle rotates counter-clockwise; the angle must be a multiple of $90$ degrees. Together, rotations and reflections generate all eight symmetries of a square grid.

\begin{compactitem}
  \item \inlinecode{transpose(a)} returns the transpose of grid \inlinecode{a}, so each output cell $(c, r)$ equals input cell $(r, c)$.
  \item \inlinecode{transpose\_in\_place(a)} transposes the square grid \inlinecode{a} in place and returns a reference to it.
  \item \inlinecode{rotate(a, degrees = 90)} returns \inlinecode{a} rotated clockwise by \inlinecode{degrees}.
  \item \inlinecode{rotate\_in\_place(a, degrees = 90)} rotates \inlinecode{a} in place and returns a reference to it. The grid must be square for rotations by $90$ or $270$ degrees.
  \item \inlinecode{reflect\_horizontal(a)} returns \inlinecode{a} reflected across its horizontal axis, reversing the row order.
  \item \inlinecode{reflect\_vertical(a)} returns \inlinecode{a} reflected across its vertical axis, reversing every row.
\end{compactitem}

\textbf{Time Complexity}: $O(R \cdot C)$ per call, where $R$ and $C$ are the grid dimensions.

\Needspace*{3\baselineskip}
\textbf{Space Complexity}:
\begin{compactitem}
  \item $O(R \cdot C)$ for the returned grid from the non-in-place operations.
  \item $O(1)$ auxiliary for the in-place operations.
\end{compactitem}

\begin{lstlisting}[language=C++]
#include <algorithm>
#include <cassert>
#include <utility>
#include <vector>

using Grid = std::vector<std::vector<int>>;

Grid transpose(const Grid &a) {
  int rows = static_cast<int>(a.size());
  int cols = a.empty() ? 0 : static_cast<int>(a[0].size());
  Grid res(cols, std::vector<int>(rows));
  for (int r = 0; r < rows; r++) {
    for (int c = 0; c < cols; c++) {
      res[c][r] = a[r][c];
    }
  }
  return res;
}

Grid &transpose_in_place(Grid &a) {
  assert(a.empty() || a.size() == a[0].size());
  for (int r = 0; r < static_cast<int>(a.size()); r++) {
    for (int c = r + 1; c < static_cast<int>(a.size()); c++) {
      std::swap(a[r][c], a[c][r]);
    }
  }
  return a;
}

Grid rotate(const Grid &a, int degrees = 90) {
  assert(degrees % 90 == 0);
  int rows = static_cast<int>(a.size());
  int cols = a.empty() ? 0 : static_cast<int>(a[0].size());
  degrees = (degrees % 360 + 360) % 360;
  Grid res(degrees % 180 == 0 ? rows : cols, std::vector<int>(degrees % 180 == 0 ? cols : rows));
  for (int r = 0; r < rows; r++) {
    for (int c = 0; c < cols; c++) {
      if (degrees == 90) {
        res[c][rows - r - 1] = a[r][c];
      } else if (degrees == 180) {
        res[rows - r - 1][cols - c - 1] = a[r][c];
      } else if (degrees == 270) {
        res[cols - c - 1][r] = a[r][c];
      } else {
        res[r][c] = a[r][c];
      }
    }
  }
  return res;
}

Grid &rotate_in_place(Grid &a, int degrees = 90) {
  assert(degrees % 90 == 0);
  degrees = (degrees % 360 + 360) % 360;
  assert(degrees % 180 == 0 || a.empty() || a.size() == a[0].size());
  if (degrees == 90) {
    transpose_in_place(a);
    for (auto &row : a) {
      std::reverse(row.begin(), row.end());
    }
  } else if (degrees == 180) {
    std::reverse(a.begin(), a.end());
    for (auto &row : a) {
      std::reverse(row.begin(), row.end());
    }
  } else if (degrees == 270) {
    transpose_in_place(a);
    std::reverse(a.begin(), a.end());
  }
  return a;
}

Grid reflect_horizontal(Grid a) {
  std::reverse(a.begin(), a.end());
  return a;
}

Grid reflect_vertical(Grid a) {
  for (auto &row : a) {
    std::reverse(row.begin(), row.end());
  }
  return a;
}
\end{lstlisting}
\Needspace*{4\baselineskip}
\textbf{Example Usage}\par
\begin{lstlisting}[language=C++]
#include <cassert>
using namespace std;

int main() {
  Grid a{{1, 2, 3}, {4, 5, 6}};
  assert((transpose(a) == Grid{{1, 4}, {2, 5}, {3, 6}}));
  assert((rotate(a) == Grid{{4, 1}, {5, 2}, {6, 3}}));
  assert((rotate(a, -90) == Grid{{3, 6}, {2, 5}, {1, 4}}));
  assert((rotate(a, 180) == Grid{{6, 5, 4}, {3, 2, 1}}));
  assert((reflect_horizontal(a) == Grid{{4, 5, 6}, {1, 2, 3}}));
  assert((reflect_vertical(a) == Grid{{3, 2, 1}, {6, 5, 4}}));

  Grid square{{1, 2, 3}, {4, 5, 6}, {7, 8, 9}};
  for (int degrees = -360; degrees <= 360; degrees += 90) {
    Grid rotated = square;
    assert(rotate_in_place(rotated, degrees) == rotate(square, degrees));
  }
  Grid transposed = square;
  assert(transpose_in_place(transposed) == transpose(square));

  Grid rectangular = a;
  assert(rotate_in_place(rectangular, 180) == rotate(a, 180));
  return 0;
}
\end{lstlisting}
\subsection{Grid Traversal}
Traverses a rectangular grid as an implicit unweighted graph: each unblocked cell is a node joined to its unblocked (four-way by default) orthogonal neighbors. Depth-first flood fill finds connected regions, while breadth-first search finds shortest distances because every move has unit length. Seeding the BFS with several cells computes the distance to the nearest source in one pass.

\begin{compactitem}
  \item \inlinecode{inside(r, c, rows, cols)} returns whether $0 \leq {\inlinecodemath{r}} < {\inlinecodemath{rows}}$ and $0 \leq {\inlinecodemath{c}} < {\inlinecodemath{cols}}$.
  \item \inlinecode{adj4(r, c)} returns the four orthogonally adjacent coordinates without filtering those outside the grid. For eight-way movement, also include the four diagonal neighbors.
  \item \inlinecode{component\_area\_perimeter(blocked, r, c)} uses flood fill to return the area and perimeter of the component containing (\inlinecode{r}, \inlinecode{c}). If the starting cell is blocked, both values are $0$.
  \item \inlinecode{grid\_components(blocked)} uses repeated flood fills to return the connected components of cells marked zero, with each cell represented by a pair (\inlinecode{r}, \inlinecode{c}). Movement is allowed in the four orthogonal directions.
  \item \inlinecode{grid\_bfs(blocked, sources)} returns the pair (\inlinecode{dist}, \inlinecode{pred}) from every cell in \inlinecode{sources}. Each source must be an unblocked cell represented by a pair (\inlinecode{r}, \inlinecode{c}). Passing one source gives ordinary single-source BFS. Blocked and unreachable cells have distance $-1$.
  \item \inlinecode{get\_path(pred, tr, tc)} uses the predecessor grid returned by \inlinecode{grid\_bfs()} to return the path from a nearest source to (\inlinecode{tr}, \inlinecode{tc}), including both endpoints, or an empty vector if the target is blocked or unreachable.
  \item \inlinecode{spiral\_order(a)} returns the cells of grid \inlinecode{a} in clockwise spiral order, starting at the top-left corner and walking inward.
\end{compactitem}

\Needspace*{3\baselineskip}
\textbf{Time Complexity}:
\begin{compactitem}
  \item $O(R \cdot C)$ per call to \inlinecode{spiral\_order()}.
  \item $O(R \cdot C)$ per call to \inlinecode{component\_area\_perimeter()}, \inlinecode{grid\_components()}, or \inlinecode{grid\_bfs()}, where $R$ and $C$ are the grid dimensions.
  \item $O(p)$ per call to \inlinecode{get\_path()}, where $p$ is the number of cells in the returned path.
\end{compactitem}

\Needspace*{3\baselineskip}
\textbf{Space Complexity}:
\begin{compactitem}
  \item $O(R \cdot C)$ auxiliary, including the recursion stack used by \inlinecode{component\_area\_perimeter()} and \inlinecode{grid\_components()}.
  \item $O(R \cdot C)$ for the returned \inlinecode{dist} and \inlinecode{pred}, and $O(p)$ for the path returned by \inlinecode{get\_path()}.
\end{compactitem}

\begin{lstlisting}[language=C++]
#include <algorithm>
#include <array>
#include <cassert>
#include <cstddef>
#include <queue>
#include <utility>
#include <vector>

bool inside(int r, int c, int rows, int cols) {
  return 0 <= r && r < rows && 0 <= c && c < cols;
}

std::array<std::pair<int, int>, 4> adj4(int r, int c) {
  return {{{r - 1, c}, {r + 1, c}, {r, c - 1}, {r, c + 1}}};
}

template<typename Fn>
void grid_dfs(
    const std::vector<std::vector<char>> &blocked, int r, int c,
    std::vector<std::vector<char>> &visit, const Fn &f
) {
  int rows = static_cast<int>(blocked.size());
  int cols = static_cast<int>(blocked[0].size());
  visit[r][c] = true;
  f(r, c);
  for (auto [r2, c2] : adj4(r, c)) {
    if (inside(r2, c2, rows, cols) && !blocked[r2][c2] && !visit[r2][c2]) {
      grid_dfs(blocked, r2, c2, visit, f);
    }
  }
}

std::pair<int, int> component_area_perimeter(
    const std::vector<std::vector<char>> &blocked, int r, int c
) {
  int rows = static_cast<int>(blocked.size());
  int cols = blocked.empty() ? 0 : static_cast<int>(blocked[0].size());
  assert(inside(r, c, rows, cols));
  if (blocked[r][c]) {
    return {0, 0};
  }
  int area = 0, perimeter = 0;
  std::vector<std::vector<char>> visit(rows, std::vector<char>(cols));
  grid_dfs(blocked, r, c, visit, [&](int r, int c) {
    area++;
    for (auto [r2, c2] : adj4(r, c)) {
      perimeter += !inside(r2, c2, rows, cols) || blocked[r2][c2];
    }
  });
  return {area, perimeter};
}

std::vector<std::vector<std::pair<int, int>>> grid_components(
    const std::vector<std::vector<char>> &blocked
) {
  int rows = static_cast<int>(blocked.size());
  int cols = blocked.empty() ? 0 : static_cast<int>(blocked[0].size());
  std::vector<std::vector<std::pair<int, int>>> components;
  std::vector<std::vector<char>> visit(rows, std::vector<char>(cols));
  for (int r = 0; r < rows; r++) {
    for (int c = 0; c < cols; c++) {
      if (blocked[r][c] || visit[r][c]) {
        continue;
      }
      components.emplace_back();
      grid_dfs(blocked, r, c, visit, [&](int r, int c) { components.back().emplace_back(r, c); });
    }
  }
  return components;
}

auto grid_bfs(
    const std::vector<std::vector<char>> &blocked, const std::vector<std::pair<int, int>> &sources
) {
  int rows = static_cast<int>(blocked.size());
  int cols = blocked.empty() ? 0 : static_cast<int>(blocked[0].size());
  std::vector<std::vector<int>> dist(rows, std::vector<int>(cols, -1));
  std::vector<std::vector<std::pair<int, int>>> pred(
      rows, std::vector<std::pair<int, int>>(cols, {-1, -1})
  );
  std::queue<std::pair<int, int>> q;
  for (auto [r, c] : sources) {
    assert(inside(r, c, rows, cols) && !blocked[r][c]);
    if (dist[r][c] == -1) {
      dist[r][c] = 0;
      pred[r][c] = {r, c};
      q.emplace(r, c);
    }
  }
  while (!q.empty()) {
    auto [r, c] = q.front();
    q.pop();
    for (auto [r2, c2] : adj4(r, c)) {
      if (inside(r2, c2, rows, cols) && !blocked[r2][c2] && dist[r2][c2] == -1) {
        dist[r2][c2] = dist[r][c] + 1;
        pred[r2][c2] = {r, c};
        q.emplace(r2, c2);
      }
    }
  }
  return std::pair{std::move(dist), std::move(pred)};
}

std::vector<std::pair<int, int>> get_path(
    const std::vector<std::vector<std::pair<int, int>>> &pred, int tr, int tc
) {
  int rows = static_cast<int>(pred.size());
  int cols = pred.empty() ? 0 : static_cast<int>(pred[0].size());
  assert(inside(tr, tc, rows, cols));
  if (pred[tr][tc].first == -1) {
    return {};
  }
  std::vector<std::pair<int, int>> path;
  std::pair cur{tr, tc};
  while (pred[cur.first][cur.second] != cur) {
    path.push_back(cur);
    cur = pred[cur.first][cur.second];
  }
  path.push_back(cur);
  std::reverse(path.begin(), path.end());
  return path;
}

template<typename T>
std::vector<T> spiral_order(const std::vector<std::vector<T>> &a) {
  int rows = static_cast<int>(a.size());
  int cols = a.empty() ? 0 : static_cast<int>(a[0].size());
  std::vector<T> res;
  res.reserve(static_cast<size_t>(rows) * cols);
  // Peel one ring at a time. Guard the last 2 passes so single remaining rows/cols don't repeat.
  for (int top = 0, bottom = rows - 1, left = 0, right = cols - 1; top <= bottom && left <= right;
       top++, bottom--, left++, right--) {
    for (int c = left; c <= right; c++) {
      res.push_back(a[top][c]);
    }
    for (int r = top + 1; r <= bottom; r++) {
      res.push_back(a[r][right]);
    }
    if (top < bottom) {
      for (int c = right - 1; c >= left; c--) {
        res.push_back(a[bottom][c]);
      }
    }
    if (left < right) {
      for (int r = bottom - 1; r > top; r--) {
        res.push_back(a[r][left]);
      }
    }
  }
  return res;
}
\end{lstlisting}
\Needspace*{4\baselineskip}
\textbf{Example Usage}\par
\begin{lstlisting}[language=C++]
#include <cassert>
using namespace std;

int main() {
  vector<vector<char>> blocked{
      {0, 0, 1, 0},
      {1, 0, 1, 0},
      {0, 0, 1, 0},
  };
  auto components = grid_components(blocked);
  assert(components.size() == 2);
  assert(components[0].size() == 5 && components[1].size() == 3);
  assert((component_area_perimeter(blocked, 0, 0) == pair{5, 12}));
  assert((component_area_perimeter(blocked, 0, 2) == pair{0, 0}));

  auto [dist, pred] = grid_bfs(blocked, {{0, 0}, {2, 3}});
  assert((dist == vector<vector<int>>{{0, 1, -1, 2}, {-1, 2, -1, 1}, {4, 3, -1, 0}}));
  assert((get_path(pred, 2, 1) == vector<pair<int, int>>{{0, 0}, {0, 1}, {1, 1}, {2, 1}}));
  assert((get_path(pred, 0, 3) == vector<pair<int, int>>{{2, 3}, {1, 3}, {0, 3}}));
  assert(get_path(pred, 0, 2).empty());

  vector<vector<int>> g{
      {1, 2, 3, 4},
      {5, 6, 7, 8},
      {9, 10, 11, 12},
  };
  assert((spiral_order(g) == vector<int>{1, 2, 3, 4, 8, 12, 11, 10, 9, 5, 6, 7}));
  assert((spiral_order(vector<vector<int>>{{1, 2, 3}}) == vector<int>{1, 2, 3}));
  assert((spiral_order(vector<vector<int>>{{1}, {2}, {3}}) == vector<int>{1, 2, 3}));
  return 0;
}
\end{lstlisting}
\subsection{Grid DP}
Solves basic dynamic programming problems on a rectangular grid. In grid DP, each cell represents a state and transitions usually come from already-computed neighboring cells. Processing rows from top to bottom and columns from left to right makes predecessors above and to the left available when a state is computed.

\begin{compactitem}
  \item \inlinecode{count\_grid\_paths(blocked)} returns the number of paths from the upper-left cell to the lower-right cell, moving only down or right and avoiding blocked cells marked nonzero. The state $dp(r, c)$ counts paths to cell $(r, c)$ from the counts above and left. Only the current row is retained.
  \item \inlinecode{min\_cost\_grid\_path(cost)} returns a pair (\inlinecode{min\_cost}, \inlinecode{path}) containing the minimum cost of a down/right path from the upper-left cell to the lower-right cell and the cells of one optimal path in order. The state $dp(r, c)$ stores the cheapest cost to cell $(r, c)$. Only the current row of costs is retained, while predecessor directions reconstruct one optimal path. If the grid is empty, \inlinecode{min\_cost} is $0$ and \inlinecode{path} is empty.
  \item \inlinecode{largest\_one\_square(a)} returns a tuple (\inlinecode{side}, \inlinecode{top}, \inlinecode{left}) describing a largest all-nonzero square in \inlinecode{a}. The state $dp(c)$ is the largest square ending at the current cell in column $c$; a nonzero cell extends the minimum of the states above, left, and diagonally above-left. If no such square exists, the function returns $(0, -1, -1)$.
\end{compactitem}

Overflow warning: Path counts and path costs must fit in \inlinecode{int64\_t}.

\textbf{Time Complexity}: $O(R \cdot C)$ per call, where $R$ and $C$ are the grid dimensions.

\Needspace*{3\baselineskip}
\textbf{Space Complexity}:
\begin{compactitem}
  \item $O(C)$ auxiliary for \inlinecode{count\_grid\_paths(blocked)}.
  \item $O(R \cdot C)$ auxiliary for predecessor directions, $O(C)$ for path costs, and $O(R + C)$ for the returned path from \inlinecode{min\_cost\_grid\_path(cost)}.
  \item $O(C)$ auxiliary for \inlinecode{largest\_one\_square(a)}.
\end{compactitem}

\begin{lstlisting}[language=C++]
#include <algorithm>
#include <cstdint>
#include <tuple>
#include <utility>
#include <vector>

int64_t count_grid_paths(const std::vector<std::vector<char>> &blocked) {
  int rows = static_cast<int>(blocked.size());
  int cols = blocked.empty() ? 0 : static_cast<int>(blocked[0].size());
  if (rows == 0 || cols == 0 || blocked[0][0] || blocked[rows - 1][cols - 1]) {
    return 0;
  }
  std::vector<int64_t> dp(cols);
  dp[0] = 1;
  for (int r = 0; r < rows; r++) {
    for (int c = 0; c < cols; c++) {
      if (blocked[r][c]) {
        dp[c] = 0;
      } else if (c > 0) {
        dp[c] += dp[c - 1];  // Overflow warning.
      }
    }
  }
  return dp[cols - 1];
}

std::pair<int64_t, std::vector<std::pair<int, int>>> min_cost_grid_path(
    const std::vector<std::vector<int>> &cost
) {
  int rows = static_cast<int>(cost.size());
  int cols = cost.empty() ? 0 : static_cast<int>(cost[0].size());
  if (rows == 0 || cols == 0) {
    return {0, {}};
  }
  std::vector<int64_t> dp(cols);
  std::vector<std::vector<char>> pred(rows, std::vector<char>(cols));
  dp[0] = cost[0][0];
  for (int r = 0; r < rows; r++) {
    for (int c = 0; c < cols; c++) {
      if (r == 0 && c == 0) {
        continue;
      }
      if (r > 0 && (c == 0 || dp[c] <= dp[c - 1])) {
        dp[c] += cost[r][c];  // Overflow warning.
        pred[r][c] = 'U';
      } else {
        dp[c] = dp[c - 1] + cost[r][c];
        pred[r][c] = 'L';
      }
    }
  }
  // Optional: reconstruct one optimal path.
  std::vector<std::pair<int, int>> path;
  for (int r = rows - 1, c = cols - 1;;) {
    path.emplace_back(r, c);
    if (r == 0 && c == 0) {
      break;
    }
    if (pred[r][c] == 'U') {
      r--;
    } else {
      c--;
    }
  }
  std::reverse(path.begin(), path.end());
  return {dp.back(), path};
}

std::tuple<int, int, int> largest_one_square(const std::vector<std::vector<char>> &a) {
  int rows = static_cast<int>(a.size());
  int cols = a.empty() ? 0 : static_cast<int>(a[0].size());
  std::vector<int> dp(cols);
  int best = 0, top = -1, left = -1;
  for (int r = 0; r < rows; r++) {
    int diagonal = 0;
    for (int c = 0; c < cols; c++) {
      int above = dp[c];
      dp[c] = a[r][c] ? 1 + std::min({above, c == 0 ? 0 : dp[c - 1], diagonal}) : 0;
      diagonal = above;
      if (dp[c] > best) {
        best = dp[c];
        top = r - best + 1;
        left = c - best + 1;
      }
    }
  }
  return {best, top, left};
}
\end{lstlisting}
\Needspace*{4\baselineskip}
\textbf{Example Usage}\par
\begin{lstlisting}[language=C++]
#include <cassert>
using namespace std;

int main() {
  vector<vector<char>> blocked{
      {0, 0, 0},
      {0, 1, 0},
      {0, 0, 0},
  };
  assert(count_grid_paths(blocked) == 2);

  vector<vector<int>> cost{
      {1, 3, 1},
      {1, 5, 1},
      {4, 2, 1},
  };
  auto [min_cost, path] = min_cost_grid_path(cost);
  assert(min_cost == 7 && (path == vector<pair<int, int>>{{0, 0}, {0, 1}, {0, 2}, {1, 2}, {2, 2}}));

  vector<vector<char>> binary{
      {1, 0, 1, 0, 0},
      {1, 0, 1, 1, 1},
      {1, 1, 1, 1, 1},
      {1, 0, 1, 1, 1},
  };
  assert((largest_one_square(binary) == tuple{3, 1, 2}));
  assert((largest_one_square({{0, 0}}) == tuple{0, -1, -1}));
  return 0;
}
\end{lstlisting}

\section{\texorpdfstring{Greedy and Scheduling}{1.5 Greedy and Scheduling}}
\setcounter{section}{5}
\setcounter{subsection}{0}
\subsection{Fractional Knapsack}
Solves the fractional knapsack problem: given items with positive integer weights and nonnegative integer values, choose any fraction of each item to maximize total value without exceeding a capacity limit.

The greedy algorithm processes items in decreasing order of value per unit weight, taking each item in full until only part of the next item fits. This is optimal because replacing any weight taken from a lower-density item with the same weight from a higher-density item cannot decrease the total value. Unlike 0-1 knapsack, allowing fractions makes each such exchange feasible.

\begin{compactitem}
  \item \inlinecode{fractional\_knapsack(weight, value, capacity)} returns a pair (\inlinecode{best\_value}, \inlinecode{fraction}) containing the maximum value and the fraction of each input item taken. Item \inlinecode{i} has weight \inlinecode{weight[i]} and value \inlinecode{value[i]}, \inlinecode{capacity} is the maximum total weight, and \inlinecode{fraction[i]} is in $[0, 1]$.
\end{compactitem}

\textbf{Time Complexity}: $O(n \log n)$ per call due to sorting, where $n$ is the number of items.

\textbf{Space Complexity}: $O(n)$ auxiliary and $O(n)$ for the returned fractions.

\begin{lstlisting}[language=C++]
#include <algorithm>
#include <cassert>
#include <cstdint>
#include <numeric>
#include <utility>
#include <vector>

std::pair<double, std::vector<double>> fractional_knapsack(
    const std::vector<int> &weight, const std::vector<int> &value, int capacity
) {
  int n = static_cast<int>(weight.size());
  assert(static_cast<int>(value.size()) == n && capacity >= 0);
  for (int i = 0; i < n; i++) {
    assert(weight[i] > 0 && value[i] >= 0);
  }
  std::vector<int> order(n);
  std::iota(order.begin(), order.end(), 0);
  std::sort(order.begin(), order.end(), [&](int i, int j) {
    int64_t lhs = static_cast<int64_t>(value[i]) * weight[j];
    int64_t rhs = static_cast<int64_t>(value[j]) * weight[i];
    return lhs != rhs ? lhs > rhs : i < j;
  });
  double best_value = 0;
  std::vector<double> fraction(n);
  for (int i : order) {
    int taken = std::min(capacity, weight[i]);
    fraction[i] = static_cast<double>(taken) / weight[i];
    best_value += fraction[i] * value[i];
    capacity -= taken;
    if (capacity == 0) {
      break;
    }
  }
  return {best_value, fraction};
}
\end{lstlisting}
\Needspace*{4\baselineskip}
\textbf{Example Usage}\par
\begin{lstlisting}[language=C++]
#include <cmath>
using namespace std;

int main() {
  vector<int> weight{10, 20, 30};
  vector<int> value{60, 100, 120};
  auto [best_value, fraction] = fractional_knapsack(weight, value, 50);
  assert(fabs(best_value - 240) < 1e-9 && (fraction == vector<double>{1, 1, 2.0 / 3}));
  return 0;
}
\end{lstlisting}
\subsection{Interval Scheduling}
Selects the maximum number of non-overlapping intervals using the classic earliest-finish-time greedy algorithm. This is the unweighted interval scheduling problem, also known as activity selection: every interval has the same value, so the objective is to choose as many compatible intervals as possible.

The greedy choice is safe because among all intervals that could be chosen first, taking one with the earliest finish time leaves the most room for the remaining intervals. In any optimal solution, the first chosen interval can be exchanged for an earliest-finishing compatible interval without reducing the number of intervals selected. Repeating this argument after each choice proves the greedy algorithm optimal. The weighted version in 1.5.3 needs dynamic programming instead.

Intervals are represented as half-open ranges $[{\inlinecodemath{start}}, {\inlinecodemath{finish}})$, so two intervals are compatible if the next interval's \inlinecode{start} is at least the previous interval's \inlinecode{finish}.

\begin{compactitem}
  \item \inlinecode{select\_intervals(intervals)} returns one maximum-size compatible subset of intervals, in the order they are selected, as original indices into an input vector of \inlinecode{Interval} with fields \inlinecode{start} and \inlinecode{finish}. Every interval must satisfy \inlinecode{start} \textless{} \inlinecode{finish}.
\end{compactitem}

\textbf{Time Complexity}: $O(n \log n)$ per call due to sorting, where $n$ is the number of intervals.

\textbf{Space Complexity}: $O(n)$ auxiliary and $O(n)$ for the returned indices.

\begin{lstlisting}[language=C++]
#include <algorithm>
#include <cassert>
#include <climits>
#include <numeric>
#include <vector>

struct Interval {
  int start, finish;
};

std::vector<int> select_intervals(const std::vector<Interval> &intervals) {
  for (const auto &iv : intervals) {
    assert(iv.start < iv.finish);
  }
  std::vector<int> order(intervals.size());
  std::iota(order.begin(), order.end(), 0);
  std::sort(order.begin(), order.end(), [&](int i, int j) {
    return intervals[i].finish != intervals[j].finish ? intervals[i].finish < intervals[j].finish
                                                      : intervals[i].start < intervals[j].start;
  });
  std::vector<int> selected;
  int last_finish = INT_MIN;
  for (int i : order) {
    const auto &iv = intervals[i];
    if (iv.start >= last_finish) {
      selected.push_back(i);
      last_finish = iv.finish;
    }
  }
  return selected;
}
\end{lstlisting}
\Needspace*{4\baselineskip}
\textbf{Example Usage}\par
\begin{lstlisting}[language=C++]
#include <cassert>
using namespace std;

int main() {
  vector<Interval> intervals{{1, 4}, {3, 5}, {0, 6}, {5, 7}, {8, 9}};
  // Earliest-finish greedy chooses original indices 0, 3, 4.
  assert((select_intervals(intervals) == vector<int>{0, 3, 4}));

  vector<Interval> touching{{0, 2}, {2, 4}, {4, 5}};
  assert((select_intervals(touching) == vector<int>{0, 1, 2}));
  return 0;
}
\end{lstlisting}
\subsection{Weighted Interval Scheduling}
Selects a maximum-weight subset of non-overlapping intervals using dynamic programming and binary search. Unlike unweighted interval scheduling, the earliest finish-time greedy choice is not sufficient when intervals have weights. With intervals sorted by finish time, the best total for the first $i$ intervals either skips interval $i$ or adds its weight to the best total over intervals finishing no later than its start, with that predecessor located by binary search.

Intervals are represented as half-open ranges $[{\inlinecodemath{start}}, {\inlinecodemath{finish}})$, so two intervals are compatible if the next interval's \inlinecode{start} is at least the previous interval's \inlinecode{finish}.

\begin{compactitem}
  \item \inlinecode{select\_weighted\_intervals(intervals)} returns a pair (\inlinecode{weight}, \inlinecode{selected}) containing that maximum weight and the selected intervals as original input indices in execution order, from an input vector of \inlinecode{WeightedInterval} with fields \inlinecode{start}, \inlinecode{finish}, and \inlinecode{weight}. \inlinecode{dp[i]} stores the best answer using the first \inlinecode{i} intervals after sorting by finish time. Each interval must satisfy \inlinecode{start} \textless{} \inlinecode{finish}. The empty subset is allowed, so nonpositive weights need not be selected.
\end{compactitem}

Overflow warning: Accumulated weights must fit in \inlinecode{int64\_t}.

\textbf{Time Complexity}: $O(n \log n)$ per call due to sorting and binary searching compatible intervals.

\textbf{Space Complexity}: $O(n)$ auxiliary and $O(n)$ for the returned indices.

\begin{lstlisting}[language=C++]
#include <algorithm>
#include <cassert>
#include <cstdint>
#include <numeric>
#include <utility>
#include <vector>

struct WeightedInterval {
  int start, finish;
  int64_t weight;
};

std::pair<int64_t, std::vector<int>> select_weighted_intervals(
    const std::vector<WeightedInterval> &intervals
) {
  for (const auto &iv : intervals) {
    assert(iv.start < iv.finish);
  }
  int n = static_cast<int>(intervals.size());
  std::vector<int> order(n), finish(n), prev(n);
  std::iota(order.begin(), order.end(), 0);
  std::sort(order.begin(), order.end(), [&](int i, int j) {
    return intervals[i].finish != intervals[j].finish ? intervals[i].finish < intervals[j].finish
                                                      : intervals[i].start < intervals[j].start;
  });
  for (int i = 0; i < n; i++) {
    finish[i] = intervals[order[i]].finish;
  }
  std::vector<int64_t> dp(n + 1);
  std::vector<char> take(n + 1);
  for (int i = 1; i <= n; i++) {
    int j =
        std::upper_bound(finish.begin(), finish.begin() + i - 1, intervals[order[i - 1]].start) -
        finish.begin();
    prev[i - 1] = j;
    int64_t candidate = dp[j] + intervals[order[i - 1]].weight;  // Overflow warning.
    if (dp[i - 1] < candidate) {
      dp[i] = candidate;
      take[i] = true;
    } else {
      dp[i] = dp[i - 1];
    }
  }
  // Optional: reconstruct one optimal subset of intervals.
  std::vector<int> selected;
  for (int i = n; i > 0;) {
    if (take[i]) {
      selected.push_back(order[i - 1]);
      i = prev[i - 1];
    } else {
      i--;
    }
  }
  std::reverse(selected.begin(), selected.end());
  return {dp[n], selected};
}
\end{lstlisting}
\Needspace*{4\baselineskip}
\textbf{Example Usage}\par
\begin{lstlisting}[language=C++]
#include <cassert>
using namespace std;

int main() {
  vector<WeightedInterval> intervals{{1, 3, 5}, {2, 5, 6}, {4, 6, 5}, {6, 7, 4}, {5, 8, 11}};
  auto [weight, selected] = select_weighted_intervals(intervals);
  // Taking ids 1 and 4 beats the earliest-finish unweighted-looking choices.
  assert(weight == 17 && (selected == vector<int>{1, 4}));

  vector<WeightedInterval> touching{{0, 2, 5}, {2, 4, 6}, {1, 3, 7}};
  auto [touching_weight, touching_selected] = select_weighted_intervals(touching);
  assert(touching_weight == 11 && (touching_selected == vector<int>{0, 1}));
  return 0;
}
\end{lstlisting}
\subsection{Interval Merging and Covering}
Two fundamental interval problems ask either for the union of a collection of intervals or for a minimum-size subcollection that covers a target interval. For the union, sorting by start time makes all intervals that overlap the current merged interval consecutive, so one scan can extend its finish or begin the next disjoint interval.

For minimum covering, maintain the rightmost point covered so far. Among every interval starting at or before that frontier, greedily select one with the farthest finish. Any feasible cover must next use one of those intervals, and replacing that choice with the farthest-reaching one cannot make the remaining target harder to cover. This is different from interval scheduling, which selects the compatible interval with the earliest finish.

Intervals are represented as half-open ranges $[{\inlinecodemath{start}}, {\inlinecodemath{finish}})$. Touching intervals may be merged because their union is another half-open interval with no gap.

\begin{compactitem}
  \item \inlinecode{merge\_intervals(intervals)} returns the union as disjoint intervals sorted by start. Every input interval must satisfy \inlinecode{start} \textless{} \inlinecode{finish}.
  \item \inlinecode{cover\_interval(intervals, lo, hi)} returns a minimum-size list of original input indices whose intervals cover $[{\inlinecodemath{lo}}, {\inlinecodemath{hi}})$, in the order selected, or \inlinecode{std::nullopt} if coverage is impossible. Every input interval must satisfy \inlinecode{start} \textless{} \inlinecode{finish}, and \inlinecode{lo} must not exceed \inlinecode{hi}.
\end{compactitem}

\textbf{Time Complexity}: $O(n \log n)$ per call due to sorting, where $n$ is the number of intervals.

\textbf{Space Complexity}: $O(n)$ auxiliary and $O(n)$ for the returned intervals or indices.

\begin{lstlisting}[language=C++]
#include <algorithm>
#include <cassert>
#include <numeric>
#include <optional>
#include <vector>

struct Interval {
  int start, finish;
};

std::vector<Interval> merge_intervals(std::vector<Interval> intervals) {
  for (const auto &iv : intervals) {
    assert(iv.start < iv.finish);
  }
  std::sort(intervals.begin(), intervals.end(), [](const Interval &a, const Interval &b) {
    return a.start != b.start ? a.start < b.start : a.finish < b.finish;
  });
  std::vector<Interval> merged;
  for (const auto &iv : intervals) {
    if (merged.empty() || merged.back().finish < iv.start) {
      merged.push_back(iv);
    } else {
      merged.back().finish = std::max(merged.back().finish, iv.finish);
    }
  }
  return merged;
}

std::optional<std::vector<int>> cover_interval(
    const std::vector<Interval> &intervals, int lo, int hi
) {
  assert(lo <= hi);
  for (const auto &iv : intervals) {
    assert(iv.start < iv.finish);
  }
  int n = static_cast<int>(intervals.size());
  std::vector<int> order(n);
  std::iota(order.begin(), order.end(), 0);
  std::sort(order.begin(), order.end(), [&](int i, int j) {
    if (intervals[i].start != intervals[j].start) {
      return intervals[i].start < intervals[j].start;
    }
    return intervals[i].finish != intervals[j].finish ? intervals[i].finish > intervals[j].finish
                                                      : i < j;
  });
  std::vector<int> selected;
  int pos = 0, frontier = lo;
  while (frontier < hi) {
    int best = -1, best_finish = frontier;
    while (pos < n && intervals[order[pos]].start <= frontier) {
      int i = order[pos++];
      if (best_finish < intervals[i].finish ||
          (best_finish == intervals[i].finish && best != -1 && i < best)) {
        best = i;
        best_finish = intervals[i].finish;
      }
    }
    if (best == -1) {
      return std::nullopt;
    }
    selected.push_back(best);
    frontier = best_finish;
  }
  return selected;
}
\end{lstlisting}
\Needspace*{4\baselineskip}
\textbf{Example Usage}\par
\begin{lstlisting}[language=C++]
#include <cassert>
#include <vector>
using namespace std;

int main() {
  vector<Interval> intervals{{8, 10}, {1, 3}, {2, 6}, {15, 18}, {10, 12}};
  auto merged = merge_intervals(intervals);
  assert(merged.size() == 3);
  assert(merged[0].start == 1 && merged[0].finish == 6);
  assert(merged[1].start == 8 && merged[1].finish == 12);
  assert(merged[2].start == 15 && merged[2].finish == 18);

  vector<Interval> cover{{0, 2}, {1, 5}, {4, 7}, {6, 10}, {2, 6}, {8, 11}};
  assert(cover_interval(cover, 0, 10) == vector<int>({0, 4, 3}));
  assert(!cover_interval({{0, 2}, {3, 5}}, 0, 5));
  assert(cover_interval({}, 4, 4) == vector<int>());
  return 0;
}
\end{lstlisting}
\subsection{Interval Partitioning}
Assigns intervals to the minimum number of groups so that no two intervals in the same group overlap. This is the classic interval partitioning problem, also known as finding the minimum number of rooms needed for meetings.

The greedy algorithm sorts intervals by start time and reuses the group whose current finish time is earliest whenever possible. Otherwise, it creates a new group. This is optimal because the number of groups must be at least the maximum number of intervals overlapping at any time. When the earliest-finishing active group still overlaps the next interval, every active group overlaps it, so creating a new group is unavoidable. If that group is free, reusing it preserves as many later options as possible.

Intervals are treated as half-open ranges $[{\inlinecodemath{start}}, {\inlinecodemath{finish}})$, so one interval may reuse a room that another interval vacates at the same time.

\begin{compactitem}
  \item \inlinecode{partition\_intervals(intervals)} returns a pair (\inlinecode{rooms}, \inlinecode{room}), where \inlinecode{rooms} is the minimum number of rooms and \inlinecode{room[i]} is the assigned room for input interval \inlinecode{i}. Each \inlinecode{Interval} has fields \inlinecode{start} and \inlinecode{finish}, which must satisfy \inlinecode{start} \textless{} \inlinecode{finish}.
\end{compactitem}

\textbf{Time Complexity}: $O(n \log n)$ per call due to sorting and priority queue operations.

\textbf{Space Complexity}: $O(n)$ auxiliary and $O(n)$ for the returned assignment.

\begin{lstlisting}[language=C++]
#include <algorithm>
#include <cassert>
#include <functional>
#include <numeric>
#include <queue>
#include <utility>
#include <vector>

struct Interval {
  int start, finish;
};

std::pair<int, std::vector<int>> partition_intervals(const std::vector<Interval> &intervals) {
  for (const auto &iv : intervals) {
    assert(iv.start < iv.finish);
  }
  int n = static_cast<int>(intervals.size());
  std::vector<int> order(n);
  std::iota(order.begin(), order.end(), 0);
  std::sort(order.begin(), order.end(), [&](int i, int j) {
    return intervals[i].start != intervals[j].start ? intervals[i].start < intervals[j].start
                                                    : intervals[i].finish < intervals[j].finish;
  });
  std::vector<int> room(n);
  std::priority_queue<std::pair<int, int>, std::vector<std::pair<int, int>>, std::greater<>> pq;
  int rooms = 0;
  for (int i : order) {
    if (!pq.empty() && pq.top().first <= intervals[i].start) {
      int r = pq.top().second;
      pq.pop();
      room[i] = r;
    } else {
      room[i] = rooms++;
    }
    pq.emplace(intervals[i].finish, room[i]);
  }
  return {rooms, room};
}
\end{lstlisting}
\Needspace*{4\baselineskip}
\textbf{Example Usage}\par
\begin{lstlisting}[language=C++]
#include <algorithm>
#include <cassert>
using namespace std;

int main() {
  auto [rooms, room] = partition_intervals(vector<Interval>{{0, 30}, {5, 10}, {15, 20}});
  // The long interval overlaps both short intervals, but the short intervals can share a room.
  assert(rooms == 2);
  assert(room[1] == room[2]);
  assert(room[0] != room[1]);

  vector<Interval> touching{{0, 2}, {2, 4}, {4, 5}};
  auto [touching_rooms, touching_room] = partition_intervals(touching);
  assert(touching_rooms == 1 && *max_element(touching_room.begin(), touching_room.end()) == 0);
  assert(partition_intervals({}).first == 0);
  return 0;
}
\end{lstlisting}
\subsection{Deadline Scheduling}
Schedules unit-time jobs with deadlines and profits to maximize total profit. Each job takes one time slot and must be scheduled at or before its deadline. The greedy algorithm processes jobs in decreasing profit order and places each chosen job in the latest available slot before its deadline.

The latest-slot choice is safe because it uses the current high-profit job while leaving earlier slots available for jobs with tighter deadlines. A disjoint-set forest over time slots makes this efficient: \inlinecode{find(t)} returns the latest still-free slot at or before \inlinecode{t}, and occupying slot \inlinecode{t} links it to the next candidate slot \inlinecode{t - 1}.

\begin{compactitem}
  \item \inlinecode{select\_deadline\_jobs(jobs)} returns a pair (\inlinecode{profit}, \inlinecode{slot}) containing that maximum profit and \inlinecode{slot[i]}, the assigned time slot of input job \inlinecode{i}, or $-1$ if that job is not selected, for a vector of \inlinecode{Job} with fields \inlinecode{deadline} and \inlinecode{profit}. Deadlines are positive integers, and assigned slots are numbered starting from $1$. Jobs with nonpositive profit are not selected.
\end{compactitem}

Overflow warning: The total profit must fit in \inlinecode{int64\_t}.

\textbf{Time Complexity}: $O(n \log n)$ per call due to sorting, with near-constant disjoint-set operations.

\textbf{Space Complexity}: $O(n)$ auxiliary and $O(n)$ for the returned slot assignment.

\begin{lstlisting}[language=C++]
#include <algorithm>
#include <cassert>
#include <cstdint>
#include <numeric>
#include <utility>
#include <vector>

struct Job {
  int deadline;
  int64_t profit;
};

class SlotDSU {
  std::vector<int> root;

 public:
  explicit SlotDSU(int n) : root(n + 1) { std::iota(root.begin(), root.end(), 0); }

  int find(int u) { return root[u] == u ? u : root[u] = find(root[u]); }
  void occupy(int u) { root[u] = find(u - 1); }
};

std::pair<int64_t, std::vector<int>> select_deadline_jobs(const std::vector<Job> &jobs) {
  for (const auto &job : jobs) {
    assert(job.deadline > 0);
  }
  int n = static_cast<int>(jobs.size());
  std::vector<int> order(n);
  std::iota(order.begin(), order.end(), 0);
  std::sort(order.begin(), order.end(), [&](int i, int j) {
    return jobs[i].profit != jobs[j].profit ? jobs[i].profit > jobs[j].profit
                                            : jobs[i].deadline < jobs[j].deadline;
  });
  int max_deadline = 0;
  for (const auto &j : jobs) {
    max_deadline = std::max(max_deadline, j.deadline);
  }
  max_deadline = std::min(max_deadline, n);
  SlotDSU slots(max_deadline);
  std::vector<int> assigned(n, -1);
  int64_t res = 0;
  for (int i : order) {
    if (jobs[i].profit <= 0) {
      continue;
    }
    int slot = slots.find(std::min(jobs[i].deadline, max_deadline));
    if (slot > 0) {
      res += jobs[i].profit;  // Overflow warning.
      assigned[i] = slot;
      slots.occupy(slot);
    }
  }
  return {res, assigned};
}
\end{lstlisting}
\Needspace*{4\baselineskip}
\textbf{Example Usage}\par
\begin{lstlisting}[language=C++]
#include <cassert>
using namespace std;

int main() {
  vector<Job> jobs{{2, 100}, {1, 19}, {2, 27}, {1, 25}, {3, 15}};
  auto [profit, slot] = select_deadline_jobs(jobs);
  // Jobs 0, 2, 4 fill slots 2, 1, 3 for total profit 100 + 27 + 15.
  assert(profit == 142 && (slot == vector<int>{2, -1, 1, -1, 3}));
  return 0;
}
\end{lstlisting}
\subsection{Minimum Late Jobs (Moore-Hodgson)}
Selects a maximum-size subset of jobs that can all be completed by their deadlines. Each job has a processing time and deadline, and all jobs are available at time $0$. The Moore-Hodgson algorithm sorts by deadline and keeps the accepted jobs in a max-heap by processing time; whenever the schedule becomes late, it removes the longest accepted job.

After considering jobs up to some deadline, if the accepted set no longer fits, at least one accepted job must be removed. Removing the longest one leaves the most remaining time while keeping the same number of removed jobs, so it is never worse than removing a shorter accepted job. Repeating this repair after each deadline leaves a largest feasible accepted set.

\begin{compactitem}
  \item \inlinecode{max\_on\_time\_jobs(jobs)} returns the selected jobs as original input indices in a feasible earliest-deadline-first execution order, for an input vector of \inlinecode{TimedJob} with fields \inlinecode{duration} and \inlinecode{deadline}. Durations and deadlines must be nonnegative integers.
\end{compactitem}

Overflow warning: The total duration must fit in \inlinecode{int64\_t}.

\textbf{Time Complexity}: $O(n \log n)$ per call due to sorting and heap operations.

\textbf{Space Complexity}: $O(n)$ auxiliary and $O(n)$ for the returned indices.

\begin{lstlisting}[language=C++]
#include <algorithm>
#include <cassert>
#include <cstdint>
#include <numeric>
#include <queue>
#include <utility>
#include <vector>

struct TimedJob {
  int duration, deadline;
};

std::vector<int> max_on_time_jobs(const std::vector<TimedJob> &jobs) {
  for (const auto &job : jobs) {
    assert(job.duration >= 0 && job.deadline >= 0);
  }
  int n = static_cast<int>(jobs.size());
  std::vector<int> order(n);
  std::iota(order.begin(), order.end(), 0);
  std::sort(order.begin(), order.end(), [&](int i, int j) {
    return jobs[i].deadline != jobs[j].deadline ? jobs[i].deadline < jobs[j].deadline
                                                : jobs[i].duration < jobs[j].duration;
  });
  std::priority_queue<std::pair<int, int>> accepted;
  std::vector<char> selected(n);
  int64_t time = 0;
  for (int i : order) {
    time += jobs[i].duration;  // Overflow warning.
    accepted.emplace(jobs[i].duration, i);
    selected[i] = true;
    if (time > jobs[i].deadline) {
      time -= accepted.top().first;
      selected[accepted.top().second] = false;
      accepted.pop();
    }
  }
  std::vector<int> res;
  for (int i : order) {
    if (selected[i]) {
      res.push_back(i);
    }
  }
  return res;
}
\end{lstlisting}
\Needspace*{4\baselineskip}
\textbf{Example Usage}\par
\begin{lstlisting}[language=C++]
#include <cassert>
using namespace std;

int main() {
  vector<TimedJob> jobs{{3, 4}, {2, 3}, {1, 2}, {2, 7}};
  assert((max_on_time_jobs(jobs) == vector<int>{2, 1, 3}));
  return 0;
}
\end{lstlisting}
\subsection{Weighted Completion Time}
Given jobs that must run one at a time on a single machine, each with a processing time and a weight, order them to minimize the weighted sum of completion times $\sum w_j C_j$, where $C_j$ is the moment job $j$ finishes. Unlike the deadline problems of sections 1.5.6 and 1.5.7, nothing is rejected and no deadline is missed; every order is feasible, and only the total cost differs.

Smith's rule sorts by the ratio of processing time to weight, running the job with the smallest $p_j / w_j$ first. An exchange argument proves it optimal. Consider any schedule that runs job $b$ immediately before job $a$ with $p_a / w_a < p_b / w_b$, and swap the two. Every other job finishes at the same moment, since the pair occupies the same total time, so the change in cost is $p_a w_b - p_b w_a$, which the ratio inequality makes negative. Any schedule violating the ratio order therefore improves under a swap, so an optimal schedule must be sorted by it. With equal weights the rule reduces to shortest processing time first, the classic way to minimize the average completion time.

\begin{compactitem}
  \item \inlinecode{min\_weighted\_completion\_time(jobs)} returns the pair (\inlinecode{cost}, \inlinecode{order}), where \inlinecode{cost} is the minimum weighted sum of completion times and \inlinecode{order} lists the job indices in an optimal sequence. Each job is a pair (\inlinecode{time}, \inlinecode{weight}) with a positive weight and a nonnegative time.
\end{compactitem}

Ties break by index, and the comparison cross-multiplies rather than divides so that it stays exact on integers. Parallel machines are a different problem, NP-hard for general weights.

Overflow warning: the products in the comparator and the running completion time are accumulated in \inlinecode{int64\_t}, so the total processing time multiplied by the total weight must fit in that type.

\textbf{Time Complexity}: $O(n \log n)$ per call, where $n$ is the number of jobs.

\textbf{Space Complexity}: $O(n)$ auxiliary and $O(n)$ for the returned order.

\begin{lstlisting}[language=C++]
#include <algorithm>
#include <cassert>
#include <cstdint>
#include <numeric>
#include <utility>
#include <vector>

std::pair<int64_t, std::vector<int>> min_weighted_completion_time(
    const std::vector<std::pair<int, int>> &jobs
) {
  for (auto [time, weight] : jobs) {
    assert(time >= 0 && weight > 0);
  }
  int n = static_cast<int>(jobs.size());
  std::vector<int> order(n);
  std::iota(order.begin(), order.end(), 0);
  std::sort(order.begin(), order.end(), [&](int a, int b) {
    // Sort by time / weight ascending, cross-multiplied to stay exact.
    int64_t lhs = static_cast<int64_t>(jobs[a].first) * jobs[b].second;
    int64_t rhs = static_cast<int64_t>(jobs[b].first) * jobs[a].second;
    return lhs != rhs ? lhs < rhs : a < b;
  });
  int64_t clock = 0, cost = 0;
  for (int id : order) {
    clock += jobs[id].first;
    cost += clock * jobs[id].second;  // Overflow warning.
  }
  return {cost, order};
}
\end{lstlisting}
\Needspace*{4\baselineskip}
\textbf{Example Usage}\par
\begin{lstlisting}[language=C++]
#include <cassert>
using namespace std;

int main() {
  // (time, weight) pairs. Ratios are 3, 1, and 2, so job 1 runs first, then job 2, then job 0.
  auto [cost, order] = min_weighted_completion_time({{3, 1}, {2, 2}, {4, 2}});
  assert((order == vector<int>{1, 2, 0}));
  assert(cost == 2 * 2 + 6 * 2 + 9 * 1);  // Completion times 2, 6, and 9.

  // Equal weights reduce the rule to shortest processing time first.
  auto [spt_cost, spt_order] = min_weighted_completion_time({{5, 1}, {1, 1}, {3, 1}});
  assert((spt_order == vector<int>{1, 2, 0}));
  assert(spt_cost == 1 + 4 + 9);

  assert(min_weighted_completion_time({}).first == 0);
  return 0;
}
\end{lstlisting}

\section{\texorpdfstring{Streaming and Randomized Algorithms}{1.6 Streaming and Randomized Algorithms}}
\setcounter{section}{6}
\setcounter{subsection}{0}
\subsection{Fisher-Yates Shuffle}
Randomly permutes a range in-place using the Fisher-Yates shuffle. Each possible permutation is produced with equal probability, assuming the random number source is uniform.

The algorithm scans from right to left. At each position $i$, it chooses a random position $j$ from $[0, i]$ and swaps the two elements. This fixes one uniformly random remaining element at a time.

\begin{compactitem}
  \item \inlinecode{fisher\_yates\_shuffle(lo, hi)} randomly shuffles the range $[{\inlinecodemath{lo}}, {\inlinecodemath{hi}})$ using an internal random engine. The range must provide random-access iterators.
\end{compactitem}

\textbf{Time Complexity}: $O(n)$ per call, where $n$ is the distance between \inlinecode{lo} and \inlinecode{hi}.

\textbf{Space Complexity}: $O(1)$ auxiliary.

\begin{lstlisting}[language=C++]
#include <algorithm>
#include <chrono>
#include <random>

template<typename It>
void fisher_yates_shuffle(It lo, It hi) {
  static std::mt19937 rng(std::chrono::steady_clock::now().time_since_epoch().count());
  int n = static_cast<int>(hi - lo);
  for (int i = n - 1; i > 0; i--) {
    std::uniform_int_distribution<int> dist(0, i);
    int j = dist(rng);
    std::iter_swap(lo + i, lo + j);
  }
}
\end{lstlisting}
\Needspace*{4\baselineskip}
\textbf{Example Usage}\par
\begin{lstlisting}[language=C++]
#include <algorithm>
#include <cassert>
#include <vector>
using namespace std;

int main() {
  vector<int> a{1, 2, 3, 4, 5}, original(a);
  fisher_yates_shuffle(a.begin(), a.end());
  sort(a.begin(), a.end());
  assert(a == original);
  return 0;
}
\end{lstlisting}
\subsection{Reservoir Sampling}
Selects uniformly random samples from a stream whose length may be unknown in advance. This is useful when values arrive online or when storing the full dataset is unnecessary.

Reservoir sampling keeps only the chosen sample in memory while processing each element once. The first $k$ arrivals fill the reservoir, and the $i$-th element thereafter is accepted with probability $k/i$, evicting a uniformly chosen occupant. One draw performs both choices: an integer uniform on $[0, i)$ is below $k$ with probability $k/i$, and is then uniform on the $k$ occupants.

The invariant is that after $n \geq k$ arrivals the reservoir is a uniformly random $k$-subset, so each of the $\binom{n}{k}$ possibilities is equally likely. This is certain at $n = k$. Assuming it after $n - 1$, fix a $k$-subset $A$ of the first $n$. If $n \notin A$, the reservoir must already hold $A$ and reject the arrival, with probability $\binom{n-1}{k}^{-1}(n-k)/n$. If $n \in A$, it must hold $A$ with $n$ replaced by one of the $n - k$ outsiders, each then evicted with probability $(k/n)(1/k) = 1/n$, giving $\binom{n-1}{k}^{-1}(n-k)/n$ again. Both equal $\binom{n}{k}^{-1}$, since $\binom{n}{k} = \binom{n-1}{k} \cdot n/(n-k)$. Taking marginals, every element seen is in the sample with probability $k/n$, and \inlinecode{ReservoirSampleOne} is the case $k = 1$.

Each class maintains its reservoir incrementally; call \inlinecode{add(x)} once per stream element in any order, then call \inlinecode{sample()} to retrieve the result.

\begin{compactitem}
  \item \inlinecode{ReservoirSampleOne\textless{}T\textgreater{}()} constructs a single-element sampler.
  \item \inlinecode{ReservoirSampleK\textless{}T\textgreater{}(k)} constructs a \inlinecode{k}-element sampler, allocating the reservoir up front so that no later \inlinecode{add()} reallocates.
  \item \inlinecode{add(x)} incorporates one more stream element.
  \item \inlinecode{sample()} returns the current sample. For \inlinecode{ReservoirSampleOne}, \inlinecode{sample()} requires at least one \inlinecode{add()} call; for \inlinecode{ReservoirSampleK}, it returns the full reservoir, which may contain fewer than \inlinecode{k} elements if the stream was shorter.
  \item \inlinecode{count()} returns the number of stream elements seen so far.
\end{compactitem}

\Needspace*{3\baselineskip}
\textbf{Time Complexity}:
\begin{compactitem}
  \item $O(1)$ per call to \inlinecode{ReservoirSampleOne::add()} and \inlinecode{ReservoirSampleK::add()}.
  \item $O(k)$ per call to the \inlinecode{ReservoirSampleK} constructor.
  \item $O(n)$ to process a range of $n$ elements.
\end{compactitem}

\Needspace*{3\baselineskip}
\textbf{Space Complexity}:
\begin{compactitem}
  \item $O(1)$ storage for \inlinecode{ReservoirSampleOne}.
  \item $O(k)$ storage for \inlinecode{ReservoirSampleK}, allocated on construction.
\end{compactitem}

\begin{lstlisting}[language=C++]
#include <cassert>
#include <chrono>
#include <optional>
#include <random>
#include <vector>

int rand_int(int lo, int hi) {
  static std::mt19937 rng(std::chrono::steady_clock::now().time_since_epoch().count());
  return std::uniform_int_distribution<int>(lo, hi)(rng);
}

template<typename T>
class ReservoirSampleOne {
  std::optional<T> value;
  int seen = 0;

 public:
  void add(const T &x) {
    seen++;
    if (rand_int(1, seen) == 1) {
      value = x;
    }
  }

  const T &sample() const {
    assert(value.has_value());
    return *value;
  }

  int count() const { return seen; }
};

template<typename T>
class ReservoirSampleK {
  int k, seen = 0;
  std::vector<T> reservoir;

 public:
  explicit ReservoirSampleK(int k) : k(k) {
    assert(k >= 0);
    reservoir.reserve(k);
  }

  void add(const T &x) {
    ++seen;
    if (k == 0) {
      return;
    }
    if (static_cast<int>(reservoir.size()) < k) {
      reservoir.push_back(x);
    } else {
      int j = rand_int(0, seen - 1);
      if (j < k) {
        reservoir[j] = x;
      }
    }
  }

  const std::vector<T> &sample() const { return reservoir; }
  int count() const { return seen; }
};
\end{lstlisting}
\Needspace*{4\baselineskip}
\textbf{Example Usage}\par
\begin{lstlisting}[language=C++]
#include <cassert>
using namespace std;

int main() {
  vector<int> a{10, 20, 30, 40, 50};

  // Streaming interface: feed elements one at a time.
  ReservoirSampleOne<int> s1;
  for (int x : a) {
    s1.add(x);
  }
  assert(s1.count() == 5);
  int one = s1.sample();
  assert(one == 10 || one == 20 || one == 30 || one == 40 || one == 50);

  ReservoirSampleK<int> sk(3);
  for (int x : a) {
    sk.add(x);
  }
  assert(sk.count() == 5);
  assert(sk.sample().size() == 3);

  ReservoirSampleK<int> large(10);
  for (int x : a) {
    large.add(x);
  }
  assert(large.count() == 5);
  assert(large.sample().size() == 5);
  return 0;
}
\end{lstlisting}
\subsection{Frequent Elements}
Finds frequent elements in a range using cancellation algorithms that retain only a small set of candidates. Boyer-Moore finds a possible majority element, while Misra-Gries generalizes the same idea to frequencies above any threshold $n/k$.

The Boyer-Moore voting algorithm maintains one candidate and counter, incrementing on a match and decrementing on a mismatch. Each decrement cancels two different values. Removing such pairs cannot eliminate an element occurring more than $\lfloor n/2 \rfloor$ times, so a true majority survives as the final candidate. A second pass verifies that candidate.

Misra-Gries keeps at most $k - 1$ candidates. A tracked value increments its counter, an untracked value claims a free counter, and otherwise every counter is decremented. Each full decrement cancels $k$ distinct occurrences, so any value occurring more than $\lfloor n/k \rfloor$ times must remain a candidate. These candidates require a second pass for exact verification because the algorithm does not retain the stream.

\begin{compactitem}
  \item \inlinecode{majority\_element(lo, hi)} returns an iterator to the first occurrence of the majority element of $[{\inlinecodemath{lo}}, {\inlinecodemath{hi}})$, or \inlinecode{hi} if no majority element exists. The range must provide ForwardIterators, and its value type must support equality.
  \item \inlinecode{frequent\_candidates(lo, hi, k)} returns a hash table of candidate values to their residual counters. \inlinecode{k} must be at least $2$, and the value type must support equality and \inlinecode{std::hash}. Use \inlinecode{k = 2} for the unverified Boyer-Moore candidate phase.
\end{compactitem}

\Needspace*{3\baselineskip}
\textbf{Time Complexity}:
\begin{compactitem}
  \item $O(n)$ per call to \inlinecode{majority\_element(lo, hi)}, where $n$ is the range length.
  \item $O(n)$ expected per call to \inlinecode{frequent\_candidates(lo, hi, k)}: a full-table decrement cancels $k$ occurrences, so all decrement sweeps take $O(n)$ total.
  \item $O(n \cdot k)$ per call to \inlinecode{frequent\_candidates()} in the collision-heavy worst case.
\end{compactitem}

\Needspace*{3\baselineskip}
\textbf{Space Complexity}:
\begin{compactitem}
  \item $O(1)$ auxiliary for \inlinecode{majority\_element()}.
  \item $O(k)$ for the candidates returned by \inlinecode{frequent\_candidates()}.
\end{compactitem}

\begin{lstlisting}[language=C++]
#include <cassert>
#include <iterator>
#include <unordered_map>

template<typename It>
It majority_element(It lo, It hi) {
  int count = 0;
  It candidate = lo;
  for (It it = lo; it != hi; ++it) {
    if (count == 0) {
      candidate = it;
      count = 1;
    } else if (*it == *candidate) {
      count++;
    } else {
      count--;
    }
  }
  // Second pass to verify. If a majority is guaranteed, skip and return candidate directly instead.
  int n = 0;
  count = 0;
  It first = hi;
  for (It it = lo; it != hi; ++it) {
    n++;
    if (*it == *candidate) {
      if (first == hi) {
        first = it;
      }
      count++;
    }
  }
  if (count <= n / 2) {
    return hi;
  }
  return first;
}

template<typename It>
auto frequent_candidates(It lo, It hi, int k) {
  using T = typename std::iterator_traits<It>::value_type;
  assert(k >= 2);
  std::unordered_map<T, int> count;
  for (It it = lo; it != hi; ++it) {
    if (auto found = count.find(*it); found != count.end()) {
      found->second++;
    } else if (static_cast<int>(count.size()) < k - 1) {
      count[*it] = 1;
    } else {
      for (auto jt = count.begin(); jt != count.end();) {
        if (--jt->second == 0) {
          jt = count.erase(jt);
        } else {
          ++jt;
        }
      }
    }
  }
  return count;
}
\end{lstlisting}
\Needspace*{4\baselineskip}
\textbf{Example Usage}\par
\begin{lstlisting}[language=C++]
#include <vector>
using namespace std;

int main() {
  vector<int> a{3, 2, 3, 1, 3};
  assert(majority_element(a.begin(), a.end()) == a.begin());
  vector<int> b{2, 3, 3, 3, 2, 1};
  assert(majority_element(b.begin(), b.end()) == b.end());

  vector<int> c{1, 2, 1, 3, 1, 2, 1, 4, 2, 2, 2};
  auto candidates = frequent_candidates(c.begin(), c.end(), 3);
  assert(candidates.count(1));
  assert(candidates.count(2));
  assert(!candidates.count(3));
  return 0;
}
\end{lstlisting}
\subsection{Online Statistics}
Online stream summaries update useful statistics as values arrive without storing the full stream. This section uses Welford's algorithm to maintain the mean and variance of a stream in constant space, then merges Welford summaries inside a queue to maintain the same statistics over a sliding window. For the median of a stream or of a trailing window, see section 1.6.5.

The naive variance formula $\operatorname{Var}(X) = E[X^2] - E[X]^2$ can lose precision when the two terms are large and close together. Welford instead maintains the current mean $\bar{x}_n$ (stored as \inlinecode{avg}) and the unnormalized second central moment $M_{2,n} = \sum_{i=1}^{n}(x_i - \bar{x}_n)^2$ (stored as \inlinecode{m2}). Each new value shifts the running mean by its deviation divided by the new count, then adds the product of its deviations from the old and new means to \inlinecode{m2}.

\begin{compactitem}
  \item \inlinecode{OnlineStatistics()} constructs an empty summary.
  \item \inlinecode{add(x)} incorporates one more value \inlinecode{x} into the summary.
  \item \inlinecode{count()} returns the number of values seen so far.
  \item \inlinecode{mean()} returns the current arithmetic mean, or $0$ if the summary is empty.
  \item \inlinecode{var\_pop()} returns population variance, dividing by $n$, or $0$ if the summary is empty.
  \item \inlinecode{var\_samp()} returns sample variance, dividing by $n - 1$, or $0$ if fewer than two values have been added.
\end{compactitem}

\textbf{Time Complexity}: $O(1)$ per call to \inlinecode{add()} and each query.

\textbf{Space Complexity}: $O(1)$ auxiliary for all operations.

\begin{lstlisting}[language=C++]
#include <cassert>
#include <utility>
#include <vector>

class OnlineStatistics {
  int n;
  double avg, m2;

 public:
  OnlineStatistics() : n(0), avg(0), m2(0) {}

  void add(double x) {
    n++;
    double delta = x - avg;
    avg += delta / n;
    double delta2 = x - avg;
    m2 += delta * delta2;
  }

  int count() const { return n; }
  double mean() const { return avg; }
  double var_pop() const { return n == 0 ? 0 : m2 / n; }
  double var_samp() const { return n <= 1 ? 0 : m2 / (n - 1); }
};
\end{lstlisting}
To additionally support FIFO removals, Welford's update cannot simply be run backwards to drop the oldest value: reversing it turns \inlinecode{m2} into a difference, which loses the very property that made the forward pass stable, and the resulting error has no way to correct itself as the window advances.

Summaries are mergeable instead. Two summaries of disjoint groups combine by shifting the mean toward the larger group and adding the spread between the two group means, so that $M_2 = M_{2,A} + M_{2,B} + \delta^2 n_A n_B / n$, where $\delta$ is the gap between the group means. That merge is associative, which makes summaries a monoid and puts them within reach of the sliding window aggregation (SWAG) technique of section 1.2.5: a queue of two stacks, each entry storing the fold of everything between it and its end of the queue, reports the fold of the entire queue in $O(1)$ amortized time. The class below is that structure specialized to this summary, so section 1.2.5 is the place to look for the general version over any associative operation, and for why flipping the back stack onto the front stack costs only $O(1)$ amortized per value.

Every value still enters through a forward Welford update of a one-element summary, and the folds only ever add nonnegative terms, so the statistics of a window are as stable as those of a whole stream.

\begin{compactitem}
  \item \inlinecode{StatisticsQueue()} constructs an empty summary.
  \item \inlinecode{add(x)} incorporates \inlinecode{x} as the newest value of the queue.
  \item \inlinecode{remove()} drops the oldest value of the queue, which must be nonempty.
  \item \inlinecode{count()} returns the number of values currently in the queue.
  \item \inlinecode{mean()}, \inlinecode{var\_pop()}, and \inlinecode{var\_samp()} return the same statistics as above, over the values currently in the queue.
\end{compactitem}

The window is left to the caller rather than fixed by the constructor. A trailing window of width $w$ is one call to \inlinecode{add()} followed by a call to \inlinecode{remove()} whenever \inlinecode{count()} exceeds $w$, while a variable-width window from the two-pointer patterns of section 1.2.3 instead removes values until its own condition holds again.

\textbf{Time Complexity}: $O(1)$ amortized per call to \inlinecode{add()} and \inlinecode{remove()}, and $O(1)$ per query.

\textbf{Space Complexity}: $O(n)$ storage, where $n$ is the number of values currently in the queue.

\begin{lstlisting}[language=C++]
class StatisticsQueue {
  struct Summary {
    int n;
    double mean, m2;
  };

  // Chan's parallel merge of two disjoint groups, which never subtracts a value out of m2.
  static Summary merge(const Summary &a, const Summary &b) {
    int n = a.n + b.n;
    double delta = b.mean - a.mean;
    return {n, a.mean + delta * b.n / n, a.m2 + b.m2 + delta * delta * a.n * b.n / n};
  }

  // Each stack entry is (summary of one value, fold), where the fold merges that value with all
  // entries between it and its end of the queue (the front for front_stack, back for back_stack).
  std::vector<std::pair<Summary, Summary>> front_stack, back_stack;

  Summary fold() const {
    if (front_stack.empty()) {
      return back_stack.empty() ? Summary{} : back_stack.back().second;
    }
    if (back_stack.empty()) {
      return front_stack.back().second;
    }
    return merge(front_stack.back().second, back_stack.back().second);
  }

 public:
  void add(double x) {
    Summary value{1, x, 0};
    back_stack.emplace_back(
        value, back_stack.empty() ? value : merge(back_stack.back().second, value)
    );
  }

  void remove() {
    assert(count() > 0);
    if (front_stack.empty()) {
      // Flip the back stack onto the front stack, reversing order so the oldest value ends up on
      // top, and recompute the front folds in queue order.
      while (!back_stack.empty()) {
        Summary oldest = back_stack.back().first;
        back_stack.pop_back();
        front_stack.emplace_back(
            oldest, front_stack.empty() ? oldest : merge(oldest, front_stack.back().second)
        );
      }
    }
    front_stack.pop_back();
  }

  int count() const { return static_cast<int>(front_stack.size() + back_stack.size()); }
  double mean() const { return fold().mean; }

  double var_pop() const {
    Summary s = fold();
    return s.n == 0 ? 0 : s.m2 / s.n;
  }

  double var_samp() const {
    Summary s = fold();
    return s.n <= 1 ? 0 : s.m2 / (s.n - 1);
  }
};
\end{lstlisting}
\Needspace*{4\baselineskip}
\textbf{Example Usage}\par
\begin{lstlisting}[language=C++]
#include <cassert>
#include <cmath>
#include <vector>
using namespace std;

bool EQ(double a, double b) {
  return fabs(a - b) < 1e-9;
}

int main() {
  OnlineStatistics empty;
  assert(empty.count() == 0 && EQ(empty.mean(), 0.0));
  assert(EQ(empty.var_pop(), 0.0) && EQ(empty.var_samp(), 0.0));

  OnlineStatistics singleton;
  singleton.add(7);
  assert(singleton.count() == 1 && EQ(singleton.mean(), 7.0));
  assert(EQ(singleton.var_pop(), 0.0) && EQ(singleton.var_samp(), 0.0));

  OnlineStatistics stats;
  for (int x = 1; x <= 5; x++) {
    stats.add(x);
  }
  assert(stats.count() == 5 && EQ(stats.mean(), 3.0));
  assert(EQ(stats.var_pop(), 2.0) && EQ(stats.var_samp(), 2.5));

  OnlineStatistics shifted;
  shifted.add(1e12 + 1);
  shifted.add(1e12 + 2);
  shifted.add(1e12 + 3);
  assert(EQ(shifted.mean(), 1e12 + 2));
  assert(EQ(shifted.var_pop(), 2.0 / 3) && EQ(shifted.var_samp(), 1.0));

  StatisticsQueue empty_queue;
  assert(empty_queue.count() == 0 && EQ(empty_queue.mean(), 0.0));
  assert(EQ(empty_queue.var_pop(), 0.0) && EQ(empty_queue.var_samp(), 0.0));

  // A trailing window of width 3, maintained by the caller.
  StatisticsQueue window;
  for (int x = 1; x <= 5; x++) {
    window.add(x);
    if (window.count() > 3) {
      window.remove();
    }
  }
  assert(window.count() == 3);  // The window holds 3, 4, and 5.
  assert(EQ(window.mean(), 4.0));
  assert(EQ(window.var_pop(), 2.0 / 3) && EQ(window.var_samp(), 1.0));

  // The same shifted values as above, now with an extra value removed from the front first.
  StatisticsQueue shifted_window;
  shifted_window.add(1e12);
  shifted_window.add(1e12 + 1);
  shifted_window.add(1e12 + 2);
  shifted_window.add(1e12 + 3);
  shifted_window.remove();
  assert(EQ(shifted_window.mean(), 1e12 + 2));
  assert(EQ(shifted_window.var_pop(), 2.0 / 3) && EQ(shifted_window.var_samp(), 1.0));

  // A variable-width window, which a fixed size cannot express: extend while the population
  // variance stays at most 1, shrinking from the front whenever it does not.
  StatisticsQueue spread;
  int longest = 0;
  for (double x : {1.0, 1.0, 2.0, 8.0, 8.0, 9.0}) {
    spread.add(x);
    while (spread.var_pop() > 1) {
      spread.remove();
    }
    if (spread.count() > longest) {
      longest = spread.count();
    }
  }
  assert(longest == 3);  // The values 1, 1, and 2 have population variance 2/9.

  return 0;
}
\end{lstlisting}
\subsection{Online Median}
The median of a stream can be maintained as values arrive, without ever sorting the stream. This section keeps the median of a growing stream in two heaps, and the median of a changing collection in a multiset, the difference between them being whether values ever have to leave again. For the mean and variance of a stream, see section 1.6.4.

Two heaps hold the values in halves, using linear storage. A max-heap stores the lower half and a min-heap stores the upper half, balanced so that the lower heap has either the same number of values as the upper heap or one extra value. Their roots are therefore the middle value or values of the stream.

\begin{compactitem}
  \item \inlinecode{OnlineMedian\textless{}T\textgreater{}()} constructs an empty median tracker for an ordered numeric type \inlinecode{T}.
  \item \inlinecode{empty()} returns whether the tracker has no values.
  \item \inlinecode{count()} returns the number of values in the tracker.
  \item \inlinecode{add(x)} inserts \inlinecode{x} into the tracker.
  \item \inlinecode{lower\_median()} and \inlinecode{upper\_median()} return the lower and the upper of the two middle values, or the middle value itself for an odd count. The tracker must be nonempty.
  \item \inlinecode{median()} returns the middle value for an odd count or the arithmetic mean of the two middle values for an even count. The tracker must be nonempty.
\end{compactitem}

Since \inlinecode{median()} averages in \inlinecode{double}, an even count over a type whose values exceed the exactly representable range of \inlinecode{double} should be averaged by the caller from \inlinecode{lower\_median()} and \inlinecode{upper\_median()} instead.

\textbf{Time Complexity}: $O(\log n)$ per call to \inlinecode{add()} and $O(1)$ per query, where $n$ is the number of values.

\textbf{Space Complexity}: $O(n)$ storage.

\begin{lstlisting}[language=C++]
#include <cassert>
#include <functional>
#include <iterator>
#include <queue>
#include <set>
#include <vector>

template<typename T>
class OnlineMedian {
  std::priority_queue<T> lower;
  std::priority_queue<T, std::vector<T>, std::greater<T>> upper;

 public:
  bool empty() const { return lower.empty(); }
  int count() const { return static_cast<int>(lower.size() + upper.size()); }

  void add(const T &x) {
    if (lower.empty() || x <= lower.top()) {
      lower.push(x);
    } else {
      upper.push(x);
    }
    if (lower.size() < upper.size()) {
      lower.push(upper.top());
      upper.pop();
    } else if (lower.size() > upper.size() + 1) {
      upper.push(lower.top());
      lower.pop();
    }
  }

  const T &lower_median() const {
    assert(!empty());
    return lower.top();
  }

  const T &upper_median() const {
    assert(!empty());
    return lower.size() != upper.size() ? lower.top() : upper.top();
  }

  double median() const {
    return static_cast<double>(lower_median()) / 2 + static_cast<double>(upper_median()) / 2;
  }
};
\end{lstlisting}
To additionally support arbitrary removals, a multiset (balanced binary search tree) holds the values in sorted order while an iterator tracks the lower median. Each insertion or erasure shifts the middle of the multiset by at most one position, so the iterator only ever steps one element in either direction: whether it moves depends on which side of the median the value lies and on the parity of the resulting size. Unlike the two heaps above, and unlike the statistics queue of section 1.6.4 that can only drop its oldest value, a multiset erases any value it holds, so removals may come in any order.

\begin{compactitem}
  \item \inlinecode{DeletableMedian\textless{}T\textgreater{}()} constructs an empty median tracker for an ordered numeric type \inlinecode{T}.
  \item \inlinecode{add(x)} inserts \inlinecode{x} into the median tracker.
  \item \inlinecode{remove(x)} erases one copy of \inlinecode{x}, which must currently be present.
  \item \inlinecode{count()} and \inlinecode{empty()} return the number of values and whether the tracker is empty.
  \item \inlinecode{lower\_median()}, \inlinecode{upper\_median()}, and \inlinecode{median()} answer as above, over the values currently held. The tracker must be nonempty.
\end{compactitem}

A trailing window of width $w$ over an array is one call to \inlinecode{add()} followed by a call to \inlinecode{remove()} on the value leaving the window, and a two-pointer scan from section 1.2.3 removes by its own condition instead. Because removals are unordered, the same structure also serves the add and remove operations of Mo's algorithm in section 2.5.4, which answers median queries over arbitrary ranges offline. Replacing the parked iterator with an order-statistic tree, such as the \inlinecode{\_\_gnu\_pbds} tree of section 8.6, generalizes further to any rank rather than just the middle.

\textbf{Time Complexity}: $O(\log n)$ per call to \inlinecode{add()} and \inlinecode{remove()}, and $O(1)$ per query, where $n$ is the number of values currently held.

\textbf{Space Complexity}: $O(n)$ storage.

\begin{lstlisting}[language=C++]
template<typename T>
class DeletableMedian {
  std::multiset<T> sorted;
  typename std::multiset<T>::iterator mid;

 public:
  bool empty() const { return sorted.empty(); }
  int count() const { return static_cast<int>(sorted.size()); }

  void add(const T &x) {
    if (sorted.empty()) {
      mid = sorted.insert(x);
      return;
    }
    bool below = x < *mid;
    sorted.insert(x);
    if (below) {
      if (sorted.size() % 2 == 0) {
        --mid;
      }
    } else if (sorted.size() % 2 == 1) {
      ++mid;
    }
  }

  void remove(const T &x) {
    assert(!empty());
    if (sorted.size() == 1) {
      assert(x == *mid);
      sorted.clear();  // The iterator is meaningless while empty, and add() resets it.
      return;
    }
    // Erasing the median's own node requires stepping off it first. Any other node holding x
    // lies strictly on one side of the median, so erasing it can never invalidate the iterator.
    auto target = (x == *mid) ? mid : sorted.find(x);
    assert(target != sorted.end());
    if (x <= *mid && sorted.size() % 2 == 0) {
      ++mid;
    } else if (x >= *mid && sorted.size() % 2 == 1) {
      --mid;
    }
    sorted.erase(target);
  }

  const T &lower_median() const {
    assert(!empty());
    return *mid;
  }

  const T &upper_median() const {
    assert(!empty());
    return sorted.size() % 2 == 1 ? *mid : *std::next(mid);
  }

  double median() const {
    return static_cast<double>(lower_median()) / 2 + static_cast<double>(upper_median()) / 2;
  }
};
\end{lstlisting}
\Needspace*{4\baselineskip}
\textbf{Example Usage}\par
\begin{lstlisting}[language=C++]
#include <cmath>
using namespace std;

bool EQ(double a, double b) {
  return fabs(a - b) < 1e-9;
}

int main() {
  OnlineMedian<int> medians;
  assert(medians.empty() && medians.count() == 0);
  medians.add(5);
  assert(EQ(medians.median(), 5));
  medians.add(1);
  assert(medians.lower_median() == 1 && medians.upper_median() == 5);
  assert(EQ(medians.median(), 3));
  medians.add(9);
  assert(medians.lower_median() == 5 && medians.upper_median() == 5);
  medians.add(3);
  assert(medians.lower_median() == 3 && medians.upper_median() == 5);
  assert(EQ(medians.median(), 4));

  // A trailing window of width 3 over an array, maintained by the caller.
  vector<int> a{5, 1, 9, 3, 3};
  DeletableMedian<int> window;
  assert(window.empty() && window.count() == 0);
  for (int hi = 0, lo = 0; hi < static_cast<int>(a.size()); hi++) {
    window.add(a[hi]);
    if (window.count() > 3) {
      window.remove(a[lo++]);
    }
  }
  assert(window.count() == 3);  // The window holds 9, 3, and 3.
  assert(window.lower_median() == 3 && EQ(window.median(), 3));

  // Removals need not follow arrival order, and may target the median's own value.
  DeletableMedian<int> tracker;
  for (int x : {1, 2, 3, 4, 5}) {
    tracker.add(x);
  }
  assert(EQ(tracker.median(), 3));
  tracker.remove(3);  // Erasing the median itself steps the iterator off its own node first.
  assert(tracker.lower_median() == 2 && tracker.upper_median() == 4);
  tracker.remove(5);
  assert(EQ(tracker.median(), 2));

  // Duplicates of the median are erased one copy at a time.
  DeletableMedian<int> repeats;
  for (int x : {7, 7, 7}) {
    repeats.add(x);
  }
  repeats.remove(7);
  assert(repeats.count() == 2 && EQ(repeats.median(), 7));
  repeats.remove(7);
  assert(repeats.count() == 1 && repeats.lower_median() == 7);
  return 0;
}
\end{lstlisting}
\subsection{Probabilistic Sketches}
A sketch answers questions about a stream using far less memory than the stream itself, in exchange for a bounded, one-sided error. This section implements a Bloom filter, which tests set membership, and a count-min sketch, which estimates how often each value occurred. Both spread every value over several independent positions of a fixed-size table, and both derive those positions from one pair of base hashes: the Kirsch-Mitzenmacher technique uses $h_i(x) = h_1(x) + i \cdot h_2(x)$ instead of computing an unrelated hash for every position, which matches the false-positive behavior of truly independent hashes while paying for only two. Both base hashes come from the SplitMix64 mixer of section 3.2.1, and the second is forced odd so that repeated steps never stand still.

A Bloom filter represents a set as a bit array. Inserting a value sets the $k$ bits it hashes to, and a membership test reports that a value is present only if all $k$ of its bits are set. Bits are never cleared, so a value that was inserted always tests positive: the error is one-sided, and only a false positive is possible, produced when unrelated insertions happen to have set all $k$ bits. For $n$ values in $m$ bits, the false-positive probability is minimized by $k = (m/n) \ln 2$, which yields $m = -n \ln p / (\ln 2)^2$ bits for a target rate $p$.

\begin{compactitem}
  \item \inlinecode{BloomFilter\textless{}T\textgreater{}(capacity, false\_positive\_rate)} constructs an empty filter sized for \inlinecode{capacity} insertions at the given target rate, where \inlinecode{capacity} is positive and the rate lies in $(0, 1)$. \inlinecode{BloomFilter\textless{}T, Hash\textgreater{}(capacity, false\_positive\_rate, hash)} instead stores the supplied hasher.
  \item \inlinecode{insert(x)} adds \inlinecode{x} to the filter.
  \item \inlinecode{contains(x)} returns whether \inlinecode{x} may be in the filter. It never returns \inlinecode{false} for a value that was inserted, but may return \inlinecode{true} for one that was not.
  \item \inlinecode{count()} returns the number of insertions performed, counting repeated values separately.
  \item \inlinecode{bit\_count()} and \inlinecode{hash\_count()} return the derived table size $m$ and number of positions $k$.
  \item \inlinecode{false\_positive\_rate()} returns the current estimated rate $(1 - e^{-kn/m})^k$, which rises as the filter fills past its planned capacity.
\end{compactitem}

Deletion is not supported, since clearing a bit could hide an unrelated value that also set it. Use a counting Bloom filter, which replaces each bit with a small counter, when values must leave.

\Needspace*{3\baselineskip}
\textbf{Time Complexity}:
\begin{compactitem}
  \item $O(1)$ per call to \inlinecode{count()}, \inlinecode{bit\_count()}, \inlinecode{hash\_count()}, and \inlinecode{false\_positive\_rate()}.
  \item $O(k)$ per call to \inlinecode{insert()} and \inlinecode{contains()}, where $k$ is \inlinecode{hash\_count()}, plus the cost of one call to the hasher.
  \item $O(m)$ per call to the constructor, where $m$ is \inlinecode{bit\_count()}.
\end{compactitem}

\Needspace*{3\baselineskip}
\textbf{Space Complexity}:
\begin{compactitem}
  \item $O(m)$ for storage of the filter, where $m$ is \inlinecode{bit\_count()}.
  \item $O(1)$ auxiliary for all operations.
\end{compactitem}

\begin{lstlisting}[language=C++]
#include <algorithm>
#include <cassert>
#include <cmath>
#include <cstdint>
#include <functional>
#include <utility>
#include <vector>

uint64_t splitmix64(uint64_t x) {
  x += 0x9e3779b97f4a7c15ULL;
  x = (x ^ (x >> 30)) * 0xbf58476d1ce4e5b9ULL;
  x = (x ^ (x >> 27)) * 0x94d049bb133111ebULL;
  return x ^ (x >> 31);
}

template<typename T, typename Hash>
std::pair<uint64_t, uint64_t> hash_pair(const Hash &hash, const T &x) {
  uint64_t h1 = splitmix64(static_cast<uint64_t>(hash(x)));
  return {h1, splitmix64(h1 + 1) | 1};
}

template<typename T, typename Hash = std::hash<T>>
class BloomFilter {
  std::vector<uint64_t> words;
  uint64_t nbits;
  int hashes, inserted;
  Hash hash;

 public:
  BloomFilter(int capacity, double false_positive_rate, Hash hash = Hash{})
      : nbits(0), hashes(0), inserted(0), hash(hash) {
    assert(capacity > 0 && false_positive_rate > 0 && false_positive_rate < 1);
    double ln2 = std::log(2.0);
    nbits =
        static_cast<uint64_t>(std::ceil(-capacity * std::log(false_positive_rate) / (ln2 * ln2)));
    hashes = std::max(1, static_cast<int>(std::lround(nbits * ln2 / capacity)));
    words.resize((nbits + 63) / 64);
  }

  void insert(const T &x) {
    auto [h1, h2] = hash_pair(hash, x);
    for (int i = 0; i < hashes; i++) {
      uint64_t bit = (h1 + static_cast<uint64_t>(i) * h2) % nbits;
      words[bit >> 6] |= 1ULL << (bit & 63);
    }
    inserted++;
  }

  bool contains(const T &x) const {
    auto [h1, h2] = hash_pair(hash, x);
    for (int i = 0; i < hashes; i++) {
      uint64_t bit = (h1 + static_cast<uint64_t>(i) * h2) % nbits;
      if (((words[bit >> 6] >> (bit & 63)) & 1) == 0) {
        return false;
      }
    }
    return true;
  }

  int count() const { return inserted; }
  uint64_t bit_count() const { return nbits; }
  int hash_count() const { return hashes; }

  double false_positive_rate() const {
    double filled = 1 - std::exp(-1.0 * hashes * inserted / nbits);
    return std::pow(filled, hashes);
  }
};
\end{lstlisting}
A count-min sketch estimates the total weight added under each value. It stores an $R$ by $C$ table of counters and, for each row, adds the weight to the one counter that the value hashes to in that row. Every counter that a value touches therefore holds that value's true total plus the totals of whatever else collided there, so with nonnegative weights each row is an overestimate and the minimum over the rows is the sharpest one available. Choosing $C = \lceil e/\epsilon \rceil$ and $R = \lceil \ln(1/\delta) \rceil$ bounds the overestimate by $\epsilon N$ with probability at least $1 - \delta$, where $N$ is the total weight added.

\begin{compactitem}
  \item \inlinecode{CountMinSketch\textless{}T\textgreater{}(eps, delta)} constructs an empty sketch whose estimates exceed the true count by at most \inlinecode{eps} times the total weight, with probability at least $1 - {\inlinecodemath{delta}}$. Both parameters must lie in $(0, 1)$. \inlinecode{CountMinSketch\textless{}T, Hash\textgreater{}(eps, delta, hash)} instead stores the supplied hasher.
  \item \inlinecode{add(x, c = 1)} adds weight \inlinecode{c} to the count of \inlinecode{x}.
  \item \inlinecode{count(x)} returns the estimated total weight of \inlinecode{x}, which is never an underestimate.
  \item \inlinecode{total()} returns the total weight added over all values.
  \item \inlinecode{num\_rows()} and \inlinecode{num\_cols()} return the derived table dimensions.
\end{compactitem}

The error bound is relative to the total weight of the whole stream, so a sketch estimates heavy values well and light ones poorly. That makes it a natural companion to the deterministic heavy-hitter candidates of section 1.6.3: run both, then use the sketch to rank the candidates. Negative weights break the guarantee, since a row is then no longer an overestimate; deletions require the count-mean-min variant or a conservative sketch.

\Needspace*{3\baselineskip}
\textbf{Time Complexity}:
\begin{compactitem}
  \item $O(R)$ per call to \inlinecode{add()} and \inlinecode{count()}, plus the cost of one call to the hasher.
  \item $O(R \cdot C)$ per call to the constructor.
  \item $O(1)$ per call to \inlinecode{total()}, \inlinecode{num\_rows()}, and \inlinecode{num\_cols()}.
\end{compactitem}

\Needspace*{3\baselineskip}
\textbf{Space Complexity}:
\begin{compactitem}
  \item $O(w \cdot d)$ for storage of the sketch.
  \item $O(1)$ auxiliary for all operations.
\end{compactitem}

\begin{lstlisting}[language=C++]
template<typename T, typename Hash = std::hash<T>>
class CountMinSketch {
  std::vector<std::vector<int64_t>> table;
  int rows, cols;
  int64_t weight;
  Hash hash;

 public:
  CountMinSketch(double eps, double delta, Hash hash = Hash{}) : weight(0), hash(hash) {
    assert(eps > 0 && eps < 1 && delta > 0 && delta < 1);
    rows = std::max(1, static_cast<int>(std::ceil(std::log(1 / delta))));
    cols = static_cast<int>(std::ceil(std::exp(1.0) / eps));
    table.assign(rows, std::vector<int64_t>(cols));
  }

  void add(const T &x, int64_t c = 1) {
    auto [h1, h2] = hash_pair(hash, x);
    for (int r = 0; r < rows; r++) {
      table[r][(h1 + static_cast<uint64_t>(r) * h2) % cols] += c;
    }
    weight += c;  // Overflow warning.
  }

  int64_t count(const T &x) const {
    auto [h1, h2] = hash_pair(hash, x);
    int64_t res = INT64_MAX;
    for (int r = 0; r < rows; r++) {
      res = std::min(res, table[r][(h1 + static_cast<uint64_t>(r) * h2) % cols]);
    }
    return res;
  }

  int64_t total() const { return weight; }
  int num_rows() const { return rows; }
  int num_cols() const { return cols; }
};
\end{lstlisting}
\Needspace*{4\baselineskip}
\textbf{Example Usage}\par
\begin{lstlisting}[language=C++]
#include <string>
using namespace std;

int main() {
  BloomFilter<string> filter(1000, 0.01);
  assert(filter.bit_count() == 9586 && filter.hash_count() == 7);
  for (int i = 0; i < 1000; i++) {
    filter.insert("key" + to_string(i));
  }
  assert(filter.count() == 1000);
  for (int i = 0; i < 1000; i++) {
    assert(filter.contains("key" + to_string(i)));  // Inserted values never test negative.
  }
  int false_positives = 0;
  for (int i = 1000; i < 2000; i++) {
    false_positives += filter.contains("key" + to_string(i));
  }
  assert(false_positives < 40);  // Absent values may test positive, about 10 of these 1000.
  assert(filter.false_positive_rate() < 0.02);

  CountMinSketch<string> sketch(0.01, 0.01);
  sketch.add("apple", 3);
  sketch.add("banana");
  assert(sketch.total() == 4);
  assert(sketch.count("apple") == 3 && sketch.count("banana") == 1);
  // A value that was never added reads 0 unless it collides in every row, and an estimate is never
  // lower than the truth.
  assert(sketch.count("cherry") <= sketch.total());
  return 0;
}
\end{lstlisting}

\section{\texorpdfstring{Bit Manipulation}{1.7 Bit Manipulation}}
\setcounter{section}{7}
\setcounter{subsection}{0}
\subsection{Bit Operations and Masks}
Integers are all fixed-size sets of bits, where the $i$-th bit (counting from the least significant bit where $i = 0$) is ``present'' when set to $1$. Treating an \inlinecode{int} as a bitmask makes set membership, insertion, deletion, and iteration into single machine instructions, which is why bitmasks are the backbone of subset dynamic programming and many low-level tricks. The helpers below operate on \inlinecode{Mask}, which is \inlinecode{uint32\_t} by default; change its alias to \inlinecode{uint64\_t} to use 64-bit masks. The loop implementations below exist for education only; production code should prefer the matching GCC built-ins when available, or on C++20 systems the generic \inlinecode{constexpr} equivalents from the \inlinecode{\textless{}bit\textgreater{}} header, which are additionally well-defined at $0$.

\begin{compactitem}
  \item \inlinecode{test\_bit(x, i)} returns whether the \inlinecode{i}-th bit of \inlinecode{x} is set, where $0 \leq {\inlinecodemath{i}} < {\inlinecodemath{MASK\_BITS}}$.
  \item \inlinecode{set\_bit(x, i)}, \inlinecode{clear\_bit(x, i)}, and \inlinecode{toggle\_bit(x, i)} return \inlinecode{x} with the \inlinecode{i}-th respectively forced to $1$, forced to $0$, or flipped, where $0 \leq {\inlinecodemath{i}} < {\inlinecodemath{MASK\_BITS}}$.
  \item \inlinecode{lowest\_set\_bit(x)} returns the value of the lowest set bit of \inlinecode{x} (a power of two), or $0$ if \inlinecode{x} is $0$. \inlinecode{clear\_lowest\_set\_bit(x)} returns \inlinecode{x} with its lowest set bit removed.
  \item \inlinecode{popcount(x)} returns the number of set bits, analogous to C++20's \inlinecode{std::popcount()}.
  \item \inlinecode{parity(x)} returns $1$ if the number of set bits is odd and $0$ otherwise.
  \item \inlinecode{ctz(x)} returns the number of trailing 0-bits for \inlinecode{x \textgreater{} 0}, analogous to C++20's \inlinecode{std::countr\_zero()}.
  \item \inlinecode{ffs(x)} returns one plus the position of the lowest 1-bit of \inlinecode{x}, or $0$ if \inlinecode{x} is $0$.
  \item \inlinecode{clz(x)} returns the number of leading 0-bits for \inlinecode{x \textgreater{} 0}, analogous to C++20's \inlinecode{std::countl\_zero()}.
  \item \inlinecode{is\_pow2(x)} returns whether \inlinecode{x} has exactly one set bit, analogous to C++20's \inlinecode{std::has\_single\_bit()}.
  \item \inlinecode{floor\_pow2(x)} returns the largest power of two that is $\leq$ \inlinecode{x} (for \inlinecode{x} $> 0$), analogous to C++20's \inlinecode{std::bit\_floor()}.
  \item \inlinecode{ceil\_pow2(x)} returns the smallest power of two that is $\geq$ \inlinecode{x}, for $0 < {\inlinecodemath{x}} \leq 2^{b - 1}$ where $b$ is \inlinecode{MASK\_BITS}, analogous to C++20's \inlinecode{std::bit\_ceil()}.
  \item \inlinecode{for\_each\_set\_bit(x, f)} calls \inlinecode{f(i)} once for each set bit position \inlinecode{i} of \inlinecode{x}, in increasing order.
\end{compactitem}

\Needspace*{3\baselineskip}
\textbf{Time Complexity}:
\begin{compactitem}
  \item $O(b)$ worst case per call to \inlinecode{popcount()}, \inlinecode{parity()}, \inlinecode{ctz()}, \inlinecode{ffs()}, \inlinecode{clz()}, and \inlinecode{for\_each\_set\_bit()}, where $b$ is \inlinecode{MASK\_BITS}.
  \item $O(1)$ per call to every other function.
\end{compactitem}

\textbf{Space Complexity}: $O(1)$ auxiliary for all operations.

\begin{lstlisting}[language=C++]
#include <cassert>
#include <climits>
#include <cstdint>

using Mask = uint32_t;
const int MASK_BITS = sizeof(Mask) * CHAR_BIT;

bool test_bit(Mask x, int i) { return (x >> i) & 1; }
Mask set_bit(Mask x, int i) { return x | (Mask{1} << i); }
Mask clear_bit(Mask x, int i) { return x & ~(Mask{1} << i); }
Mask toggle_bit(Mask x, int i) { return x ^ (Mask{1} << i); }
Mask lowest_set_bit(Mask x) { return x & -x; }
Mask clear_lowest_set_bit(Mask x) { return x & (x - 1); }

// popcount(), parity(), ctz(), ffs(), and clz() are spelled out here for educational purposes.

int popcount(Mask x) {  // std::popcount(x) in C++20.
  int count = 0;
  for (; x != 0; x = clear_lowest_set_bit(x)) {
    count++;
  }
  return count;
}

int parity(Mask x) {
  return popcount(x) & 1;
}

int ctz(Mask x) {  // std::countr_zero(x) in C++20.
  assert(x != 0);
  int count = 0;
  for (; (x & 1) == 0; x >>= 1) {
    count++;
  }
  return count;
}

int ffs(Mask x) {
  return x == 0 ? 0 : ctz(x) + 1;
}

int clz(Mask x) {  // std::countl_zero(x) in C++20.
  assert(x != 0);
  int count = 0;
  for (Mask mask = Mask{1} << (MASK_BITS - 1); (x & mask) == 0; mask >>= 1) {
    count++;
  }
  return count;
}

bool is_pow2(Mask x) {  // std::has_single_bit(x) in C++20.
  return x != 0 && (x & (x - 1)) == 0;
}

Mask floor_pow2(Mask x) {  // std::bit_floor(x) in C++20.
  assert(x != 0);
  return Mask{1} << (MASK_BITS - 1 - clz(x));
}

Mask ceil_pow2(Mask x) {  // std::bit_ceil(x) in C++20.
  assert(0 < x && x <= (Mask{1} << (MASK_BITS - 1)));
  return is_pow2(x) ? x : Mask{1} << (MASK_BITS - clz(x));
}

template<typename Fn>
void for_each_set_bit(Mask x, Fn f) {
  while (x != 0) {
    f(ctz(x));
    x = clear_lowest_set_bit(x);
  }
}
\end{lstlisting}
\Needspace*{4\baselineskip}
\textbf{Example Usage}\par
\begin{lstlisting}[language=C++]
#include <cassert>
#include <vector>
using namespace std;

int main() {
  Mask x = 0b101100;
  assert(test_bit(x, 2) == true);
  assert(test_bit(x, 0) == false);
  assert(set_bit(x, 0) == 0b101101U);
  assert(clear_bit(x, 2) == 0b101000U);
  assert(toggle_bit(x, 3) == 0b100100U);

  assert(lowest_set_bit(x) == 0b100U);
  assert(clear_lowest_set_bit(x) == 0b101000U);
  assert(lowest_set_bit(0) == 0U);
  assert(clear_lowest_set_bit(0) == 0U);
  assert(popcount(x) == 3);
  assert(popcount(0) == 0);
  assert(parity(x) == 1);
  assert(parity(0b101101U) == 0);
  assert(parity(0) == 0);
  assert(ctz(x) == 2);
  assert(ffs(x) == 3);
  assert(ffs(0) == 0);
  assert(clz(x) == MASK_BITS - 6);

  assert(is_pow2(16) == true);
  assert(is_pow2(24) == false);
  assert(floor_pow2(20) == 16U);
  assert(ceil_pow2(20) == 32U);
  assert(ceil_pow2(16) == 16U);
  assert(ceil_pow2((Mask{1} << (MASK_BITS - 1)) - 1) == (Mask{1} << (MASK_BITS - 1)));

  vector<int> bits;
  for_each_set_bit(x, [&](int b) { bits.push_back(b); });
  assert((bits == vector<int>{2, 3, 5}));
  return 0;
}
\end{lstlisting}
\subsection{Subset and Submask Iteration}
Many bitmask algorithms must visit related masks in a structured order: every submask of a fixed mask, every mask with a fixed number of set bits, or every mask while accumulating contributions from its submasks. Each pattern has a short, well-known idiom.

The classic submask loop \inlinecode{for (int s = m; s \textgreater{} 0; s = (s - 1) \& m)} walks every nonzero submask of \inlinecode{m} in decreasing order; subtracting 1 borrows through the zero bits of \inlinecode{m}, and the \inlinecode{\& m} masks them back off. Summed over all masks \inlinecode{m} of \inlinecode{n} bits, the total work is $\sum_k \binom{n}{k} 2^k = 3^n$, since each of the \inlinecode{n} bit positions is independently absent, present in \inlinecode{m} only, or present in both \inlinecode{m} and \inlinecode{s}.

Gosper's hack advances to the next mask with the same popcount. It isolates the lowest set bit and adds it to move the lowest movable 1 one position left, clearing the block of 1-bits beneath it. The remaining expression counts those cleared bits and packs them into the least-significant positions, producing the smallest larger integer with the same number of set bits.

\begin{compactitem}
  \item \inlinecode{submasks(m)} returns all submasks of \inlinecode{m}, including \inlinecode{m} itself and 0, in decreasing order.
  \item \inlinecode{next\_same\_popcount(x)} returns the smallest integer greater than \inlinecode{x} with the same number of set bits (Gosper's hack), for \inlinecode{x \textgreater{} 0} when that next mask is representable in \inlinecode{Mask}.
  \item \inlinecode{masks\_with\_popcount(n, k)} returns all \inlinecode{k}-bit subsets of an \inlinecode{n}-bit universe as masks, in increasing order, where $0 \leq k \leq n < {\inlinecodemath{MASK\_BITS}}$.
  \item \inlinecode{subset\_sum\_transform(f)} overwrites \inlinecode{f} (indexed by mask over \inlinecode{n} bits) so that \inlinecode{f[m]} becomes the sum of the original \inlinecode{f[s]} over all submasks \inlinecode{s} of \inlinecode{m}. This ``sum over subsets'' (SOS) DP is the bitmask analog of a prefix sum, accumulating one bit dimension at a time. The value type must support \inlinecode{+=}.
\end{compactitem}

Overflow warning: All intermediate sums must be representable in the value type.

\Needspace*{3\baselineskip}
\textbf{Time Complexity}:
\begin{compactitem}
  \item $O(2^{p})$ per call to \inlinecode{submasks(m)}, where $p$ is \inlinecode{popcount(m)}.
  \item $O(1)$ per call to \inlinecode{next\_same\_popcount()}.
  \item $O(\binom{n}{k})$ per call to \inlinecode{masks\_with\_popcount(n, k)}.
  \item $O(n \cdot 2^{n})$ per call to \inlinecode{subset\_sum\_transform(f)}, where \inlinecode{f} has $2^n$ entries.
\end{compactitem}

\Needspace*{3\baselineskip}
\textbf{Space Complexity}:
\begin{compactitem}
  \item $O(1)$ auxiliary for \inlinecode{next\_same\_popcount()} and \inlinecode{subset\_sum\_transform()}.
  \item $O(1)$ auxiliary and $O(s)$ for the returned masks from \inlinecode{submasks(m)} and \inlinecode{masks\_with\_popcount(n, k)}, where $s$ is the result size.
\end{compactitem}

\begin{lstlisting}[language=C++]
#include <cassert>
#include <climits>
#include <cstdint>
#include <vector>

using Mask = uint32_t;
const int MASK_BITS = sizeof(Mask) * CHAR_BIT;

std::vector<Mask> submasks(Mask m) {
  std::vector<Mask> res;
  Mask s = m;
  while (true) {
    res.push_back(s);
    if (s == 0) {
      break;
    }
    s = (s - 1) & m;
  }
  return res;
}

Mask next_same_popcount(Mask x) {
  assert(x != 0);
  Mask c = x & (Mask{0} - x), r = x + c;
  assert(r != 0);
  return r | (((x ^ r) >> 2) / c);
}

std::vector<Mask> masks_with_popcount(int n, int k) {
  assert(0 <= k && k <= n && n < MASK_BITS);
  std::vector<Mask> res;
  if (k == 0) {
    return {0};
  }
  Mask limit = Mask{1} << n;
  for (Mask x = (Mask{1} << k) - 1; x < limit; x = next_same_popcount(x)) {
    res.push_back(x);
  }
  return res;
}

template<typename T>
void subset_sum_transform(std::vector<T> &f) {
  int size = static_cast<int>(f.size());
  assert(size > 0 && (size & (size - 1)) == 0);
  int n = __builtin_ctz(static_cast<unsigned>(size));  // size = 2^n.
  for (int b = 0; b < n; b++) {
    for (int m = 0; m < size; m++) {
      if (m & (1U << b)) {
        f[m] += f[m ^ (1U << b)];
      }
    }
  }
}
\end{lstlisting}
\Needspace*{4\baselineskip}
\textbf{Example Usage}\par
\begin{lstlisting}[language=C++]
using namespace std;

int main() {
  assert((submasks(0b101) == vector<Mask>{0b101, 0b100, 0b001, 0b000}));

  assert(next_same_popcount(0b0110) == 0b1001U);
  assert(
      (masks_with_popcount(4, 2) == vector<Mask>{0b0011, 0b0101, 0b0110, 0b1001, 0b1010, 0b1100})
  );

  // f[m] = 1 for every mask; after the transform f[m] counts submasks of m = 2^popcount(m).
  vector<int> f(1 << 3, 1);
  subset_sum_transform(f);
  assert(f[0b000] == 1);
  assert(f[0b101] == 4);
  assert(f[0b111] == 8);
  return 0;
}
\end{lstlisting}
\subsection{Bitmask DP}
Bitmask dynamic programming uses a bitmask to represent a small set or frontier, then treats that mask as the DP state. This is useful when one dimension is small enough for $2^n$ states.

\begin{lstlisting}[language=C++]
#include <algorithm>
#include <cassert>
#include <cstdint>
#include <utility>
#include <vector>
\end{lstlisting}
The assignment problem matches workers one-to-one with jobs while minimizing the total cost. The workers are processed in a fixed order: \inlinecode{dp[mask]} is the minimum cost after assigning the first \inlinecode{popcount(mask)} workers to the selected jobs, and the next worker tries every unused job.

\begin{compactitem}
  \item \inlinecode{min\_cost\_assignment(cost)} returns a pair (\inlinecode{sum}, \inlinecode{job}), where \inlinecode{sum} is the minimum cost of assigning each worker to a distinct job, and \inlinecode{job[i]} is the job assigned to worker \inlinecode{i}. The input is a square matrix with one row per worker and one column per job.
\end{compactitem}

Overflow warning: The accumulated cost must fit in \inlinecode{int64\_t}.

\textbf{Time Complexity}: $O(n \cdot 2^{n})$ per call to \inlinecode{min\_cost\_assignment()}, where $n$ is the number of workers and jobs.

\textbf{Space Complexity}: $O(2^{n})$ auxiliary and $O(n)$ for the returned assignment from \inlinecode{min\_cost\_assignment()}.

\begin{lstlisting}[language=C++]
std::pair<int64_t, std::vector<int>> min_cost_assignment(
    const std::vector<std::vector<int64_t>> &cost
) {
  int n = static_cast<int>(cost.size());
  assert(n < 31);
  int states = 1 << n;
  std::vector<int64_t> dp(states);
  std::vector<int> parent(states, -1);
  for (int mask = 0; mask < states; mask++) {
    int worker = __builtin_popcount(static_cast<unsigned>(mask));
    if (worker == n) {
      continue;
    }
    for (int job = 0; job < n; job++) {
      if ((mask & (1 << job)) == 0) {
        int next = mask | (1 << job);
        int64_t candidate = dp[mask] + cost[worker][job];  // Overflow warning.
        if (parent[next] == -1 || dp[next] > candidate) {
          dp[next] = candidate;
          parent[next] = job;
        }
      }
    }
  }
  // Optional: reconstruct one optimal assignment.
  std::vector<int> job(n);
  for (int mask = states - 1, worker = n - 1; worker >= 0; worker--) {
    job[worker] = parent[mask];
    mask ^= 1 << job[worker];
  }
  return {dp[states - 1], job};
}
\end{lstlisting}
The minimum set cover problem asks for the fewest input sets whose union contains every element of a fixed universe. Both input sets and covered subsets are represented as masks: \inlinecode{dp[mask]} is the minimum number of chosen sets needed to cover exactly the elements in \inlinecode{mask}, and adding one set moves to the union mask.

\begin{compactitem}
  \item \inlinecode{set\_cover(sets, universe\_size)} returns a pair (\inlinecode{count}, \inlinecode{chosen}), where \inlinecode{count} is the minimum number of sets needed to cover all elements in $[0, u)$, where $u$ is \inlinecode{universe\_size}, and \inlinecode{chosen} contains one optimal list of set indices. If no cover exists, \inlinecode{count} is $-1$ and \inlinecode{chosen} is empty. Each input set is represented as a bitmask.
\end{compactitem}

\textbf{Time Complexity}: $O(s \cdot 2^{u})$ per call, where $s$ is the number of sets and $u$ is \inlinecode{universe\_size}.

\textbf{Space Complexity}: $O(2^{u})$ auxiliary and $O(s)$ for the returned indices.

\begin{lstlisting}[language=C++]
std::pair<int, std::vector<int>> set_cover(const std::vector<int> &sets, int universe_size) {
  assert(0 <= universe_size && universe_size < 31);
  int states = 1 << universe_size;
  int full = states - 1;
  for (int subset : sets) {
    assert((subset & ~full) == 0);
  }
  int n = static_cast<int>(sets.size());
  int inf = n + 1;
  std::vector<int> dp(states, inf), parent(states, -1), prev(states, -1);
  dp[0] = 0;
  for (int mask = 0; mask < states; mask++) {
    if (dp[mask] == inf) {
      continue;
    }
    for (int i = 0; i < n; i++) {
      int next = mask | sets[i];
      if (dp[next] > dp[mask] + 1) {
        dp[next] = dp[mask] + 1;
        parent[next] = i;
        prev[next] = mask;
      }
    }
  }
  if (dp[full] == inf) {
    return {-1, {}};
  }
  // Optional: reconstruct one minimum set cover.
  std::vector<int> chosen;
  for (int mask = full; mask != 0; mask = prev[mask]) {
    int i = parent[mask];
    chosen.push_back(i);
  }
  std::reverse(chosen.begin(), chosen.end());
  return {dp[full], chosen};
}
\end{lstlisting}
The minimum-cost set partition problem divides all elements into disjoint nonempty groups when the cost of every possible group is known. For each set \inlinecode{mask}, choosing a submask as one group leaves \inlinecode{mask \textasciicircum{} sub} to be partitioned, giving the recurrence \inlinecode{dp[mask] = min(dp[mask \textasciicircum{} sub] + cost[sub])}. Over all masks, this has $O(3^{n})$ transitions.

\begin{compactitem}
  \item \inlinecode{min\_cost\_partition(group\_cost)} returns the minimum total cost to partition all elements into disjoint nonempty groups, where \inlinecode{group\_cost[mask]} is the cost of taking \inlinecode{mask} as one group.
\end{compactitem}

Overflow warning: The accumulated cost must fit in \inlinecode{int64\_t}.

\textbf{Time Complexity}: $O(3^{n})$ per call, where $n$ is the number of elements.

\textbf{Space Complexity}: $O(2^{n})$ auxiliary.

\begin{lstlisting}[language=C++]
int64_t min_cost_partition(const std::vector<int64_t> &group_cost) {
  int states = static_cast<int>(group_cost.size());
  assert(states > 0 && (states & (states - 1)) == 0);
  std::vector<int64_t> dp(states);
  for (int mask = 1; mask < states; mask++) {
    dp[mask] = group_cost[mask];
    for (int sub = (mask - 1) & mask; sub > 0; sub = (sub - 1) & mask) {
      int64_t candidate = dp[mask ^ sub] + group_cost[sub];  // Overflow warning.
      dp[mask] = std::min(dp[mask], candidate);
    }
  }
  return dp[states - 1];
}
\end{lstlisting}
Domino tiling asks for the number of ways to cover a rectangle with horizontal ($1 \times 2$) and vertical ($2 \times 1$) dominoes. Processing cells in row-major order, a profile mask records which cells on the current frontier are already occupied. An occupied current cell is cleared; otherwise, placing a horizontal domino marks the next column, while placing a vertical one carries the current column into the next row. This technique is often called DP on a broken profile or plug DP.

\begin{compactitem}
  \item \inlinecode{count\_domino\_tilings(rows, cols)} returns the number of ways to tile a \inlinecode{rows} by \inlinecode{cols} rectangle with $1 \times 2$ and $2 \times 1$ dominoes.
\end{compactitem}

Overflow warning: The tiling count must fit in \inlinecode{int64\_t}.

\textbf{Time Complexity}: $O(R \cdot C \cdot 2^{C})$ per call, where $R$ and $C$ are the number of rows and columns, respectively.

\textbf{Space Complexity}: $O(2^{C})$ auxiliary.

\begin{lstlisting}[language=C++]
int64_t count_domino_tilings(int rows, int cols) {
  assert(rows >= 0 && cols >= 0);
  if (rows == 0 || cols == 0) {
    return 1;
  }
  if (rows < cols) {
    std::swap(rows, cols);
  }
  assert(cols < 31);
  int states = 1 << cols;
  std::vector<int64_t> dp(states), next(states);
  dp[0] = 1;
  for (int r = 0; r < rows; r++) {
    for (int c = 0; c < cols; c++) {
      next.assign(states, 0);
      for (int mask = 0; mask < states; mask++) {
        if (dp[mask] == 0) {
          continue;
        }
        if (mask & (1 << c)) {
          next[mask ^ (1 << c)] += dp[mask];  // Overflow warning.
        } else {
          if (c + 1 < cols && (mask & (1 << (c + 1))) == 0) {
            next[mask | (1 << (c + 1))] += dp[mask];
          }
          if (r + 1 < rows) {
            next[mask | (1 << c)] += dp[mask];
          }
        }
      }
      dp.swap(next);
    }
  }
  return dp[0];
}
\end{lstlisting}
\Needspace*{4\baselineskip}
\textbf{Example Usage}\par
\begin{lstlisting}[language=C++]
#include <cassert>
using namespace std;

int main() {
  vector<vector<int64_t>> cost{{9, 2, 7}, {6, 4, 3}, {5, 8, 1}};
  auto [assignment_cost, job] = min_cost_assignment(cost);
  assert(assignment_cost == 9 && (job == vector<int>{1, 0, 2}));

  auto [cover_count, chosen] = set_cover(vector<int>{0b0011, 0b0110, 0b1100, 0b1000}, 4);
  assert(cover_count == 2 && (chosen == vector<int>{0, 2}));

  auto [impossible_count, impossible_chosen] = set_cover(vector<int>{0b001, 0b010}, 3);
  assert(impossible_count == -1 && impossible_chosen.empty());

  vector<int64_t> group_cost{
      0,        // Empty group is unused.
      4, 6, 7,  // Singletons.
      6, 2, 3,  // Pairs.
      9,        // All three together.
  };
  assert(min_cost_partition(group_cost) == 7);  // Groups {0} and {1, 2}.

  assert(count_domino_tilings(0, 1000000000) == 1);
  assert(count_domino_tilings(2, 3) == 3);
  assert(count_domino_tilings(3, 3) == 0);
  assert(count_domino_tilings(4, 4) == 36);
  return 0;
}
\end{lstlisting}
\subsection{Gray Code}
The binary reflected Gray code is an ordering of all $2^n$ bitmasks in which consecutive masks differ in exactly one bit. It is useful whenever the cost of moving between adjacent states depends on a single bit flip: enumerating subsets while incrementally maintaining a value, hardware where only one bit may toggle at a time, and constructions like Hamiltonian paths on the hypercube.

The Gray code at 0-based rank \inlinecode{k} is obtained from the ordinary binary value \inlinecode{k} by \inlinecode{k \textasciicircum{} (k \textgreater{}\textgreater{} 1)}. Incrementing \inlinecode{k} flips one binary bit and every lower bit, but XORing each bit with its higher neighbor cancels all but the highest of those changes, so consecutive Gray codes differ in exactly one bit. The map is a bijection: each binary bit is the prefix XOR of the Gray bits at and above it. The inverse loop computes those prefix XORs in doubling steps to recover the original rank.

\begin{compactitem}
  \item \inlinecode{gray\_code(k)} returns the mask at 0-based rank \inlinecode{k} in Gray code order.
  \item \inlinecode{inverse\_gray\_code(g)} returns the 0-based rank \inlinecode{k} such that \inlinecode{gray\_code(k)} equals \inlinecode{g}.
  \item \inlinecode{gray\_sequence(n)} returns all $2^n$ masks in Gray code order; consecutive entries (and the last with the first when \inlinecode{n \textgreater{} 0}) differ in exactly one bit, where $0 \leq {\inlinecodemath{n}} < {\inlinecodemath{MASK\_BITS}}$.
\end{compactitem}

\Needspace*{3\baselineskip}
\textbf{Time Complexity}:
\begin{compactitem}
  \item $O(1)$ per call to \inlinecode{gray\_code()}.
  \item $O(\log b)$ per call to \inlinecode{inverse\_gray\_code()}, where $b$ is \inlinecode{MASK\_BITS}.
  \item $O(2^{n})$ per call to \inlinecode{gray\_sequence(n)}.
\end{compactitem}

\Needspace*{3\baselineskip}
\textbf{Space Complexity}:
\begin{compactitem}
  \item $O(1)$ auxiliary for \inlinecode{gray\_code()} and \inlinecode{inverse\_gray\_code()}.
  \item $O(1)$ auxiliary and $O(2^{n})$ for the returned sequence from \inlinecode{gray\_sequence(n)}.
\end{compactitem}

\begin{lstlisting}[language=C++]
#include <cassert>
#include <climits>
#include <cstdint>
#include <vector>

using Mask = uint32_t;
const int MASK_BITS = sizeof(Mask) * CHAR_BIT;

Mask gray_code(Mask k) {
  return k ^ (k >> 1);
}

Mask inverse_gray_code(Mask g) {
  for (int shift = 1; shift < MASK_BITS && (g >> shift) != 0; shift <<= 1) {
    g ^= g >> shift;
  }
  return g;
}

std::vector<Mask> gray_sequence(int n) {
  assert(0 <= n && n < MASK_BITS);
  std::vector<Mask> res(Mask{1} << n);
  for (Mask k = 0; k < static_cast<Mask>(res.size()); k++) {
    res[k] = gray_code(k);
  }
  return res;
}
\end{lstlisting}
\Needspace*{4\baselineskip}
\textbf{Example Usage}\par
\begin{lstlisting}[language=C++]
#include <cassert>
using namespace std;

int main() {
  assert(gray_code(0) == 0b000U);
  assert(gray_code(1) == 0b001U);
  assert(gray_code(2) == 0b011U);
  assert(gray_code(3) == 0b010U);

  for (Mask k = 0; k < 256; k++) {
    assert(inverse_gray_code(gray_code(k)) == k);
  }
  assert(inverse_gray_code(gray_code(0xdeadbeefU)) == 0xdeadbeefU);

  vector<Mask> seq = gray_sequence(3);
  assert((seq == vector<Mask>{0b000, 0b001, 0b011, 0b010, 0b110, 0b111, 0b101, 0b100}));
  for (int i = 0; i < static_cast<int>(seq.size()); i++) {
    Mask diff = seq[i] ^ seq[(i + 1) % seq.size()];
    assert((diff & (diff - 1)) == 0);  // Cyclically, one bit changes per step.
  }
  return 0;
}
\end{lstlisting}
\subsection{XOR Basis}
An XOR basis (also called a linear basis over the field GF(2)) is the analog of a vector space basis for integers under the bitwise XOR operation. Each integer is treated as a vector of bits, where XOR plays the role of vector addition. Inserting a set of integers one at a time and keeping only those that are linearly independent yields a basis whose XOR-combinations (subset XORs) span exactly the same set of values as the original integers. From a basis of size $r$, exactly $2^r$ distinct values are reachable, and queries such as the maximum attainable XOR become a simple greedy scan.

The basis is kept in row-echelon form: the $i$-th basis element is either $0$ or a value whose highest set bit is bit $i$, and no two basis elements share the same highest bit. Reducing a value means XORing it against the basis elements from the highest bit down, eliminating each set bit that a basis element covers.

\begin{compactitem}
  \item \inlinecode{XorBasis()} constructs an empty basis over \inlinecode{Mask}, which is \inlinecode{uint64\_t} by default.
  \item \inlinecode{insert(x)} adds \inlinecode{x} to the basis, returning \inlinecode{true} if \inlinecode{x} was linearly independent of the current basis (increasing its size), or \inlinecode{false} if \inlinecode{x} was already representable as a subset XOR.
  \item \inlinecode{contains(x)} returns whether \inlinecode{x} is representable as the XOR of some subset of the inserted values, i.e. whether it reduces to $0$ against the basis.
  \item \inlinecode{max\_xor(base = 0)} returns the maximum value of \inlinecode{base} XORed with some subset XOR of the basis. With the default \inlinecode{base} of $0$, this is the maximum subset XOR.
  \item \inlinecode{size()} returns the number of elements in the basis (its rank).
\end{compactitem}

\Needspace*{3\baselineskip}
\textbf{Time Complexity}:
\begin{compactitem}
  \item $O(b)$ per call to \inlinecode{insert()}, \inlinecode{contains()}, and \inlinecode{max\_xor()}, where $b$ is the number of bits.
  \item $O(1)$ per call to \inlinecode{size()}.
\end{compactitem}

\Needspace*{3\baselineskip}
\textbf{Space Complexity}:
\begin{compactitem}
  \item $O(b)$ for storage of the basis.
  \item $O(1)$ auxiliary for all operations.
\end{compactitem}

\begin{lstlisting}[language=C++]
#include <array>
#include <climits>
#include <cstdint>

using Mask = uint64_t;
const int MASK_BITS = sizeof(Mask) * CHAR_BIT;

class XorBasis {
  std::array<Mask, MASK_BITS> basis{};  // basis[i] != 0 has its highest set bit at bit i.
  int basis_size = 0;

 public:
  bool insert(Mask x) {
    for (int i = MASK_BITS - 1; i >= 0; i--) {
      if (!((x >> i) & 1)) {
        continue;
      }
      if (basis[i] == 0) {
        basis[i] = x;
        basis_size++;
        return true;
      }
      x ^= basis[i];
    }
    return false;
  }

  bool contains(Mask x) const {
    for (int i = MASK_BITS - 1; i >= 0; i--) {
      if (((x >> i) & 1) && basis[i] != 0) {
        x ^= basis[i];
      }
    }
    return x == 0;
  }

  Mask max_xor(Mask base = 0) const {
    for (int i = MASK_BITS - 1; i >= 0; i--) {
      if ((base ^ basis[i]) > base) {
        base ^= basis[i];
      }
    }
    return base;
  }

  int size() const { return basis_size; }
};
\end{lstlisting}
\Needspace*{4\baselineskip}
\textbf{Example Usage}\par
\begin{lstlisting}[language=C++]
#include <cassert>

int main() {
  XorBasis b;
  assert(b.insert(1));   // 001
  assert(b.insert(2));   // 010
  assert(!b.insert(3));  // 011 = 1 ^ 2, already representable.
  assert(b.size() == 2);

  assert(b.contains(3));
  assert(b.contains(0));
  assert(!b.contains(4));

  assert(b.max_xor() == 3);   // Best subset XOR over {0, 1, 2, 3}.
  assert(b.max_xor(4) == 7);  // 4 ^ 3.
  return 0;
}
\end{lstlisting}
\subsection{Binary Trie}
Maintain a multiset of nonnegative $b$-bit integers as a binary trie, storing each value as a root-to-leaf path of its bits from the most significant (bit $b - 1$) down to the least significant. Every node keeps a count of the stored values passing through it, which supports deletion and counting queries, and lets the trie answer XOR-extremal queries by a greedy walk down the bits.

The classic application is the maximum-XOR query: to maximize \inlinecode{x \textasciicircum{} y} over all stored \inlinecode{y}, walk from the most significant bit. At each step, descend toward the child whose bit differs from that of \inlinecode{x} whenever such a branch exists, since setting a higher bit of the result always dominates any combination of lower bits. The minimum-XOR query is the mirror image, preferring the matching bit.

This is a multiset (a value may be inserted more than once), and is distinct from the XOR basis: the basis spans subset XORs of a fixed set, whereas this answers XOR queries against the stored elements themselves with support for insertion and deletion.

Values must lie in $[0, 2^b)$ for bit width $b$; the template parameters select the unsigned word type \inlinecode{U} and the bit width \inlinecode{BITS}, which must be less than the width of \inlinecode{U} (for example, $30$ with \inlinecode{uint32\_t} or $62$ with \inlinecode{uint64\_t}). Nodes are kept in a pool indexed by integer, where index $0$ is the root and also serves as the null child.

\begin{compactitem}
  \item \inlinecode{BinaryTrie\textless{}U, BITS\textgreater{}()} constructs an empty multiset.
  \item \inlinecode{size()} returns the number of stored values, and \inlinecode{empty()} returns whether there are none.
  \item \inlinecode{insert(x)} adds one occurrence of \inlinecode{x}.
  \item \inlinecode{count(x)} returns the number of occurrences of \inlinecode{x}.
  \item \inlinecode{erase(x)} removes one occurrence of \inlinecode{x}, returning whether one was present.
  \item \inlinecode{max\_xor(x)} and \inlinecode{min\_xor(x)} return the maximum and minimum of \inlinecode{x \textasciicircum{} y} over all stored \inlinecode{y}. The trie must be non-empty.
  \item \inlinecode{count\_xor\_less(x, k)} returns the number of stored \inlinecode{y} (with multiplicity) such that \inlinecode{x \textasciicircum{} y} is strictly less than \inlinecode{k}.
\end{compactitem}

\Needspace*{3\baselineskip}
\textbf{Time Complexity}:
\begin{compactitem}
  \item $O(b)$ per call to \inlinecode{insert()}, \inlinecode{erase()}, \inlinecode{count()}, \inlinecode{max\_xor()}, \inlinecode{min\_xor()}, and \inlinecode{count\_xor\_less()}, where $b$ is the bit width \inlinecode{BITS}.
  \item $O(1)$ per call to \inlinecode{size()} and \inlinecode{empty()}.
\end{compactitem}

\Needspace*{3\baselineskip}
\textbf{Space Complexity}:
\begin{compactitem}
  \item $O(n \cdot b)$ for storage, where $n$ is the number of distinct values inserted so far.
  \item $O(1)$ auxiliary per operation.
\end{compactitem}

\begin{lstlisting}[language=C++]
#include <array>
#include <cassert>
#include <climits>
#include <cstdint>
#include <vector>

template<typename U = uint32_t, int BITS = 30>
class BinaryTrie {
  static_assert(0 < BITS && BITS < static_cast<int>(sizeof(U) * CHAR_BIT));

  std::vector<std::array<int, 2>> child;  // child[node][bit], where 0 means no child.
  std::vector<int> cnt;                   // Number of stored values passing through each node.

  static void assert_fits(U x) { assert((x >> BITS) == 0); }

  int new_node() {
    child.push_back({0, 0});
    cnt.push_back(0);
    return static_cast<int>(cnt.size()) - 1;
  }

 public:
  BinaryTrie() { new_node(); }  // Node 0 is the root (and doubles as the null child).

  int size() const { return cnt[0]; }
  bool empty() const { return cnt[0] == 0; }

  void insert(U x) {
    assert_fits(x);
    int node = 0;
    cnt[node]++;
    for (int b = BITS - 1; b >= 0; b--) {
      int bit = (x >> b) & 1;
      if (child[node][bit] == 0) {
        int next = new_node();
        child[node][bit] = next;
      }
      node = child[node][bit];
      cnt[node]++;
    }
  }

  int count(U x) const {
    assert_fits(x);
    int node = 0;
    for (int b = BITS - 1; b >= 0; b--) {
      int next = child[node][(x >> b) & 1];
      if (next == 0) {
        return 0;
      }
      node = next;
    }
    return cnt[node];
  }

  bool erase(U x) {
    if (count(x) == 0) {
      return false;
    }
    int node = 0;
    cnt[node]--;
    for (int b = BITS - 1; b >= 0; b--) {
      node = child[node][(x >> b) & 1];
      cnt[node]--;
    }
    return true;
  }

  U max_xor(U x) const {
    assert(!empty());
    assert_fits(x);
    int node = 0;
    U res = 0;
    for (int b = BITS - 1; b >= 0; b--) {
      int bit = (x >> b) & 1;
      int opp = child[node][bit ^ 1];
      if (opp != 0 && cnt[opp] > 0) {
        res |= U{1} << b;
        node = opp;
      } else {
        node = child[node][bit];
      }
    }
    return res;
  }

  U min_xor(U x) const {
    assert(!empty());
    assert_fits(x);
    int node = 0;
    U res = 0;
    for (int b = BITS - 1; b >= 0; b--) {
      int bit = (x >> b) & 1;
      int same = child[node][bit];
      if (same != 0 && cnt[same] > 0) {
        node = same;
      } else {
        res |= U{1} << b;
        node = child[node][bit ^ 1];
      }
    }
    return res;
  }

  int count_xor_less(U x, U k) const {
    assert_fits(x);
    if ((k >> BITS) != 0) {
      return cnt[0];  // Every XOR result is below 2^BITS, which is at most k.
    }
    int node = 0, res = 0;
    for (int b = BITS - 1; b >= 0; b--) {
      int xb = (x >> b) & 1, kb = (k >> b) & 1;
      if (kb == 1) {
        // Stored values matching x at bit b give XOR bit 0 here, hence are already below k.
        int below = child[node][xb];
        if (below != 0) {
          res += cnt[below];
        }
        node = child[node][xb ^ 1];  // Stay tied with k by taking XOR bit 1.
      } else {
        node = child[node][xb];  // XOR bit must be 0 to remain below k.
      }
      if (node == 0) {
        break;
      }
    }
    return res;
  }
};
\end{lstlisting}
\Needspace*{4\baselineskip}
\textbf{Example Usage}\par
\begin{lstlisting}[language=C++]
#include <cassert>
using namespace std;

int main() {
  BinaryTrie<> trie;
  assert(trie.empty());
  for (int x : {3, 10, 5, 25, 2}) {
    trie.insert(x);
  }
  assert(!trie.empty());
  assert(trie.size() == 5);

  // x XOR y over the stored y: {26, 19, 28, 0, 27} for x = 25.
  assert(trie.max_xor(25) == 28);                  // 25 XOR 5.
  assert(trie.min_xor(25) == 0);                   // 25 XOR 25.
  assert(trie.count_xor_less(25, 20) == 2);        // XOR values 0 and 19 are below 20.
  assert(trie.count_xor_less(25, 1U << 30) == 5);  // Bound at/above 2^BITS counts all.

  assert(trie.count(10) == 1);
  trie.insert(10);
  assert(trie.count(10) == 2);  // Multiset: two copies of 10.
  assert(trie.erase(10) && trie.count(10) == 1);
  assert(trie.erase(25) && trie.count(25) == 0);
  assert(!trie.erase(25));        // Already gone.
  assert(trie.max_xor(0) == 10);  // Largest remaining value is 10.
  return 0;
}
\end{lstlisting}
\subsection{Bitwise Convolution (FWHT)}
Bitwise convolution combines two arrays indexed by masks, grouping pairs of masks by a bitwise operation instead of by ordinary addition. For arrays \inlinecode{a} and \inlinecode{b} over an $n$-bit universe, the XOR convolution \inlinecode{c[m]} is the sum of \inlinecode{a[x]*b[y]} over all \inlinecode{x \textasciicircum{} y = m}. OR and AND convolution replace \inlinecode{x \textasciicircum{} y} with \inlinecode{x | y} or \inlinecode{x \& y}. These appear in subset DP, counting pairs by bitwise result, and algebra over the Boolean hypercube.

The fast Walsh-Hadamard transform (FWHT) diagonalizes XOR convolution in the same way that the FFT diagonalizes polynomial convolution: transform both arrays, multiply pointwise, and transform back. The OR and AND versions are Zeta/Mobius transforms over the subset lattice, but they have the same butterfly-shaped implementation and the same $O(n \log n)$ runtime for arrays of length $n = 2^k$.

\begin{compactitem}
  \item \inlinecode{xor\_transform(f, invert)} applies the XOR FWHT to \inlinecode{f} in place. The inverse transform uses \inlinecode{invert = true} and divides every coefficient by \inlinecode{f.size()}.
  \item \inlinecode{or\_transform(f, invert)} applies the OR Zeta transform in place; the inverse is Mobius inversion.
  \item \inlinecode{and\_transform(f, invert)} applies the AND Zeta transform in place; the inverse is Mobius inversion.
  \item \inlinecode{xor\_convolve(a, b)} returns the XOR convolution of \inlinecode{a} and \inlinecode{b}.
  \item \inlinecode{or\_convolve(a, b)} returns the OR convolution of \inlinecode{a} and \inlinecode{b}.
  \item \inlinecode{and\_convolve(a, b)} returns the AND convolution of \inlinecode{a} and \inlinecode{b}.
\end{compactitem}

Subset convolution restricts OR convolution to pairs of disjoint masks, so \inlinecode{c[m]} must count only \inlinecode{x | y = m} with \inlinecode{x \& y = 0}. Rank each entry by its number of set bits, transform every rank layer, and multiply the layers as polynomials in the rank. Mobius inversion then recovers pairs whose union is exactly \inlinecode{m}; among these pairs, the ranks sum to the number of set bits in \inlinecode{m} exactly when the pair is disjoint.

\begin{compactitem}
  \item \inlinecode{subset\_convolve(a, b)} returns the subset convolution of \inlinecode{a} and \inlinecode{b}, where \inlinecode{c[m]} sums \inlinecode{a[x]*b[y]} over the pairs of disjoint masks satisfying \inlinecode{x | y = m}.
\end{compactitem}

The in-place transforms require a nonempty array whose length is a power of two. A convolution returns empty if either input is empty; otherwise, both inputs are padded with zeros to the smallest power of two at least their maximum length, which is also the output length. For exact integer XOR convolution, the inverse divisions are exact; for modular arithmetic, replace the division by multiplication with the modular inverse of the transform length.

Overflow warning: All intermediate sums, differences, and products must be representable in the value type.

\Needspace*{3\baselineskip}
\textbf{Time Complexity}:
\begin{compactitem}
  \item $O(n \log n)$ per transform or bitwise convolution, where $n$ is the padded power-of-two length.
  \item $O(n \log^{2} n)$ per call to \inlinecode{subset\_convolve()}.
\end{compactitem}

\Needspace*{3\baselineskip}
\textbf{Space Complexity}:
\begin{compactitem}
  \item $O(1)$ auxiliary for the in-place transforms.
  \item $O(n)$ auxiliary and $O(n)$ for the returned array from each bitwise convolution.
  \item $O(n \log n)$ auxiliary for \inlinecode{subset\_convolve()}, from the rank layers.
\end{compactitem}

\begin{lstlisting}[language=C++]
#include <algorithm>
#include <cassert>
#include <cstdint>
#include <utility>
#include <vector>

template<typename T>
void xor_transform(std::vector<T> &f, bool invert) {
  int n = static_cast<int>(f.size());
  assert(n > 0 && (n & (n - 1)) == 0);
  for (int len = 1; len < n; len <<= 1) {
    for (int i = 0; i < n; i += len << 1) {
      for (int j = 0; j < len; j++) {
        T u = f[i + j], v = f[i + j + len];
        f[i + j] = u + v;  // Overflow warning.
        f[i + j + len] = u - v;
      }
    }
  }
  if (invert) {
    for (T &x : f) {
      x /= n;
    }
  }
}

template<typename T>
void or_transform(std::vector<T> &f, bool invert) {
  int n = static_cast<int>(f.size());
  assert(n > 0 && (n & (n - 1)) == 0);
  for (int bit = 1; bit < n; bit <<= 1) {
    for (int mask = 0; mask < n; mask++) {
      if (mask & bit) {
        if (invert) {
          f[mask] -= f[mask ^ bit];  // Overflow warning.
        } else {
          f[mask] += f[mask ^ bit];
        }
      }
    }
  }
}

template<typename T>
void and_transform(std::vector<T> &f, bool invert) {
  int n = static_cast<int>(f.size());
  assert(n > 0 && (n & (n - 1)) == 0);
  for (int bit = 1; bit < n; bit <<= 1) {
    for (int mask = 0; mask < n; mask++) {
      if ((mask & bit) == 0) {
        if (invert) {
          f[mask] -= f[mask ^ bit];  // Overflow warning.
        } else {
          f[mask] += f[mask ^ bit];
        }
      }
    }
  }
}

template<typename T, typename Transform>
std::vector<T> bitwise_convolve(std::vector<T> a, std::vector<T> b, Transform transform) {
  if (a.empty() || b.empty()) {
    return {};
  }
  int needed = static_cast<int>(std::max(a.size(), b.size()));
  int n = 1;
  while (n < needed) {
    n <<= 1;
  }
  a.resize(n);
  b.resize(n);
  transform(a, false);
  transform(b, false);
  for (int i = 0; i < n; i++) {
    a[i] *= b[i];  // Overflow warning.
  }
  transform(a, true);
  return a;
}

template<typename T>
std::vector<T> xor_convolve(std::vector<T> a, std::vector<T> b) {
  return bitwise_convolve(std::move(a), std::move(b), xor_transform<T>);
}

template<typename T>
std::vector<T> or_convolve(std::vector<T> a, std::vector<T> b) {
  return bitwise_convolve(std::move(a), std::move(b), or_transform<T>);
}

template<typename T>
std::vector<T> and_convolve(std::vector<T> a, std::vector<T> b) {
  return bitwise_convolve(std::move(a), std::move(b), and_transform<T>);
}

template<typename T>
std::vector<T> subset_convolve(std::vector<T> a, std::vector<T> b) {
  if (a.empty() || b.empty()) {
    return {};
  }
  int needed = static_cast<int>(std::max(a.size(), b.size()));
  int n = 1, bits = 0;
  while (n < needed) {
    n <<= 1;
    bits++;
  }
  a.resize(n);
  b.resize(n);
  std::vector<int> ranks(n);
  for (int mask = 1; mask < n; mask++) {
    ranks[mask] = ranks[mask >> 1] + (mask & 1);
  }
  // Split each input into layers by rank, so that layer k holds only the masks with k set bits.
  std::vector<std::vector<T>> fa(bits + 1, std::vector<T>(n)), fb = fa, fc = fa;
  for (int mask = 0; mask < n; mask++) {
    fa[ranks[mask]][mask] = a[mask];
    fb[ranks[mask]][mask] = b[mask];
  }
  for (int k = 0; k <= bits; k++) {
    or_transform(fa[k], false);
    or_transform(fb[k], false);
  }
  // Multiplying layers groups pairs by total rank. Once Mobius inversion recovers each pair's exact
  // union, total rank ranks[mask] selects precisely the disjoint pairs covering mask.
  for (int i = 0; i <= bits; i++) {
    for (int j = 0; i + j <= bits; j++) {
      for (int mask = 0; mask < n; mask++) {
        fc[i + j][mask] += fa[i][mask] * fb[j][mask];  // Overflow warning.
      }
    }
  }
  std::vector<T> res(n);
  for (int k = 0; k <= bits; k++) {
    or_transform(fc[k], true);
  }
  for (int mask = 0; mask < n; mask++) {
    res[mask] = fc[ranks[mask]][mask];
  }
  return res;
}
\end{lstlisting}
\Needspace*{4\baselineskip}
\textbf{Example Usage}\par
\begin{lstlisting}[language=C++]
#include <cassert>
using namespace std;

int main() {
  {
    vector<int64_t> f{1, 2, 3, 4}, original(f);
    xor_transform(f, false);
    xor_transform(f, true);
    assert(f == original);
  }
  {
    vector<int64_t> f{1, 2, 3, 4}, original(f);
    or_transform(f, false);
    or_transform(f, true);
    assert(f == original);
  }
  {
    vector<int64_t> f{1, 2, 3, 4}, original(f);
    and_transform(f, false);
    and_transform(f, true);
    assert(f == original);
  }

  vector<int64_t> a{1, 2, 3, 4}, b{5, 6, 7, 8};

  // XOR: c[m] sums a[x] * b[y] over pairs with x ^ y == m.
  assert((xor_convolve(a, b) == vector<int64_t>{70, 68, 62, 60}));

  // OR: c[m] sums a[x] * b[y] over pairs with x | y == m.
  assert((or_convolve(a, b) == vector<int64_t>{5, 28, 43, 184}));

  // AND: c[m] sums a[x] * b[y] over pairs with x & y == m.
  assert((and_convolve(a, b) == vector<int64_t>{103, 52, 73, 32}));

  // Subset: c[m] sums a[x] * b[y] over disjoint pairs with x | y == m. Every entry is at most the
  // corresponding OR convolution, which also counts the overlapping pairs.
  assert((subset_convolve(a, b) == vector<int64_t>{5, 16, 22, 60}));

  vector<int64_t> one_mask{0, 1, 0, 0};
  vector<int64_t> two_mask{0, 0, 1};
  assert((xor_convolve(one_mask, two_mask) == vector<int64_t>{0, 0, 0, 1}));
  assert((or_convolve(one_mask, two_mask) == vector<int64_t>{0, 0, 0, 1}));
  assert((and_convolve(one_mask, two_mask) == vector<int64_t>{1, 0, 0, 0}));
  assert((subset_convolve(one_mask, two_mask) == vector<int64_t>{0, 0, 0, 1}));

  // Overlapping masks contribute nothing at all to a subset convolution.
  assert((subset_convolve(one_mask, one_mask) == vector<int64_t>{0, 0, 0, 0}));
  return 0;
}
\end{lstlisting}

\section{\texorpdfstring{Classic Problems}{1.8 Classic Problems}}
\setcounter{section}{8}
\setcounter{subsection}{0}
\subsection{Minimum Excluded Value}
Given a multiset of integers, find the minimum excluded nonnegative value (MEX), i.e. the smallest integer $x \geq 0$ not present in the multiset. MEX often appears in array problems, game theory with Grundy numbers, and dynamic programming states. Since the MEX of $n$ values never exceeds $n$, the one-shot version tracks seen values in $[0, n]$ before scanning to find the first missing.

\begin{compactitem}
  \item \inlinecode{mex(lo, hi)} returns the MEX of the values in $[{\inlinecodemath{lo}}, {\inlinecodemath{hi}})$.
\end{compactitem}

The MEX of a dynamic multiset can be computed using counts of present values and maintaining an ordered set of currently missing candidates.

\begin{compactitem}
  \item \inlinecode{DynamicMex()} maintains the MEX of a multiset of nonnegative integers.
  \item \inlinecode{add(x)} inserts one copy of value \inlinecode{x} if \inlinecode{x} $\geq 0$.
  \item \inlinecode{remove(x)} removes one copy of value \inlinecode{x} if present.
  \item \inlinecode{mex()} returns the current MEX.
\end{compactitem}

\Needspace*{3\baselineskip}
\textbf{Time Complexity}:
\begin{compactitem}
  \item $O(n)$ per call to \inlinecode{mex(lo, hi)}, where $n$ is the number of values.
  \item $O(\log n)$ expected per call to \inlinecode{add(x)} and \inlinecode{remove(x)} for \inlinecode{DynamicMex}, and $O(n)$ in the collision-heavy worst case for the hash table.
  \item $O(1)$ per call to \inlinecode{mex()} for \inlinecode{DynamicMex}.
\end{compactitem}

\Needspace*{3\baselineskip}
\textbf{Space Complexity}:
\begin{compactitem}
  \item $O(n)$ auxiliary for \inlinecode{mex(lo, hi)}.
  \item $O(n)$ storage for \inlinecode{DynamicMex}.
\end{compactitem}

\begin{lstlisting}[language=C++]
#include <iterator>
#include <set>
#include <unordered_map>
#include <vector>

template<typename It>
int mex(It lo, It hi) {
  int n = std::distance(lo, hi);
  std::vector<char> seen(n + 1);
  for (It it = lo; it != hi; ++it) {
    if (0 <= *it && *it <= n) {
      seen[*it] = true;
    }
  }
  for (int x = 0; x <= n; x++) {
    if (!seen[x]) {
      return x;
    }
  }
  return n + 1;
}

class DynamicMex {
  std::unordered_map<int, int> count;
  std::set<int> missing;  // Missing candidates in [0, size].
  int size = 0;

 public:
  DynamicMex() { missing.insert(0); }

  void add(int x) {
    if (x < 0) {
      return;
    }
    if (!count.count(size + 1)) {
      missing.insert(size + 1);
    }
    if (count[x]++ == 0) {
      missing.erase(x);
    }
    size++;
  }

  void remove(int x) {
    auto it = count.find(x);
    if (it == count.end()) {
      return;
    }
    bool removed_last = --it->second == 0;
    if (removed_last) {
      count.erase(it);
    }
    size--;
    missing.erase(size + 1);
    if (removed_last && x <= size) {
      missing.insert(x);
    }
  }

  int mex() const { return *missing.begin(); }
};
\end{lstlisting}
\Needspace*{4\baselineskip}
\textbf{Example Usage}\par
\begin{lstlisting}[language=C++]
#include <cassert>
using namespace std;

int main() {
  vector<int> a{0, 1, 4, 2, 1};
  assert(mex(a.begin(), a.end()) == 3);
  vector<int> b{-1, 0, 1, 2};
  assert(mex(b.begin(), b.end()) == 3);  // Negative values are ignored.

  DynamicMex m;
  m.add(0);
  m.add(1);
  m.add(1);
  m.add(3);
  assert(m.mex() == 2);
  m.add(2);
  assert(m.mex() == 4);
  m.remove(1);
  assert(m.mex() == 4);  // One copy of 1 remains.
  m.remove(1);
  assert(m.mex() == 1);
  m.remove(100);
  assert(m.mex() == 1);
  return 0;
}
\end{lstlisting}
\subsection{Josephus Problem}
The Josephus problem arranges $n$ people, numbered $[0, n)$, in a circle and repeatedly removes every $k$-th remaining person until one survives. Counting starts at person $0$, which counts as the first person. After the first removal, relabeling the smaller circle from the next person gives the same problem on $n - 1$ people. Mapping its survivor back to the original labels yields the recurrence $J(1, k) = 0$ and $J(n, k) = (J(n - 1, k) + k) \bmod n$.

When $k$ is small, several removals can be processed at once. One complete pass removes $\lfloor n/k \rfloor$ people; recursively solving the compressed circle and undoing the resulting index shift reduces the running time to $O(k \log n)$. To recover the full elimination order, a Fenwick tree stores which labels remain and selects each next victim by rank.

\begin{compactitem}
  \item \inlinecode{josephus(n, k)} returns the 0-based label of the survivor, where \inlinecode{n} and \inlinecode{k} must be positive.
  \item \inlinecode{josephus\_small\_k(n, k)} returns the same survivor using the batched recurrence. It is preferable when $k$ is small relative to $n$.
  \item \inlinecode{josephus\_order(n, k)} returns all labels in elimination order, ending with the survivor.
\end{compactitem}

\Needspace*{3\baselineskip}
\textbf{Time Complexity}:
\begin{compactitem}
  \item $O(n)$ per call to \inlinecode{josephus()}.
  \item $O(\min\left(n, k \log n\right))$ per call to \inlinecode{josephus\_small\_k()}.
  \item $O(n \log n)$ per call to \inlinecode{josephus\_order()}.
\end{compactitem}

\Needspace*{3\baselineskip}
\textbf{Space Complexity}:
\begin{compactitem}
  \item $O(1)$ auxiliary for \inlinecode{josephus()}.
  \item $O(\min\left(n, k \log n\right))$ call stack space for \inlinecode{josephus\_small\_k()}.
  \item $O(n)$ auxiliary and $O(n)$ for the returned elimination order from \inlinecode{josephus\_order()}.
\end{compactitem}

\begin{lstlisting}[language=C++]
#include <cassert>
#include <cstdint>
#include <vector>

int josephus(int n, int k) {
  assert(n > 0 && k > 0);
  int survivor = 0;
  for (int size = 2; size <= n; size++) {
    survivor = static_cast<int>((static_cast<int64_t>(survivor) + k) % size);
  }
  return survivor;
}

int josephus_small_k(int n, int k) {
  assert(n > 0 && k > 0);
  if (n == 1) {
    return 0;
  }
  if (k == 1) {
    return n - 1;
  }
  if (k > n) {
    return josephus(n, k);
  }
  int removed = n / k;
  int survivor = josephus_small_k(n - removed, k);
  survivor -= n % k;
  if (survivor < 0) {
    survivor += n;
  } else {
    survivor += survivor / (k - 1);
  }
  return survivor;
}

std::vector<int> josephus_order(int n, int k) {
  assert(n > 0 && k > 0);
  std::vector<int> bit(n + 1), order;
  order.reserve(n);
  for (int i = 1; i <= n; i++) {
    bit[i] = i & -i;  // Initially every label is alive.
  }
  int max_step = 1;
  while (max_step <= n / 2) {
    max_step *= 2;
  }
  auto find_by_rank = [&](int rank) {
    int pos = 0;
    for (int step = max_step; step > 0; step /= 2) {
      int next = pos + step;
      if (next <= n && bit[next] <= rank) {
        pos = next;
        rank -= bit[next];
      }
    }
    return pos;
  };
  int rank = 0;
  for (int remaining = n; remaining > 0; remaining--) {
    rank = static_cast<int>((static_cast<int64_t>(rank) + k - 1) % remaining);
    int removed = find_by_rank(rank);
    order.push_back(removed);
    for (int i = removed + 1; i <= n; i += i & -i) {
      bit[i]--;
    }
    if (remaining > 1) {
      rank %= remaining - 1;
    }
  }
  return order;
}
\end{lstlisting}
\Needspace*{4\baselineskip}
\textbf{Example Usage}\par
\begin{lstlisting}[language=C++]
#include <cassert>
#include <vector>
using namespace std;

int main() {
  assert(josephus(1, 5) == 0);
  assert(josephus(7, 3) == 3);
  assert(josephus(10, 1) == 9);
  for (int n = 1; n <= 30; n++) {
    for (int k = 1; k <= 30; k++) {
      assert(josephus_small_k(n, k) == josephus(n, k));
      assert(josephus_order(n, k).back() == josephus(n, k));
    }
  }
  assert(josephus_order(7, 3) == vector<int>({2, 5, 1, 6, 4, 0, 3}));
  return 0;
}
\end{lstlisting}
\subsection{Tower of Hanoi}
Move a stack of $n$ disks of distinct sizes from one peg to another using a third as scratch space, moving one disk at a time and never placing a larger disk on a smaller one. The recursive solution is forced by the largest disk: it can only move once the other $n - 1$ disks are stacked on the spare peg, so the puzzle is solved by moving $n - 1$ disks aside, moving the largest disk across, and moving the $n - 1$ disks back on top of it. That gives $T(n) = 2T(n - 1) + 1$ with $T(0) = 0$, so exactly $2^n - 1$ moves are needed, and no shorter solution exists because the largest disk cannot move any earlier.

The sequence also has a closed form, for when one move is wanted rather than all of them: numbering moves from $1$, move $k$ transfers the disk given by the number of trailing zero bits of $k$, and that disk always steps in one rotational direction fixed by the parity of $n$. The smallest disk therefore moves every other turn, which is the rule an iterative solver follows.

\begin{compactitem}
  \item \inlinecode{hanoi\_moves(n, from = 0, to = 2)} returns the moves solving the puzzle for \inlinecode{n} disks, as pairs (\inlinecode{from\_peg}, \inlinecode{to\_peg}). Pegs are numbered $[0, 3)$ and the spare peg is whichever of the three is neither \inlinecode{from} nor \inlinecode{to}. The number of disks must be nonnegative.
  \item \inlinecode{hanoi\_move\_count(n)} returns $2^n - 1$, the number of moves, without generating them. It requires $n \leq 63$, the largest count that fits in \inlinecode{int64\_t}.
\end{compactitem}

With four or more pegs the problem becomes Frame-Stewart: move a prefix aside using every peg, transfer the rest with one peg fewer, bring the prefix back, and minimize over the split. That this family of strategies is optimal was only proved in 2014.

\Needspace*{3\baselineskip}
\textbf{Time Complexity}:
\begin{compactitem}
  \item $O(2^{n})$ per call to \inlinecode{hanoi\_moves()}, which is optimal since that is the size of the output.
  \item $O(1)$ per call to \inlinecode{hanoi\_move\_count()}.
\end{compactitem}

\Needspace*{3\baselineskip}
\textbf{Space Complexity}:
\begin{compactitem}
  \item $O(n)$ auxiliary stack space and $O(2^{n})$ for the returned moves from \inlinecode{hanoi\_moves()}.
  \item $O(1)$ auxiliary for \inlinecode{hanoi\_move\_count()}.
\end{compactitem}

\begin{lstlisting}[language=C++]
#include <cassert>
#include <cstdint>
#include <utility>
#include <vector>

std::vector<std::pair<int, int>> hanoi_moves(int n, int from = 0, int to = 2) {
  assert(n >= 0 && from >= 0 && from < 3 && to >= 0 && to < 3 && from != to);
  std::vector<std::pair<int, int>> moves;
  auto rec = [&](auto &&rec, int disks, int src, int dest) -> void {
    if (disks == 0) {
      return;
    }
    int spare = 3 - src - dest;
    rec(rec, disks - 1, src, spare);
    moves.emplace_back(src, dest);
    rec(rec, disks - 1, spare, dest);
  };
  rec(rec, n, from, to);
  return moves;
}

int64_t hanoi_move_count(int n) {
  assert(n >= 0 && n <= 63);
  return static_cast<int64_t>((uint64_t{1} << n) - 1);  // Unsigned, so n = 63 does not overflow.
}
\end{lstlisting}
\Needspace*{4\baselineskip}
\textbf{Example Usage}\par
\begin{lstlisting}[language=C++]
#include <cassert>
using namespace std;

int main() {
  assert(hanoi_moves(0).empty());
  assert((hanoi_moves(1) == vector<pair<int, int>>{{0, 2}}));
  assert((hanoi_moves(2) == vector<pair<int, int>>{{0, 1}, {0, 2}, {1, 2}}));
  assert(hanoi_move_count(10) == 1023);

  // Replay the moves of a larger puzzle to confirm no disk ever lands on a smaller one.
  const int n = 12;
  vector<vector<int>> pegs(3);
  for (int disk = n; disk >= 1; disk--) {
    pegs[0].push_back(disk);
  }
  auto moves = hanoi_moves(n);
  assert(static_cast<int64_t>(moves.size()) == hanoi_move_count(n));
  for (auto [src, dest] : moves) {
    assert(!pegs[src].empty());
    assert(pegs[dest].empty() || pegs[dest].back() > pegs[src].back());
    pegs[dest].push_back(pegs[src].back());
    pegs[src].pop_back();
  }
  assert(pegs[0].empty() && pegs[1].empty() && static_cast<int>(pegs[2].size()) == n);
  return 0;
}
\end{lstlisting}
\subsection{Cycle Detection (Floyd)}
Given a function $f$ mapping a finite set to itself and a starting value $x_0$, return the entry value, position, and length of the cycle reached by repeatedly applying $f$. Formally, the sequence $x_0$, $x_1 = f(x_0)$, $x_2 = f(x_1)$, $\ldots$, $x_n = f(x_{n - 1})$, $\ldots$ must eventually repeat some value. If $x_i = x_j$ for $i < j$, then the sequence from $x_i$ to $x_{j - 1}$ repeats forever. This is useful for detecting cycles in functional graphs (see 4.2.9), pseudo-random generators, Pollard rho style algorithms, degenerate linked lists, and arrays where each index points to the next index. For example, if an array $a_0, \ldots, a_{n+1}$ contains only values in $[1, n]$, with exactly one value occurring more than once, then $f(i) = a_i$ forms a functional graph; the duplicate value is the entry point of the cycle reached from $0$.

Floyd's cycle-finding algorithm, a.k.a. the ``tortoise and the hare algorithm'', is a space-efficient default choice: it moves two pointers through the sequence at different speeds, finds a meeting point inside the cycle, then locates the cycle start. Brent's algorithm is a compact variant that usually calls $f$ fewer times by only resetting the tortoise at powers of two; prefer it when evaluating $f$ is expensive and constant factors matter. In the detection phase, Brent advances one pointer with one call to $f$ per step, while Floyd advances the tortoise once and the hare twice, using three calls to $f$ per step.

For Floyd's algorithm, suppose the tail has length $m$ and the cycle has length $n$. At a meeting, the hare has traveled twice as far as the tortoise, so the tortoise's distance is a multiple of $n$. Consequently, the distance from the meeting point around the cycle to its entry is congruent to $m$ modulo $n$. Resetting one pointer to $x_0$ and advancing both one step at a time makes them meet at the cycle entry after $m$ steps. Brent first determines $n$, then advances one pointer $n$ steps; moving both together from that fixed separation likewise makes them meet at the entry.

\begin{compactitem}
  \item \inlinecode{find\_cycle\_floyd(f, x0)} repeatedly applies \inlinecode{f} starting from \inlinecode{x0} and returns the reached cycle as a tuple (\inlinecode{entry}, \inlinecode{start}, \inlinecode{length}), where \inlinecode{entry} is the first cycle value, \inlinecode{start} is the number of applications of \inlinecode{f} needed to reach it from \inlinecode{x0}, and \inlinecode{length} is the number of distinct values in the cycle.
  \item \inlinecode{find\_cycle\_brent(f, x0)} does the same with fewer calls to \inlinecode{f} in many cases.
\end{compactitem}

\textbf{Time Complexity}: $O(m + n)$ per call, where $m$ is the first index in the cycle and $n$ is the cycle length.

\textbf{Space Complexity}: $O(1)$ auxiliary.

\begin{lstlisting}[language=C++]
#include <tuple>

template<typename Fn, typename T>
std::tuple<T, int, int> find_cycle_floyd(Fn f, T x0) {
  // Find a meeting point inside the cycle.
  T tortoise = f(x0), hare = f(f(x0));
  while (tortoise != hare) {
    tortoise = f(tortoise);
    hare = f(f(hare));
  }
  // Find the cycle entry.
  int start = 0;
  tortoise = x0;
  while (tortoise != hare) {
    tortoise = f(tortoise);
    hare = f(hare);
    start++;
  }
  // Measure the cycle length.
  int length = 1;
  hare = f(tortoise);
  while (tortoise != hare) {
    hare = f(hare);
    length++;
  }
  return {tortoise, start, length};
}

template<typename Fn, typename T>
std::tuple<T, int, int> find_cycle_brent(Fn f, T x0) {
  // Find the cycle length by doubling the search block.
  int power = 1, length = 1;
  T tortoise = x0, hare = f(x0);
  while (tortoise != hare) {
    if (power == length) {
      tortoise = hare;
      power *= 2;
      length = 0;
    }
    hare = f(hare);
    length++;
  }
  // Separate the pointers by one cycle length.
  hare = x0;
  for (int i = 0; i < length; i++) {
    hare = f(hare);
  }
  // Find the cycle entry.
  int start = 0;
  tortoise = x0;
  while (tortoise != hare) {
    tortoise = f(tortoise);
    hare = f(hare);
    start++;
  }
  return {tortoise, start, length};
}
\end{lstlisting}
\Needspace*{4\baselineskip}
\textbf{Example Usage}\par
\begin{lstlisting}[language=C++]
#include <cassert>
#include <set>
#include <vector>
using namespace std;

int f(int x) {
  return (123 * x * x + 4567890) % 1337;
}

void verify(int x0, int start, int length) {
  set<int> s;
  int x = x0;
  for (int i = 0; i < start; i++) {
    assert(!s.count(x));
    s.insert(x);
    x = f(x);
  }
  int startx = x;
  s.clear();
  for (int i = 0; i < length; i++) {
    assert(!s.count(x));
    s.insert(x);
    x = f(x);
  }
  assert(startx == x);
}

// Given n + 1 integers within [1, n], where one value in [1, n] is present 2 or more times (but
// other values in [1, n] are present at most once), find the duplicate without modifying the array.
// Interpret each value as a pointer to the next index; the duplicate is the cycle entry.
int find_duplicate_integer(const vector<int> &a) {
  return get<0>(find_cycle_floyd([&](int i) { return a[i]; }, 0));
}

int main() {
  int x0 = 0;
  auto [floyd_entry, floyd_start, floyd_len] = find_cycle_floyd(f, x0);
  assert(floyd_start == 4 && floyd_len == 2);
  verify(x0, floyd_start, floyd_len);

  auto [brent_entry, brent_start, brent_len] = find_cycle_brent(f, x0);
  assert(brent_entry == floyd_entry && brent_start == floyd_start && brent_len == floyd_len);

  assert(find_duplicate_integer({1, 3, 4, 2, 2}) == 2);
  assert(find_duplicate_integer({3, 1, 3, 4, 2}) == 3);
  return 0;
}
\end{lstlisting}
\subsection{Subset Sum}
Given a collection of integers and a target, the subset sum problem asks for the maximum sum of any subset that does not exceed the target. Two algorithms below solve this problem in complementary regimes, and which one wins depends on whether the bottleneck is the number of items or the size of the weights.

The meet-in-the-middle method splits the collection in half, enumerates every subset sum for each half, sorts one side, and binary searches for the best compatible partner. It makes no assumption about the sign or size of the values, so it handles negative inputs, and it is the method of choice when the number of items $n$ is small (up to roughly $40$).

Pisinger's algorithm targets the opposite regime: nonnegative weights whose maximum value is small, with possibly many items. Let $W$ be the largest weight. It greedily takes a maximal prefix with sum $s$ that fits the target; because the next weight is at most $W$ and does not fit, $s$ is greater than \inlinecode{target} $- W$. The prefix itself is feasible, so the optimum lies between $s$ and \inlinecode{target} and only that width-$W$ window matters. A balanced sweep explores exchanges that add later items to sums below the target and remove earlier items from sums above it. For each sum, it retains the farthest prefix boundary reached, which dominates smaller boundaries because it permits at least the same future exchanges. This gives a running time linear in the number of items times $W$, independent of the target beyond that window. This routine is sometimes called ``fast knapsack,'' a slight misnomer: it maximizes a subset sum rather than packing items with distinct values, so it is really the bounded-weight case of subset sum rather than the 0-1 value knapsack.

\begin{compactitem}
  \item \inlinecode{max\_subset\_sum\_at\_most(values, target)} returns a pair (\inlinecode{sum}, \inlinecode{items}) containing that maximum sum and the selected item indices. Values may be negative, and \inlinecode{target} must be nonnegative.
  \item \inlinecode{max\_subset\_sum\_bounded(weights, target)} returns the maximum sum of any subset of \inlinecode{weights} that does not exceed \inlinecode{target}. All weights and \inlinecode{target} must be nonnegative integers. This is faster than meet-in-the-middle when the largest weight is small relative to $2^{n/2}$.
\end{compactitem}

Overflow warning: All subset sums and intermediate differences in \inlinecode{max\_subset\_sum\_at\_most()} must fit in \inlinecode{int64\_t}.

\Needspace*{3\baselineskip}
\textbf{Time Complexity}:
\begin{compactitem}
  \item $O(n \cdot 2^{n/2})$ per call to \inlinecode{max\_subset\_sum\_at\_most()}, where $n$ is the number of values.
  \item $O(n \cdot W)$ per call to \inlinecode{max\_subset\_sum\_bounded()}, where $n$ is the number of items and $W$ is the largest weight.
\end{compactitem}

\Needspace*{3\baselineskip}
\textbf{Space Complexity}:
\begin{compactitem}
  \item $O(2^{n/2})$ auxiliary for \inlinecode{max\_subset\_sum\_at\_most()}.
  \item $O(W)$ auxiliary for \inlinecode{max\_subset\_sum\_bounded()}.
\end{compactitem}

\begin{lstlisting}[language=C++]
#include <algorithm>
#include <cassert>
#include <cstdint>
#include <utility>
#include <vector>

std::pair<int64_t, std::vector<int>> max_subset_sum_at_most(
    const std::vector<int> &values, int64_t target
) {
  assert(target >= 0);
  int n = static_cast<int>(values.size());
  assert(n / 2 < 63 && n - n / 2 < 63);
  int64_t llen = 1LL << (n / 2), rlen = 1LL << (n - n / 2);
  std::vector<int64_t> lsum(llen);
  std::vector<std::pair<int64_t, int64_t>> rsum(rlen);
  for (int64_t mask = 1; mask < llen; mask++) {
    int bit = __builtin_ctzll(mask);
    lsum[mask] = lsum[mask ^ (1LL << bit)] + values[bit];  // Overflow warning.
  }
  for (int64_t mask = 1; mask < rlen; mask++) {
    int bit = __builtin_ctzll(mask);
    rsum[mask].first = rsum[mask ^ (1LL << bit)].first + values[n / 2 + bit];
    rsum[mask].second = mask;
  }
  std::sort(rsum.begin(), rsum.end());
  int64_t best = INT64_MIN, lmask = 0, rmask = 0;
  for (int64_t mask = 0; mask < llen; mask++) {
    int64_t limit = target - lsum[mask];  // Overflow warning.
    auto it = std::upper_bound(rsum.begin(), rsum.end(), std::make_pair(limit, INT64_MAX));
    if (it != rsum.begin()) {
      --it;
      int64_t candidate = lsum[mask] + it->first;
      if (best < candidate) {
        best = candidate;
        lmask = mask;
        rmask = it->second;
      }
    }
  }
  // Optional: reconstruct one optimal subset.
  std::vector<int> items;
  for (int i = 0; i < n / 2; i++) {
    if ((lmask >> i) & 1) {
      items.push_back(i);
    }
  }
  for (int i = 0; i < n - n / 2; i++) {
    if ((rmask >> i) & 1) {
      items.push_back(n / 2 + i);
    }
  }
  return {best, items};
}

int max_subset_sum_bounded(const std::vector<int> &weights, int target) {
  assert(target >= 0);
  for (int w : weights) {
    assert(w >= 0);
  }
  int n = static_cast<int>(weights.size());
  // Greedily take a prefix while it still fits; the optimum then lies within one weight of target.
  int sum = 0, b = 0;
  while (b < n && weights[b] <= target - sum) {
    sum += weights[b++];
  }
  if (b == n) {
    return sum;  // Every item fits, so the whole collection is the best subset.
  }
  int m = 0;
  for (int w : weights) {
    m = std::max(m, w);
  }
  // reach[s - target + m] = largest prefix boundary that can realize a balanced subset of sum s.
  // Only sums within m of target are tracked, which is where the optimum must lie.
  std::vector<int> reach(2 * m, -1);
  reach[sum - target + m] = b;
  for (int i = b; i < n; i++) {
    std::vector<int> prev = reach;
    for (int x = 0; x < m; x++) {  // Add item i to sums currently below target.
      reach[x + weights[i]] = std::max(reach[x + weights[i]], prev[x]);
    }
    for (int x = 2 * m - 1; x > m; x--) {  // Drop an earlier item from sums above target.
      for (int j = std::max(0, prev[x]); j < reach[x]; j++) {
        reach[x - weights[j]] = std::max(reach[x - weights[j]], j);
      }
    }
  }
  int s = target;
  while (reach[s - target + m] < 0) {
    s--;
  }
  return s;
}
\end{lstlisting}
\Needspace*{4\baselineskip}
\textbf{Example Usage}\par
\begin{lstlisting}[language=C++]
#include <cassert>
using namespace std;

int main() {
  vector<int> a{9, 1, 5, 0, 1, 11, 5};
  auto [sum, items] = max_subset_sum_at_most(a, 8);
  int64_t selected_sum = 0;
  for (int i : items) {
    selected_sum += a[i];
  }
  assert(sum == 7 && selected_sum == sum);
  vector<int> b{-7, -3, -2, 5, 8};
  assert(max_subset_sum_at_most(b, 0).first == 0);

  // The bounded-weight method agrees with meet-in-the-middle on nonnegative inputs.
  assert(max_subset_sum_bounded(a, 8) == 7);
  vector<int> c{3, 34, 4, 12, 5, 2};
  assert(max_subset_sum_bounded(c, 9) == 9);     // 4 + 5.
  assert(max_subset_sum_bounded(c, 10) == 10);   // 3 + 5 + 2.
  assert(max_subset_sum_bounded(c, 100) == 60);  // All items fit.
  return 0;
}
\end{lstlisting}
\subsection{N-Queens}
The N-queens problem asks for a placement of $n$ queens on an $n$-by-$n$ chessboard such that no two share a row, column, or diagonal. Backtracking places one queen in each successive row and tracks which columns and diagonals are occupied. A conflicting placement is rejected immediately; after a failed recursive branch, its column and diagonals are released before trying the next column.

\begin{compactitem}
  \item \inlinecode{n\_queens(n)} returns a placement as a vector \inlinecode{cols} of length \inlinecode{n}, where \inlinecode{cols[row]} is the queen's column in that row. It returns an empty vector if no placement exists.
\end{compactitem}

\textbf{Time Complexity}: $O(n!)$ per call.

\textbf{Space Complexity}: $O(n)$ auxiliary and $O(n)$ for the returned placement.

\begin{lstlisting}[language=C++]
#include <cassert>
#include <vector>

std::vector<int> n_queens(int n) {
  assert(n > 0);
  std::vector<int> cols(n);
  std::vector<char> used_col(n), used_diag1(2 * n - 1), used_diag2(2 * n - 1);
  auto rec = [&](auto &&rec, int row) -> bool {
    if (row == n) {
      return true;
    }
    for (int col = 0; col < n; col++) {
      int diag1 = row - col + n - 1, diag2 = row + col;
      if (used_col[col] || used_diag1[diag1] || used_diag2[diag2]) {
        continue;
      }
      cols[row] = col;
      used_col[col] = used_diag1[diag1] = used_diag2[diag2] = true;
      if (rec(rec, row + 1)) {
        return true;
      }
      used_col[col] = used_diag1[diag1] = used_diag2[diag2] = false;
      cols[row] = -1;
    }
    return false;
  };
  return rec(rec, 0) ? cols : std::vector<int>{};
}
\end{lstlisting}
\Needspace*{4\baselineskip}
\textbf{Example Usage}\par
\begin{lstlisting}[language=C++]
#include <cassert>
#include <vector>
using namespace std;

int main() {
  assert(n_queens(4) == vector<int>({1, 3, 0, 2}));
  assert(n_queens(2).empty());
  assert(n_queens(1) == vector<int>({0}));
  return 0;
}
\end{lstlisting}
\subsection{Knights Tour}
A knight's tour visits every square of an $n$-by-$n$ board exactly once, moving as a chess knight. Plain backtracking over the eight moves is hopeless past a small board, since the search tree has branching factor up to eight and depth $n^2$.

Warnsdorff's rule fixes the move order: from each square, try the onward squares in increasing order of how many unvisited moves they themselves have. The intuition is that squares with few remaining exits, such as corners, become unreachable if they are not used early, so the rule visits the most constrained squares first and leaves the flexible ones for later. This is the same minimum-remaining-values idea that drives the exact cover search of section 1.8.8.

Warnsdorff's rule alone is a heuristic and does fail on some boards, so the search below keeps backtracking as a fallback and merely orders the moves by the rule. That ordering is strong enough that a tour is normally found along the first path tried, with backtracking almost never invoked.

\begin{compactitem}
  \item \inlinecode{knight\_tour(n, sr = 0, sc = 0)} returns an \inlinecode{n}-by-\inlinecode{n} grid whose entry is the step at which the knight visits that square, numbered $[0, {\inlinecodemath{n}}^2)$, starting from square (\inlinecode{sr}, \inlinecode{sc}). It returns an empty grid when no tour exists, which is the case for ${\inlinecodemath{n}} \in \{2, 3, 4\}$. The board size must be positive and the starting square must lie on the board.
\end{compactitem}

A tour is closed when the knight can leap from its final square back to its start, which is possible only for even $n$ at least $6$. Requiring one means rejecting a completed tour whose last square is not a knight's move from the first, so the same search finds it by adding that test at the end.

\textbf{Time Complexity}: $O(n^{2})$ per call in practice, since the move ordering finds a tour without backtracking on the boards where one exists. The worst case remains exponential.

\textbf{Space Complexity}: $O(n^{2})$ auxiliary stack space and $O(n^{2})$ for the returned grid.

\begin{lstlisting}[language=C++]
#include <algorithm>
#include <array>
#include <cassert>
#include <utility>
#include <vector>

const std::array<std::pair<int, int>, 8> KNIGHT_MOVES = {
    {{-2, -1}, {-2, 1}, {-1, -2}, {-1, 2}, {1, -2}, {1, 2}, {2, -1}, {2, 1}}
};

std::vector<std::vector<int>> knight_tour(int n, int sr = 0, int sc = 0) {
  assert(n > 0 && sr >= 0 && sr < n && sc >= 0 && sc < n);
  std::vector<std::vector<int>> visit(n, std::vector<int>(n, -1));
  auto onward_moves = [&](int r, int c) {
    int count = 0;
    for (auto [dr, dc] : KNIGHT_MOVES) {
      int nr = r + dr, nc = c + dc;
      if (nr >= 0 && nr < n && nc >= 0 && nc < n && visit[nr][nc] == -1) {
        count++;
      }
    }
    return count;
  };
  auto rec = [&](auto &&rec, int r, int c, int step) -> bool {
    visit[r][c] = step;
    if (step == n * n - 1) {
      return true;
    }
    // Order the onward squares by how few exits they have left, which is Warnsdorff's rule.
    std::vector<std::pair<int, std::pair<int, int>>> candidates;
    for (auto [dr, dc] : KNIGHT_MOVES) {
      int r2 = r + dr, c2 = c + dc;
      if (r2 >= 0 && r2 < n && c2 >= 0 && c2 < n && visit[r2][c2] == -1) {
        candidates.emplace_back(onward_moves(r2, c2), std::make_pair(r2, c2));
      }
    }
    std::sort(candidates.begin(), candidates.end());
    for (auto [degree, square] : candidates) {
      if (rec(rec, square.first, square.second, step + 1)) {
        return true;
      }
    }
    visit[r][c] = -1;
    return false;
  };
  return rec(rec, sr, sc, 0) ? visit : std::vector<std::vector<int>>{};
}
\end{lstlisting}
\Needspace*{4\baselineskip}
\textbf{Example Usage}\par
\begin{lstlisting}[language=C++]
#include <cassert>
using namespace std;

int main() {
  assert(knight_tour(1) == vector<vector<int>>{{0}});
  assert(knight_tour(3).empty());  // No tour exists on boards of size 2, 3, or 4.
  assert(knight_tour(4).empty());

  for (int n : {5, 6, 8}) {
    auto visit = knight_tour(n, n / 2, n / 2);
    assert(!visit.empty());
    vector<pair<int, int>> at(n * n);
    for (int r = 0; r < n; r++) {
      for (int c = 0; c < n; c++) {
        assert(visit[r][c] >= 0 && visit[r][c] < n * n);
        at[visit[r][c]] = {r, c};  // Every step number is used exactly once.
      }
    }
    assert(at[0] == make_pair(n / 2, n / 2));
    for (int step = 1; step < n * n; step++) {
      int dr = abs(at[step].first - at[step - 1].first);
      int dc = abs(at[step].second - at[step - 1].second);
      assert(dr * dc == 2);  // Each consecutive pair differs by a knight's move.
    }
  }
  return 0;
}
\end{lstlisting}
\subsection{Exact Cover (Dancing Links)}
Given a binary matrix, the exact cover problem asks for a set of rows in which every column contains exactly one $1$. Many search problems reduce to it: a sudoku grid becomes a matrix whose columns demand that each cell hold one digit and that each digit appear once per row, column, and box, while a tiling becomes a matrix whose columns demand that each cell be covered by exactly one piece placement.

Knuth's Algorithm X solves it by repeatedly choosing an uncovered column, trying each row that covers it, and deleting from the matrix every column that row satisfies along with every other row conflicting with it. Dancing links is the representation that makes the deletions cheap. Each $1$ of the matrix is a node in two circular doubly linked lists, one along its row and one along its column, and every column has a header node holding the count of its remaining nodes. Removing a node \inlinecode{x} from a list is \inlinecode{left[right[x]] = left[x]} and \inlinecode{right[left[x]] = right[x]}, which leaves \inlinecode{x} itself still pointing at its old neighbors; restoring it on backtrack is therefore just \inlinecode{left[right[x]] = x} and \inlinecode{right[left[x]] = x}, with no allocation and no search. Undoing the covering operations in the exact reverse of the order that performed them restores the matrix exactly, which is what makes the whole search run in place on one static set of nodes.

The column to branch on is the one with the fewest remaining rows. This is the minimum-remaining- values heuristic, and it fails as early as possible: a column with no remaining rows ends the branch immediately, and a column with one forces the choice.

Secondary columns spare a problem with optional requirements from padding its matrix with one slack row per optional column. In the N-queens reduction below, every board row and column must hold exactly one queen and so is primary, while a diagonal may hold none and so is secondary. Asking for two solutions is the usual test of uniqueness, as a sudoku generator does when validating a puzzle.

\begin{compactitem}
  \item \inlinecode{ExactCover(primary, secondary = 0)} constructs an empty matrix with \inlinecode{primary + secondary} columns, both counts nonnegative. Columns in $[0, {\inlinecodemath{primary}})$ must be covered exactly once, while the optional columns in $[{\inlinecodemath{primary}}, {\inlinecodemath{primary}} + {\inlinecodemath{secondary}})$ may be covered at most once.
  \item \inlinecode{add\_row(cols)} appends a row containing a $1$ in each column of \inlinecode{cols}, which must all be in $[0, {\inlinecodemath{primary}} + {\inlinecodemath{secondary}})$, and returns the row's ID. The IDs of the first $r$ rows added are $[0, r)$, in the order they were added.
  \item \inlinecode{solve(limit = 1)} returns up to \inlinecode{limit} solutions, where \inlinecode{limit} is positive. Each solution is the sorted list of the IDs of its rows, and an empty result means no exact cover exists. The matrix is left unchanged, so it may be queried again.
\end{compactitem}

Solutions come in the order the search finds them, which follows the branching heuristic rather than row IDs.

\Needspace*{3\baselineskip}
\textbf{Time Complexity}:
\begin{compactitem}
  \item $O(k)$ per call to \inlinecode{add\_row(cols)}, where $k$ is the size of \inlinecode{cols}.
  \item Exponential per call to \inlinecode{solve()}. Exact cover is NP-complete, and dancing links prunes the search rather than avoiding it: the cost is proportional to the number of nodes the search unlinks and relinks, which the branching heuristic reduces but does not bound polynomially.
\end{compactitem}

\Needspace*{3\baselineskip}
\textbf{Space Complexity}:
\begin{compactitem}
  \item $O(n + c)$ for storage, where $n$ is the number of $1$s in the matrix and $c$ is the number of columns.
  \item $O(d)$ auxiliary stack space per call to \inlinecode{solve()}, where $d$ is the number of rows in the deepest partial solution reached, plus $O(s \cdot d)$ for the returned solutions, where $s$ is the number found.
\end{compactitem}

\begin{lstlisting}[language=C++]
#include <algorithm>
#include <cassert>
#include <vector>

class ExactCover {
  int columns, rows;
  std::vector<int> left, right, up, down, col_of, row_of, col_size;
  std::vector<int> partial;

  int add_node(int c, int r) {
    left.push_back(0);
    right.push_back(0);
    up.push_back(0);
    down.push_back(0);
    col_of.push_back(c);
    row_of.push_back(r);
    return static_cast<int>(col_of.size()) - 1;
  }

  // Unlinks column c and every row meeting it, so both vanish from all remaining traversals.
  void cover(int c) {
    right[left[c]] = right[c];
    left[right[c]] = left[c];
    for (int i = down[c]; i != c; i = down[i]) {
      for (int j = right[i]; j != i; j = right[j]) {
        up[down[j]] = up[j];
        down[up[j]] = down[j];
        col_size[col_of[j]]--;
      }
    }
  }

  // Relinks in the mirror order of cover(), which is what restores the matrix exactly.
  void uncover(int c) {
    for (int i = up[c]; i != c; i = up[i]) {
      for (int j = left[i]; j != i; j = left[j]) {
        col_size[col_of[j]]++;
        up[down[j]] = j;
        down[up[j]] = j;
      }
    }
    right[left[c]] = c;
    left[right[c]] = c;
  }

  bool search(int limit, std::vector<std::vector<int>> &found) {
    if (right[0] == 0) {  // Every primary column is covered.
      found.push_back(partial);
      std::sort(found.back().begin(), found.back().end());
      return static_cast<int>(found.size()) >= limit;
    }
    int best = right[0];
    for (int c = right[0]; c != 0; c = right[c]) {
      if (col_size[c] < col_size[best]) {
        best = c;
      }
    }
    if (col_size[best] == 0) {
      return false;  // No row can cover this column, so this branch is dead.
    }
    cover(best);
    for (int r = down[best]; r != best; r = down[r]) {
      partial.push_back(row_of[r]);
      for (int j = right[r]; j != r; j = right[j]) {
        cover(col_of[j]);
      }
      bool done = search(limit, found);
      for (int j = left[r]; j != r; j = left[j]) {
        uncover(col_of[j]);
      }
      partial.pop_back();
      if (done) {
        break;
      }
    }
    uncover(best);
    return static_cast<int>(found.size()) >= limit;
  }

 public:
  explicit ExactCover(int primary, int secondary = 0) : columns(primary + secondary), rows(0) {
    assert(primary >= 0 && secondary >= 0);
    col_size.assign(columns + 1, 0);
    for (int c = 0; c <= columns; c++) {  // Node 0 is the root, and node c + 1 heads column c.
      int header = add_node(c, -1);
      up[header] = down[header] = header;
      left[header] = right[header] = header;
    }
    int prev = 0;  // Secondary headers stay outside the root ring, so they are never branched on.
    for (int c = 1; c <= primary; c++) {
      right[prev] = c;
      left[c] = prev;
      prev = c;
    }
    right[prev] = 0;
    left[0] = prev;
  }

  int add_row(const std::vector<int> &cols) {
    for (int c : cols) {
      assert(c >= 0 && c < columns);
    }
    int row = rows++, first = -1;
    for (int c : cols) {
      int header = c + 1, node = add_node(header, row);
      up[node] = up[header];
      down[node] = header;
      down[up[header]] = node;
      up[header] = node;
      col_size[header]++;
      if (first == -1) {
        first = node;
        left[node] = right[node] = node;
      } else {
        left[node] = left[first];
        right[node] = first;
        right[left[first]] = node;
        left[first] = node;
      }
    }
    return row;
  }

  std::vector<std::vector<int>> solve(int limit = 1) {
    assert(limit > 0);
    std::vector<std::vector<int>> found;
    partial.clear();
    search(limit, found);
    return found;
  }
};
\end{lstlisting}
\Needspace*{4\baselineskip}
\textbf{Example Usage}\par
\begin{lstlisting}[language=C++]
#include <cassert>
using namespace std;

int main() {
  // Knuth's example matrix over 7 columns, whose unique cover is rows 1, 3, and 5.
  //     col: 0123456
  //   row 0: 1001001
  //   row 1: 1001000
  //   row 2: 0001101
  //   row 3: 0010110
  //   row 4: 0110011
  //   row 5: 0100001
  ExactCover matrix(7);
  matrix.add_row(vector<int>{0, 3, 6});
  matrix.add_row(vector<int>{0, 3});
  matrix.add_row(vector<int>{3, 4, 6});
  matrix.add_row(vector<int>{2, 4, 5});
  matrix.add_row(vector<int>{1, 2, 5, 6});
  matrix.add_row(vector<int>{1, 6});
  assert((matrix.solve() == vector<vector<int>>{{1, 3, 5}}));
  assert(matrix.solve(2).size() == 1);  // Asking for a second solution proves uniqueness.

  ExactCover impossible(2);
  impossible.add_row(vector<int>{0});
  assert(impossible.solve().empty());

  // The 8-queens problem. Each board row and each board column must hold exactly one queen, while
  // each diagonal may hold at most one, so the 2*n board lines are primary and the 2*(2*n - 1)
  // diagonals are secondary. One matrix row per square places a queen there.
  const int n = 8;
  ExactCover queens(2 * n, 2 * (2 * n - 1));
  for (int r = 0; r < n; r++) {
    for (int c = 0; c < n; c++) {
      queens.add_row(vector<int>{r, n + c, 2 * n + r + c, 2 * n + (2 * n - 1) + r - c + n - 1});
    }
  }
  assert(queens.solve().size() == 1);
  assert(queens.solve(1000).size() == 92);
  return 0;
}
\end{lstlisting}
